% Lesson 48 — Writing: Writes compositions on campaigning and voting for leaders — Anglais Tle A — Cours signé OnBuch+
% Programme officiel MINESEC. Compiler avec : tectonic EN48-writing-writes-compositions-on-campaigning-and-voting-f.tex
\newcommand{\DOCMATIERE}{Anglais}
\newcommand{\DOCNIVEAU}{Tle A}
\newcommand{\DOCMODULE}{Chapter 4 : Citizenship and Human Rights}
\newcommand{\DOCLECON}{EN48}
\newcommand{\DOCTITRE}{Lesson 48 — Writing: Writes compositions on campaigning and voting for leaders}
% OnBuch+ — préambule « cours signés » ENRICHI (compilé avec tectonic / XeLaTeX)
% Système visuel « Pop » d'OnBuch+ : fond crème (jamais blanc), bordures encre
% épaisses, ombres « dures » décalées, titres Archivo Black en capitales.
% Version MATHÉMATIQUES : graphes (pgfplots/tikz), boîtes pédagogiques
% (définition, propriété, méthode, attention, le savais-tu, exercice résolu,
% exercice, TP, bilan), page de garde et signature OnBuch+ ancrée en bas.
%
% Macros attendues dans main.tex AVANT \input{preamble} :
%   \DOCMATIERE  \DOCNIVEAU  \DOCMODULE  \DOCLECON (numéro)  \DOCTITRE
\documentclass[11pt]{article}

\usepackage{fontspec}
\usepackage[english,french]{babel} % français = langue principale ; l'anglais via \eng{..} / textebox
\usepackage[a4paper,margin=2cm,top=3.4cm,bottom=2.8cm,headheight=1.4cm,footskip=1.3cm]{geometry}
\usepackage{amsmath,amssymb,amsfonts}
\usepackage{mathtools} % \xmapsto, \coloneqq... souvent utilisés par les modèles
\usepackage{cancel} % simplifications barrées (bilans de forces)
% \abs{z} (module, valeur absolue) et \norm{u} (norme), utilisés par les modèles
\providecommand{\abs}{}\let\abs\relax\DeclarePairedDelimiter{\abs}{\lvert}{\rvert}
\providecommand{\norm}{}\let\norm\relax\DeclarePairedDelimiter{\norm}{\lVert}{\rVert}
\usepackage[table]{xcolor} % [table] : fournit aussi \rowcolors (alternance de lignes)
\usepackage{pagecolor}
\usepackage{tikz}
\usetikzlibrary{arrows.meta,positioning,calc,shapes.geometric,shapes.misc,shapes.arrows,decorations.pathmorphing,decorations.pathreplacing,decorations.markings,patterns,shadows,mindmap,trees,fit,backgrounds,matrix,intersections,angles,quotes}
\usepackage{pgfplots}
\usepackage[version=4]{mhchem} % équations nucléaires \ce{^{235}_{92}U}
\usepackage[european,siunitx]{circuitikz} % schémas électriques
\pgfplotsset{compat=1.18}
\usepgfplotslibrary{fillbetween,groupplots} % aires entre courbes (calcul intégral)
\usepackage{enumitem}
\usepackage{fancyhdr}
% [most] charge toutes les bibliothèques tcolorbox (listings, minted, raster,
% documentation...) et gonfle inutilement la mémoire TeX sur un document long
% (~300 boîtes) : on ne charge que ce qui est réellement utilisé.
\usepackage[skins,breakable]{tcolorbox}
\usepackage{graphicx}
\usepackage{colortbl}
\usepackage{booktabs}
\usepackage{tabularx}
\usepackage{diagbox} % case d'en-tête diagonale des tableaux à double entrée
% Colonnes extensibles centrée / gauche / droite (C, L, R), souvent utilisées par les modèles
\newcolumntype{C}{>{\centering\arraybackslash}X}
\newcolumntype{L}{>{\raggedright\arraybackslash}X}
\newcolumntype{R}{>{\raggedleft\arraybackslash}X}
\usepackage{array}
\usepackage{multicol}
\usepackage{multirow}
\usepackage{siunitx}
% parse-numbers=false : accepte aussi \SI{2,5 \times 10^{-3}}{...}
\sisetup{locale=FR,output-decimal-marker={,},group-minimum-digits=4,parse-numbers=false}
\DeclareSIUnit\ohm{\ensuremath{\symup{\Omega}}} % U+2126 absent de la police
\DeclareSIUnit\division{div} % réglages d'oscilloscope (ms/div, V/div)
\DeclareSIUnit\tr{tr} % tours (vitesse de rotation, ex. \SI{1500}{\tr\per\minute})
\DeclareSIUnit\molar{M} % concentration molaire (ex. \SI{3}{\molar}, \SI{1}{\micro\molar})
\DeclareSIUnit\an{an} % année (le modèle confond parfois avec \year, un compteur TeX) — ex. \SI{5}{\kilo\meter\per\an}
\DeclareSIUnit\FCFA{FCFA} % franc CFA (ex. \SI{500}{\FCFA\per\kilo\gram})
\DeclareSIUnit\degreeBrix{\textdegree Brix} % teneur en sucre d'un sirop/jus (réfractométrie)
\DeclareSIUnit\month{mois} % le modèle utilise parfois \month, un compteur TeX, comme unité de durée
\usepackage{qrcode}
% \degree : commande fréquemment utilisée par les modèles pour un angle ou une
% température en dehors de \SI{}{} (non fournie par les paquets chargés).
\providecommand{\degree}{^{\circ}}
% \qed (fin de démonstration), utilisé par les modèles sans amsthm
\providecommand{\qed}{\hfill$\square$}
% \barre : double barre des valeurs interdites dans un tableau de variations (tkz-tab)
\providecommand{\barre}{\ensuremath{\|}}
% environnement proof (amsthm non chargé) utilisé par les modèles
\ifdefined\proof\else\newenvironment{proof}[1][Démonstration]{\par\noindent\textit{#1.}\ }{\qed\par}\fi
\usepackage{lastpage}
\usepackage{hyperref}
\hypersetup{hidelinks}

% ============================================================
% Palette « Pop » OnBuch+ — mêmes hex que la classe P (plus_ui.dart)
% ============================================================
\definecolor{poporange}{HTML}{F59321}
\definecolor{popdark}{HTML}{E07A0C}
\definecolor{popcream}{HTML}{FFF6E8}
\definecolor{popink}{HTML}{1C1714}
\definecolor{poppurple}{HTML}{7A5AE0}
\definecolor{popgreen}{HTML}{2E9E6A}
\definecolor{poppink}{HTML}{E0457B}
\definecolor{popblue}{HTML}{2F7DE1}
\definecolor{popgold}{HTML}{C98A12}
\definecolor{popmuted}{HTML}{8A7F76}
\definecolor{popline}{HTML}{EADFD2}
\definecolor{popgray}{HTML}{8A7F76}
\colorlet{popgrey}{popgray}
% Alias tolérants (le modèle nomme parfois les couleurs différemment)
\colorlet{popwhite}{white}
\colorlet{popblack}{popink}
\colorlet{popred}{poppink}
\colorlet{popyellow}{popgold}
% Teintes claires (fonds de tableaux/encadrés) — le modèle nomme parfois
% les couleurs avec un suffixe L (ex. popblueL) sans les avoir définies.
\colorlet{poporangeL}{poporange!20}
\colorlet{popdarkL}{popdark!20}
\colorlet{poppurpleL}{poppurple!20}
\colorlet{popgreenL}{popgreen!20}
\colorlet{poppinkL}{poppink!20}
\colorlet{popblueL}{popblue!20}
\colorlet{popgoldL}{popgold!20}
\colorlet{popredL}{popred!20}
\colorlet{popgrayL}{popgray!20}
% Teintes claires pour fonds de tableaux / zones
\colorlet{popblueL}{popblue!10!white}
\colorlet{poporangeL}{poporange!14!white}
\colorlet{popgreenL}{popgreen!12!white}
\colorlet{poppurpleL}{poppurple!12!white}
\colorlet{poppinkL}{poppink!12!white}
\colorlet{popgoldL}{popgold!14!white}

\pagecolor{popcream}

% ============================================================
% Typographie
% ============================================================
\defaultfontfeatures{Ligatures=TeX}
\setmainfont{PlusJakartaSans-Medium.ttf}[Path=../../../fonts/,BoldFont=PlusJakartaSans-Bold.ttf]
% Maths en sans-serif (Fira Math) avec chiffres et lettres droites de Plus
% Jakarta, pour que les formules se fondent dans le texte.
\usepackage[math-style=ISO,bold-style=ISO]{unicode-math}
\setmathfont{FiraMath-Regular.otf}[Path=../../../fonts/]
\setmathfont{PlusJakartaSans-Medium.ttf}[Path=../../../fonts/,range={up/num,up/{Latin,latin},\mathup}]
\usepackage{icomma} % virgule décimale sans espace en mode math
% Caractères mathématiques Unicode absents des polices de texte (Plus Jakarta,
% Archivo Black) : rendus par leur équivalent mathématique, en texte comme en
% formule — sinon ils s'affichent en carré vide (ex. « Arithmétique dans ℤ »).
\usepackage{newunicodechar}
\newunicodechar{ℝ}{\ensuremath{\mathbb{R}}}
\newunicodechar{ℤ}{\ensuremath{\mathbb{Z}}}
\newunicodechar{ℂ}{\ensuremath{\mathbb{C}}}
\newunicodechar{ℕ}{\ensuremath{\mathbb{N}}}
\newunicodechar{ℚ}{\ensuremath{\mathbb{Q}}}
\newunicodechar{𝔻}{\ensuremath{\mathbb{D}}}
\newunicodechar{α}{\ensuremath{\alpha}}
\newunicodechar{β}{\ensuremath{\beta}}
\newunicodechar{γ}{\ensuremath{\gamma}}
\newunicodechar{δ}{\ensuremath{\delta}}
\newunicodechar{ε}{\ensuremath{\varepsilon}}
\newunicodechar{θ}{\ensuremath{\theta}}
\newunicodechar{λ}{\ensuremath{\lambda}}
\newunicodechar{μ}{\ensuremath{\mu}}
\newunicodechar{σ}{\ensuremath{\sigma}}
\newunicodechar{τ}{\ensuremath{\tau}}
\newunicodechar{φ}{\ensuremath{\varphi}}
\newunicodechar{ψ}{\ensuremath{\psi}}
\newunicodechar{ω}{\ensuremath{\omega}}
\newunicodechar{Σ}{\ensuremath{\Sigma}}
% majuscules produites par \MakeUppercase dans les titres (α-aminés -> Α)
\newunicodechar{Α}{\ensuremath{\alpha}}
\newunicodechar{Β}{\ensuremath{\beta}}
\newunicodechar{Γ}{\ensuremath{\Gamma}}
\newunicodechar{Δ}{\ensuremath{\Delta}}
\newunicodechar{Ε}{\ensuremath{\varepsilon}}
\newunicodechar{Θ}{\ensuremath{\Theta}}
\newunicodechar{Λ}{\ensuremath{\Lambda}}
\newunicodechar{Μ}{\ensuremath{\mu}}
\newunicodechar{Τ}{\ensuremath{\tau}}
\newunicodechar{Φ}{\ensuremath{\Phi}}
\newunicodechar{Ψ}{\ensuremath{\Psi}}
\newunicodechar{Ω}{\ensuremath{\Omega}}
\newunicodechar{↦}{\ensuremath{\mapsto}}
\newunicodechar{∈}{\ensuremath{\in}}
\newunicodechar{∉}{\ensuremath{\notin}}
\newunicodechar{∧}{\ensuremath{\wedge}}
\newunicodechar{⇔}{\ensuremath{\Leftrightarrow}}
\newunicodechar{⟺}{\ensuremath{\Longleftrightarrow}}
\newunicodechar{⇒}{\ensuremath{\Rightarrow}}
\newunicodechar{⟹}{\ensuremath{\Longrightarrow}}
\newunicodechar{∘}{\ensuremath{\circ}}
\newunicodechar{∪}{\ensuremath{\cup}}
\newunicodechar{∩}{\ensuremath{\cap}}
\newunicodechar{≡}{\ensuremath{\equiv}}
\newunicodechar{⊂}{\ensuremath{\subset}}
\newunicodechar{⊥}{\ensuremath{\perp}}
\newunicodechar{∃}{\ensuremath{\exists}}
\newunicodechar{∀}{\ensuremath{\forall}}
\newunicodechar{∛}{\ensuremath{\sqrt[3]{\,}}}
\newunicodechar{∜}{\ensuremath{\sqrt[4]{\,}}}
\newunicodechar{⃗}{\ensuremath{{}^{\to}}}
\newunicodechar{ⁿ}{\ifmmode^{n}\else\textsuperscript{\ensuremath{n}}\fi}
\newunicodechar{ˣ}{\ifmmode^{x}\else\textsuperscript{\ensuremath{x}}\fi}
\newunicodechar{ᵇ}{\ifmmode^{b}\else\textsuperscript{\ensuremath{b}}\fi}
\newunicodechar{ʳ}{\ifmmode^{r}\else\textsuperscript{\ensuremath{r}}\fi}
\newunicodechar{ᵘ}{\ifmmode^{u}\else\textsuperscript{\ensuremath{u}}\fi}
\newunicodechar{ʸ}{\ifmmode^{y}\else\textsuperscript{\ensuremath{y}}\fi}
\newunicodechar{ᵃ}{\ifmmode^{a}\else\textsuperscript{\ensuremath{a}}\fi}
\newunicodechar{ᵉ}{\ifmmode^{e}\else\textsuperscript{\ensuremath{e}}\fi}
\newunicodechar{ᵏ}{\ifmmode^{k}\else\textsuperscript{\ensuremath{k}}\fi}
\newunicodechar{ᵖ}{\ifmmode^{p}\else\textsuperscript{\ensuremath{p}}\fi}
\newunicodechar{ᴺ}{\ifmmode^{N}\else\textsuperscript{\ensuremath{N}}\fi}
\newunicodechar{ᶜ}{\ifmmode^{c}\else\textsuperscript{\ensuremath{c}}\fi}
\newunicodechar{ʰ}{\ifmmode^{h}\else\textsuperscript{\ensuremath{h}}\fi}
\newunicodechar{ᵠ}{\ifmmode^{q}\else\textsuperscript{\ensuremath{q}}\fi}
\newunicodechar{ₙ}{\ifmmode_{n}\else\textsubscript{\ensuremath{n}}\fi}
\newunicodechar{ₐ}{\ifmmode_{a}\else\textsubscript{\ensuremath{a}}\fi}
\newunicodechar{ᵢ}{\ifmmode_{i}\else\textsubscript{\ensuremath{i}}\fi}
\newunicodechar{ᵣ}{\ifmmode_{r}\else\textsubscript{\ensuremath{r}}\fi}
\newunicodechar{ₖ}{\ifmmode_{k}\else\textsubscript{\ensuremath{k}}\fi}
\newunicodechar{ᵦ}{\ifmmode_{\beta}\else\textsubscript{\ensuremath{\beta}}\fi}
\newunicodechar{ₓ}{\ifmmode_{x}\else\textsubscript{\ensuremath{x}}\fi}
\newfontfamily\popdisplay{ArchivoBlack-Regular.ttf}[Path=../../../fonts/]
\newfontfamily\poplogo{SpaceGrotesk-Bold.ttf}[Path=../../../fonts/]
\newfontfamily\popsemi{PlusJakartaSans-SemiBold.ttf}[Path=../../../fonts/]
\newfontfamily\popxbold{PlusJakartaSans-ExtraBold.ttf}[Path=../../../fonts/]
\newfontfamily\popxboldls{PlusJakartaSans-ExtraBold.ttf}[Path=../../../fonts/,LetterSpace=3.0]

\setlist[enumerate]{leftmargin=1.7em,itemsep=5pt,topsep=4pt}
\setlist[itemize]{leftmargin=1.7em,itemsep=4pt,topsep=4pt,label={\color{poporange}\textbullet}}
\setlength{\parindent}{0pt}
\setlength{\parskip}{5pt}
\renewcommand{\arraystretch}{1.25}

% pgfplots : style Pop par défaut
\pgfplotsset{
  every axis/.append style={axis line style={popink,line width=1pt},tick style={popink},
    grid style={popline,line width=0.5pt},label style={font=\small},tick label style={font=\footnotesize}},
  popcurve/.style={poporange,line width=1.8pt,no markers,smooth},
}
% Même style disponible dans un \draw TikZ simple (hors pgfplots)
\tikzset{popcurve/.style={poporange,line width=1.8pt}}
\tikzset{popgrid/.style={popline,very thin}}
\colorlet{popcurve}{poporange} % le nom du style est parfois utilisé comme couleur

% ============================================================
% Pill / squiggle
% ============================================================
\newcommand{\poppill}[3]{%
  \tikz[baseline=-0.55ex]\node[draw=popink,line width=1.2pt,rounded corners=2.6pt,%
    fill=#2,inner xsep=6.5pt,inner ysep=3pt]{%
    \popxboldls\fontsize{7}{7}\selectfont\color{#3}\MakeUppercase{#1}};%
}
\newcommand{\popsquiggle}[1]{%
  \tikz[baseline=-0.4ex]{
    \draw[#1,line width=2.6pt,line cap=round]
      plot [smooth] coordinates {(0,0.09) (0.34,0.01) (0.68,0.21) (1.02,0.01) (1.36,0.10)};
  }%
}
% Mot-clé surligné (marqueur orange sous le texte)
\newcommand{\cle}[1]{{\popxbold\color{popdark}#1}}

% ============================================================
% En-tête / pied
% ============================================================
\pagestyle{fancy}
\fancyhf{}
\renewcommand{\headrulewidth}{0pt}
\renewcommand{\footrulewidth}{0pt}
\lhead{\raisebox{-0.15ex}{\poplogo\fontsize{13}{13}\selectfont\color{popink}OnBuch\color{poporange}+}}
\rhead{\poppill{\DOCMATIERE\ \textbullet\ \DOCNIVEAU\ \textbullet\ Leçon \DOCLECON}{white}{popink}}
\lfoot{\popxbold\fontsize{7.6}{7.6}\selectfont\color{popmuted}\MakeUppercase{Cours signé OnBuch+}\ \textbullet\ {\popsemi\DOCTITRE}}
\rfoot{\tikz[baseline=-0.6ex]\node[rounded corners=8.5pt,fill=popink,minimum height=17pt,minimum width=17pt,inner xsep=5pt,inner sep=0]{\popxbold\fontsize{8}{8}\selectfont\color{popcream}\thepage\,/\,\pageref*{LastPage}};}

% ============================================================
% Titres
% ============================================================
\newcommand{\chaptitre}[2]{%
  \begin{tcolorbox}[enhanced,colback=poporange,colframe=popink,boxrule=2.6pt,arc=11pt,
      drop shadow=popink,left=15pt,right=15pt,top=12pt,bottom=11pt,before skip=3pt,after skip=18pt]
    \poppill{#2}{popink}{popcream}\par\vspace{9pt}
    {\raggedright\hyphenpenalty=10000\exhyphenpenalty=10000\popdisplay\fontsize{22}{24}\selectfont\color{popcream}\MakeUppercase{#1}\par}
  \end{tcolorbox}
}
% Section numérotée : pastille orange avec le numéro + titre Archivo
\newcounter{popsec}
\newcommand{\coursec}[1]{%
  \stepcounter{popsec}\setcounter{popsub}{0}%
  \par\vspace{16pt}\noindent
  \begin{minipage}{\linewidth}
  \tikz[baseline=-0.7ex]\node[draw=popink,line width=1.6pt,fill=poporange,rounded corners=4pt,minimum size=22pt,inner sep=0]{\popdisplay\fontsize{12}{12}\selectfont\color{popcream}\thepopsec};%
  \hspace{8pt}{\popdisplay\fontsize{15}{17}\selectfont\color{popink}\MakeUppercase{#1}}\par
  \vspace{4pt}\noindent{\color{poporange}\rule{36pt}{3.6pt}}
  \end{minipage}\par\nopagebreak\vspace{8pt}
}
\newcounter{popsub}
\newcommand{\courssub}[1]{%
  \stepcounter{popsub}%
  \par\vspace{9pt}\noindent{\popxbold\fontsize{11.5}{13}\selectfont\color{popink}{\color{poporange}\thepopsec.\thepopsub}\hspace{6pt}#1}\par\nopagebreak\vspace{4pt}
}

% ============================================================
% Boîtes pédagogiques — même recette : fond blanc, bordure épaisse colorée,
% ombre dure encre, badge plein. Titre optionnel [#1].
% ============================================================
\tcbset{popbase/.style={breakable,enhanced,colback=white,boxrule=2.3pt,arc=9pt,
  shadow={3pt}{-3pt}{0pt}{popink},left=14pt,right=14pt,top=11pt,bottom=11pt,before skip=11pt,after skip=15pt,
  fontupper=\popsemi\fontsize{10.6}{14.5}\selectfont\color{popink},
  fontlower=\popsemi\fontsize{10.6}{14.5}\selectfont\color{popink}}}
% Coche dessinée (le glyphe ✓ n'existe pas dans les polices du document)
\newcommand{\popcheck}[1]{\tikz[baseline=-0.5ex]\draw[#1,line width=1.6pt,line cap=round,line join=round](-0.13,0.0)--(-0.04,-0.09)--(0.13,0.1);}
\providecommand{\coche}{\popcheck{popgreen}}
\providecommand{\resultat}[1]{\par\smallskip\noindent\fcolorbox{popgreen}{popgreenL}{\parbox{\dimexpr\linewidth-2\fboxsep-2\fboxrule\relax}{\textbf{Résultat.} #1}}\par\smallskip}
\newcommand{\popboxhead}[3]{\poppill{#1}{#2}{white}\ifx\relax#3\relax\else\hspace{6pt}{\popxbold\fontsize{10.6}{12}\selectfont\color{popink}#3}\fi\par\vspace{7pt}}

\newtcolorbox{aretenir}[1][]{popbase,colframe=popgreen,before upper={\popboxhead{À retenir}{popgreen}{#1}}}
% Texte en anglais (document de lecture, dialogue, extrait) : cadre
% d'étude. Un titre optionnel (source, auteur) s'affiche dans l'en-tête.
\newtcolorbox{textebox}[1][]{popbase,colframe=popblue,before upper={\popboxhead{Text}{popblue}{#1}\begin{otherlanguage}{english}},after upper={\end{otherlanguage}}}
% Marques ✓ / ✗ (forme correcte / fautive) souvent utilisées par les modèles
\providecommand{\cross}{\textcolor{poppink}{\ensuremath{\times}}}
\providecommand{\xmark}{\cross}\providecommand{\crossmark}{\cross}\providecommand{\cmark}{\ensuremath{\checkmark}}
% Ligne de réponse / justification à compléter (inventée par les modèles)
\providecommand{\justification}[1]{\par\noindent\textit{Justification~:} \dotfill\par}
% \textipa (package tipa non chargé) : la police ne couvre pas l'API -> simple passage
\providecommand{\textipa}[1]{#1}
% Soulignement (ulem non chargé) : \uline, \ul
\providecommand{\uline}[1]{\underline{#1}}\providecommand{\ul}[1]{\underline{#1}}
% \glossaire{\item ...} inventée par les modèles : liste de glossaire
\providecommand{\glossaire}[1]{\begin{itemize}#1\end{itemize}}
% \alert (beamer) inventée par les modèles : texte mis en évidence
\providecommand{\alert}[1]{\textbf{\textcolor{poppink}{#1}}}
% variantes ASCII d'ordinaux écrites en macro
\providecommand{\iere}{\textsuperscript{re}}\providecommand{\ieme}{\textsuperscript{e}}\providecommand{\ere}{\textsuperscript{re}}\providecommand{\eme}{\textsuperscript{e}}\providecommand{\er}{\textsuperscript{er}}\providecommand{\ie}{\textsuperscript{e}}
% Ordinaux français écrits en macro par les modèles : 1\ière, 3\ième
\newcommand{\ière}{\textsuperscript{re}}\newcommand{\ième}{\textsuperscript{e}}
% \exemple (inventée par les modèles) : étiquette « Exemple : »
\providecommand{\exemple}{\textbf{Exemple~:}\ }
% \square utilisable aussi hors mode math (cases à cocher)
\AtBeginDocument{\let\squareMath\square\renewcommand{\square}{\ensuremath{\squareMath}}}
% Mot ou phrase en anglais à l'intérieur d'un paragraphe français
\newcommand{\eng}[1]{\ifmmode\text{\textit{#1}}\else\begingroup\selectlanguage{english}\itshape #1\endgroup\fi}
\let\esp\eng % alias toléré
% flèches d'intonation inventées par les modèles
\providecommand{\textsearrow}{}\renewcommand{\textsearrow}{\ensuremath{\searrow}}\providecommand{\textnearrow}{}\renewcommand{\textnearrow}{\ensuremath{\nearrow}}
% \fr{..} : mot français (inventé par les modèles)
\providecommand{\fr}[1]{\textit{#1}}
\newtcolorbox{exemplebox}[1][]{popbase,colframe=poporange,before upper={\popboxhead{Exemple}{poporange}{#1}}}
\newtcolorbox{definition}[1][]{popbase,colframe=popblue,before upper={\popboxhead{Définition}{popblue}{#1}}}
\newtcolorbox{propriete}[1][]{popbase,colframe=poppurple,before upper={\popboxhead{Propriété}{poppurple}{#1}}}
\newtcolorbox{methode}[1][]{popbase,colframe=popgold,before upper={\popboxhead{Méthode}{popgold}{#1}}}
\newtcolorbox{attention}[1][]{popbase,colframe=poppink,before upper={\popboxhead{Attention}{poppink}{#1}}}
\newtcolorbox{experience}[1][]{popbase,colframe=popblue,colback=popblueL,before upper={\popboxhead{Expérience}{popblue}{#1}}}
% « Le savais-tu ? » : carte violette pleine (contexte, culture, Cameroun)
\newtcolorbox{savaistu}[1][]{popbase,colback=poppurple,colframe=popink,
  fontupper=\popsemi\fontsize{10.4}{14.2}\selectfont\color{white},
  before upper={\poppill{Le savais-tu\,?}{white}{poppurple}\ifx\relax#1\relax\else\hspace{6pt}{\popxbold\color{white}#1}\fi\par\vspace{7pt}}}
% Exercice résolu : énoncé en haut, \tcblower puis la solution sur fond crème
\newcounter{popexo}\newcounter{popexr}
\newtcolorbox{exoresolu}[1][]{popbase,colframe=poporange,colbacklower=popcream,
  segmentation style={popink,line width=1.2pt,dashed},
  before upper={\stepcounter{popexr}\popboxhead{Exercice résolu \thepopexr}{poporange}{#1}},
  before lower={\poppill{Solution}{popink}{popcream}\par\vspace{6pt}}}
% Exercice (non résolu en place) — la correction est dans la partie Corrigés
\newtcolorbox{exercice}[1][]{popbase,colframe=popink,
  before upper={\stepcounter{popexo}\popboxhead{Exercice \thepopexo}{popink}{#1}}}
% Corrigé
\newtcolorbox{corrige}[1][]{popbase,colframe=popgreen,colback=popgreenL,
  before upper={\popboxhead{Corrigé}{popgreen}{#1}}}
% Objectifs / prérequis / bilan
\newtcolorbox{objectifs}{popbase,colframe=popink,colback=poporangeL,
  before upper={\popboxhead{Objectifs de la leçon}{popink}{}}}
\newtcolorbox{prerequis}{popbase,colframe=popink,colback=popblueL,
  before upper={\popboxhead{Prérequis}{popblue}{}}}
\newtcolorbox{bilan}{popbase,colframe=popink,colback=white,boxrule=2.8pt,
  before upper={\popboxhead{Fiche bilan}{poporange}{L'essentiel à connaître pour l'examen}}}
% Figure encadrée avec légende
\newtcolorbox{popfigure}[1][]{enhanced,breakable=false,colback=white,colframe=popline,boxrule=1.4pt,arc=8pt,
  left=10pt,right=10pt,top=10pt,bottom=8pt,before skip=10pt,after skip=12pt,halign=center,
  lower separated=false,halign lower=center,
  fontlower=\popsemi\fontsize{9}{11.5}\selectfont\color{popmuted},
  #1}
\newcommand{\legende}[1]{\par\vspace{6pt}{\popsemi\fontsize{9}{11.5}\selectfont\color{popmuted}#1}}

% ============================================================
% Signature OnBuch+
% ============================================================
\newcommand{\poplogoinline}[1]{{\poplogo\fontsize{#1}{#1}\selectfont\color{popink}OnBuch\color{poporange}+}}
\newcommand{\signatureonbuch}{%
  \begin{tcolorbox}[enhanced,colback=popink,colframe=popink,boxrule=2.2pt,arc=10pt,
      drop shadow=poporange,left=14pt,right=14pt,top=10pt,bottom=10pt,before skip=0pt,after skip=0pt]
    \tikz[baseline=-0.7ex]\node[circle,fill=popgreen,draw=popcream,line width=1.3pt,minimum size=22pt,inner sep=0]{\popcheck{white}};%
    \hspace{9pt}\begin{minipage}[c]{0.62\linewidth}
      {\popxbold\fontsize{12}{13}\selectfont\color{popcream}Signé\ \textbullet\ {\poplogo OnBuch\color{poporange}+}}\par\vspace{2pt}
      {\raggedright\popsemi\fontsize{8}{10.5}\selectfont\color{popline}\MakeUppercase{Équipe pédagogique OnBuch+}\\\MakeUppercase{Conforme au programme officiel MINESEC}\par}
    \end{minipage}\hfill
    {\poplogo\fontsize{22}{22}\selectfont\color{popcream}OnBuch\color{poporange}+}
  \end{tcolorbox}%
}

% ============================================================
% Page de garde
% #1 titre · #2 matière · #3 niveau · #4 module · #5 n° leçon · #6 durée ·
% #7 description / objectifs (texte ou itemize)
% ============================================================
\newcommand{\pagegarde}[7]{%
  \thispagestyle{empty}
  \newgeometry{a4paper,margin=1.7cm,top=1.6cm,bottom=1.4cm}%
  \begin{tikzpicture}[remember picture,overlay]
    \fill[poporange!18] ([xshift=-3.2cm,yshift=-3.6cm]current page.north east) circle (4.6cm);
    \fill[poppurple!14] ([xshift=2.4cm,yshift=7.6cm]current page.south west) circle (3.8cm);
  \end{tikzpicture}%
  {\poplogo\fontsize{30}{30}\selectfont\color{popink}OnBuch\color{poporange}+}\hfill
  \raisebox{0.6ex}{\poppill{Cours signé \textbullet\ Premium}{popink}{popcream}}\par
  \vspace{1pt}\popsquiggle{poppurple}\par
  \vspace{8pt}
  \begin{tcolorbox}[enhanced,colback=poporange,colframe=popink,boxrule=2.8pt,arc=13pt,
      drop shadow=popink,left=18pt,right=18pt,top=12pt,bottom=13pt,after skip=10pt]
    \poppill{#2 \textbullet\ #3}{popink}{popcream}\hspace{5pt}\poppill{#4}{white}{popink}\par\vspace{8pt}
    {\popdisplay\fontsize{14}{14}\selectfont\color{popink}LEÇON #5}\par\vspace{4pt}
    {\raggedright\hyphenpenalty=10000\exhyphenpenalty=10000\popdisplay\fontsize{25}{27}\selectfont\color{popcream}\MakeUppercase{#1}\par}
    \vspace{8pt}
    {\popxbold\fontsize{9.5}{9.5}\selectfont\color{popink}\MakeUppercase{Durée conseillée : #6}}
  \end{tcolorbox}
  \begin{tcolorbox}[enhanced,colback=white,colframe=popink,boxrule=2.2pt,arc=10pt,
      drop shadow=popink,left=16pt,right=16pt,top=10pt,bottom=10pt,after skip=8pt]
    \poppill{Dans cette leçon}{poppurple}{white}\par\vspace{5pt}
    \popsemi\fontsize{9.3}{12.6}\selectfont\color{popink}#7
  \end{tcolorbox}
  \noindent\begin{minipage}[t]{0.487\linewidth}
    \begin{tcolorbox}[enhanced,colback=popgreen,colframe=popink,boxrule=2.2pt,arc=10pt,
        drop shadow=popink,left=11pt,right=11pt,top=10pt,bottom=10pt,height=3.1cm,valign=center]
      \tcbox[colback=white,colframe=popink,boxrule=1.3pt,arc=5pt,boxsep=0pt,nobeforeafter,
        left=3pt,right=3pt,top=3pt,bottom=3pt,valign=center]{\qrcode[height=1.9cm]{https://play.google.com/store/apps/details?id=com.luvvix.onbuch&hl=fr}}%
      \hspace{8pt}\begin{minipage}[c]{0.5\linewidth}
        {\popdisplay\fontsize{11}{12.5}\selectfont\color{white}\MakeUppercase{Télécharge OnBuch}}\par\vspace{4pt}
        {\raggedright\popsemi\fontsize{8.4}{11}\selectfont\color{white}Gratuit \textbullet\ Google Play\\Tuteur IA, annales, concours, résultats\par}
      \end{minipage}
    \end{tcolorbox}
  \end{minipage}\hfill
  \begin{minipage}[t]{0.487\linewidth}
    \begin{tcolorbox}[enhanced,colback=poppurple,colframe=popink,boxrule=2.2pt,arc=10pt,
        drop shadow=popink,left=11pt,right=11pt,top=10pt,bottom=10pt,height=3.1cm,valign=center]
      \tikz[baseline=-0.7ex]\node[circle,fill=white,draw=popink,line width=1.3pt,minimum size=1.6cm,inner sep=0]{\popdisplay\fontsize{24}{24}\selectfont\color{poppurple}+};%
      \hspace{8pt}\begin{minipage}[c]{0.6\linewidth}
        {\popdisplay\fontsize{11}{12.5}\selectfont\color{white}\MakeUppercase{Passe à OnBuch+}}\par\vspace{4pt}
        {\raggedright\popsemi\fontsize{8.4}{11}\selectfont\color{white}Tous les cours signés, Tuteur IA illimité, fiches PDF — onglet OnBuch+ de l'app\par}
      \end{minipage}
    \end{tcolorbox}
  \end{minipage}
  \vspace{8pt}
  \popcontenu
  \vfill
  \signatureonbuch
  \restoregeometry
  \newpage
}

% Rangée « Ce cours contient » (page de garde)
\newcommand{\poptile}[3]{% #1 couleur #2 gros texte #3 légende
  \begin{tcolorbox}[enhanced,colback=#1,colframe=popink,boxrule=2pt,arc=9pt,shadow={2.5pt}{-2.5pt}{0pt}{popink},
      left=6pt,right=6pt,top=9pt,bottom=9pt,halign=center,valign=center,height=1.75cm,width=\linewidth,nobeforeafter]
    {\popdisplay\fontsize{12.5}{13}\selectfont\color{white}\MakeUppercase{#2}}\par\vspace{3pt}
    {\centering\popsemi\fontsize{7.8}{9.5}\selectfont\color{white}#3\par}
  \end{tcolorbox}}
\newcommand{\popcontenu}{%
  \par\noindent{\popxboldls\fontsize{7.5}{7.5}\selectfont\color{popmuted}CE COURS SIGNÉ CONTIENT}\par\vspace{7pt}
  \noindent\begin{minipage}[t]{0.235\linewidth}\poptile{popblue}{Cours}{complet, illustré, pas à pas}\end{minipage}\hfill
  \begin{minipage}[t]{0.235\linewidth}\poptile{poporange}{Méthodes}{et exercices résolus}\end{minipage}\hfill
  \begin{minipage}[t]{0.235\linewidth}\poptile{poppink}{Exercices}{type examen + corrigés}\end{minipage}\hfill
  \begin{minipage}[t]{0.235\linewidth}\poptile{popgreen}{Fiche bilan}{l'essentiel à retenir}\end{minipage}\par}

% Sommaire
\newtcolorbox{sommaire}{enhanced,colback=white,colframe=popink,boxrule=2.2pt,arc=10pt,
  drop shadow=popink,left=18pt,right=18pt,top=13pt,bottom=13pt,before skip=6pt,after skip=20pt,
  before upper={\poppill{Sommaire}{poppurple}{white}\par\vspace{9pt}\popsemi\fontsize{10.6}{14.5}\selectfont\color{popink}}}

% Signature de fin de cours — poussée tout en bas de la dernière page
\newcommand{\findecours}{%
  \par\vfill\nopagebreak
  \begin{center}\popsquiggle{poporange}\end{center}\vspace{4pt}
  \signatureonbuch
}
\AtBeginDocument{\def\llbracket{\mathopen{[\mkern-3.5mu[}}\def\rrbracket{\mathclose{]\mkern-3.5mu]}}} % intervalles d'entiers (glyphe ⟦ absent de la police)
% Glyphes absents de Fira Math : redéfinis à partir de glyphes disponibles
\AtBeginDocument{%
  \def\setminus{\mathbin{\backslash}}\let\smallsetminus\setminus
  \def\vdots{\mathord{\vbox{\baselineskip3.2pt\lineskiplimit0pt\kern4pt\hbox{.}\hbox{.}\hbox{.}}}}%
  \def\ddots{\mathinner{\mkern1mu\raise6.5pt\hbox{.}\mkern2mu\raise3.6pt\hbox{.}\mkern2mu\raise0.7pt\hbox{.}\mkern1mu}}%
  \def\triangle{\mathord{\scalebox{0.9}{$\symup{\Delta}$}}}%
  \def\nabla{\mathord{\raisebox{\depth}{\scalebox{1}[-1]{$\symup{\Delta}$}}}}%
  \def\checkmark{\popcheck{popdark}}%
  \def\star{\mathbin{\ast}}\def\bigstar{\mathord{\ast}}%
  \def\diamond{\mathbin{\scalebox{0.6}{$\Diamond$}}}%
  \def\bigcup{\mathop{\mathchoice{\raisebox{-0.3em}{\scalebox{1.7}{$\cup$}}}{\scalebox{1.25}{$\cup$}}{\cup}{\cup}}\displaylimits}%
  \def\bigcap{\mathop{\mathchoice{\raisebox{-0.3em}{\scalebox{1.7}{$\cap$}}}{\scalebox{1.25}{$\cap$}}{\cap}{\cap}}\displaylimits}%
  \def\lesseqqgtr{\mathrel{\lessgtr}}\def\bigoplus{\mathop{\oplus}\displaylimits}%
}
\providecommand{\ssi}{si et seulement si} % abréviation fréquente
\AtBeginDocument{\providecommand{\wideparen}{\overparen}} % arc de cercle

%% FIGFIX-BEGIN v4 — correctifs globaux des figures (docs/onbuch-generation/tools/figfix)
\usepackage{adjustbox}\usepackage{etoolbox}
% toute figure TikZ/pgfplots plus large que la ligne est réduite à la largeur disponible
\BeforeBeginEnvironment{tikzpicture}{\begin{adjustbox}{max width=\linewidth}}
\AfterEndEnvironment{tikzpicture}{\end{adjustbox}}
% légendes pgfplots lisibles par-dessus les courbes
\pgfplotsset{every axis/.append style={legend style={fill=white,fill opacity=0.92,text opacity=1,draw=black!15,inner sep=3pt}}}
% étiquettes d'arêtes (syntaxe edge["..."]) avec fond blanc
\tikzset{every edge quotes/.append style={fill=white,inner sep=1.2pt}}
%% FIGFIX-END
\begin{document}

\pagegarde{Lesson 48 — Writing: Writes compositions on campaigning and voting for leaders}{Anglais}{Tle A}{Chapter 4 : Citizenship and Human Rights}{EN48}{2 h}{Cette leçon de 2 heures guide l'élève de la compréhension du thème "Campagne et vote" à la rédaction autonome d'une composition argumentative. Elle couvre le vocabulaire thématique, la structure type (introduction, développement, conclusion), les connecteurs logiques, l'analyse d'un modèle commenté, les pièges à éviter et la grille d'auto-évaluation alignée sur les attendus du Bac.
\vspace{4pt}
\begin{multicols}{2}\small
\begin{itemize}
\item À la fin de cette leçon, je sais identifier et utiliser le vocabulaire spécifique aux élections et à la campagne électorale.
\item Je sais structurer une composition en trois parties distinctes : introduction (accroche, définition, thèse, plan), développement (arguments illustrés), conclusion (bilan, ouverture).
\item Je sais sélectionner et employer les connecteurs logiques adaptés (addition, opposition, cause, conséquence, illustration) pour assurer la cohérence du texte.
\item Je sais conjuguer les temps verbaux appropriés (present simple, modals, conditionnel, futur) selon la visée argumentative.
\item Je sais analyser un modèle de composition pour en repérer les forces structurelles et langagières.
\item Je sais planifier ma rédaction grâce à un brouillon organisé (mind-mapping, outline).
\end{itemize}\end{multicols}}

\begin{sommaire}
\begin{multicols}{2}
\begin{enumerate}
\item 1. Vocabulaire thématique : Elections \& Campaigning
\item 2. Structure type de la composition argumentative
\item 3. Connecteurs logiques et Temps verbaux : Outils de cohérence
\item 4. Modèle commenté : "Why Voting Matters"
\item 5. Atelier de rédaction guidée : Du brouillon à la copie finale
\item 6. Grille d'évaluation \& Erreurs types des francophones (Pièges du Bac)
\item Activité d'intégration
\item Méthodes et exercices résolus
\item Exercices
\item Corrigés des exercices
\item Fiche bilan
\end{enumerate}
\end{multicols}
\end{sommaire}

\begin{objectifs}
\begin{itemize}
\item À la fin de cette leçon, je sais identifier et utiliser le vocabulaire spécifique aux élections et à la campagne électorale.
\item Je sais structurer une composition en trois parties distinctes : introduction (accroche, définition, thèse, plan), développement (arguments illustrés), conclusion (bilan, ouverture).
\item Je sais sélectionner et employer les connecteurs logiques adaptés (addition, opposition, cause, conséquence, illustration) pour assurer la cohérence du texte.
\item Je sais conjuguer les temps verbaux appropriés (present simple, modals, conditionnel, futur) selon la visée argumentative.
\item Je sais analyser un modèle de composition pour en repérer les forces structurelles et langagières.
\item Je sais planifier ma rédaction grâce à un brouillon organisé (mind-mapping, outline).
\item Je sais relire et corriger ma production en utilisant une grille d'évaluation (contenu, organisation, langue, vocabulaire).
\item Je sais éviter les erreurs fréquentes des francophones (faux-amis, syntaxe, accord, articles).
\end{itemize}
\end{objectifs}

\begin{prerequis}
\begin{itemize}
\item Maîtrise de la structure de base du paragraphe (topic sentence, supporting sentences, concluding sentence) vue en Première.
\item Connaissance des principaux connecteurs logiques (however, therefore, furthermore, for example) et de leur ponctuation.
\item Utilisation des modaux d'obligation, conseil, possibilité (must, should, can, could) pour exprimer l'opinion et l'argumentation.
\item Notions de base sur la citoyenneté et les droits de l'homme (Chapitre 4, leçons antérieures).
\item Technique de l'introduction en entonnoir (général → spécifique → thèse + plan).
\end{itemize}
\end{prerequis}

\newpage

\coursec{Vocabulaire thématique : Elections \& Campaigning}

Au Cameroun, comme au Nigeria, au Ghana, au Royaume-Uni ou aux États-Unis, les périodes d'élections animent les débats dans les maquis, les marchés et sur les réseaux sociaux. Chaque citoyen a son avis sur le candidat idéal, mais peu savent structurer leurs arguments par écrit pour convaincre un jury d'examen. Imaginez que vous deviez rédiger une tribune pour le journal \eng{Cameroon Tribune} ou \eng{The Guardian Post} afin d'expliquer pourquoi le vote est un devoir citoyen ou comment une campagne doit se dérouler. Cette leçon vous donne les clés pour transformer vos opinions en une composition structurée, cohérente et grammaticalement correcte, digne du Baccalauréat.

\courssub{Les acteurs du jeu électoral (People)}

Avant de parler d'actions, il faut identifier \textbf{qui} fait quoi. L'anglais distingue avec précision les rôles, là où le français utilise parfois le même mot (« candidat », « électeur »).

\begin{textebox}[Extrait de manifesto : « A New Dawn for Our Constituency »]
Our \textbf{manifesto} pledges to fight \textbf{corruption} and ensure \textbf{transparency} in public funds. We \textbf{urge} every \textbf{eligible voter} to \textbf{cast their ballot} on \textbf{election day}. The \textbf{incumbent} has failed his \textbf{constituency}; it is time for a \textbf{running mate} who understands the \textbf{grassroots}. Let us secure a \textbf{landslide victory} for accountability.
\end{textebox}

\begin{definition}[Acteurs clés : Vocabulaire de base]
\begin{itemize}
    \item \cle{Candidate} (n.) : Candidat (ne prend jamais de « e » à la fin, contrairement au français « candidat(e) »).
    \item \cle{Voter} (n.) : Électeur (personne qui vote), synonyme : \cle{elector}.
    \item \cle{Electorate} (n.) : Corps électoral / Électorat (l'\textbf{ensemble} des électeurs inscrits). \textbf{Attention :} ne désigne pas une personne seule.
    \item \cle{Constituency} (n.) : Circonscription (territoire) \textbf{ou} ensemble des électeurs d'une circonscription.
    \item \cle{Incumbent} (n./adj.) : Sortant (député, président en fonction). Ex: \eng{The incumbent president}.
    \item \cle{Running mate} (n.) : Colistier (candidat à la vice-présidence ou suppléant).
    \item \cle{Pollster} (n.) : Sondeur (spécialiste des sondages d'opinion).
\end{itemize}
\end{definition}

\begin{center}
\begin{tabularx}{\linewidth}{|X|X|X|}
\hline
\rowcolor{popblueL} \textbf{English} & \textbf{Français} & \textbf{Note / Contexte camerounais} \\
\hline
Candidate & Candidat & Invariable en genre (pas de *candidate* pour une femme). \\
Voter / Elector & Électeur (individu) & \eng{Registered voters} = électeurs inscrits. \\
Electorate & Corps électoral (collectif) & \eng{The electorate is angry.} = L'électorat est en colère. \\
Constituency & Circonscription / Électeurs locaux & \eng{MP for the Northwest constituency.} \\
Incumbent & Sortant (titulaire du poste) & \eng{The incumbent mayor of Douala.} \\
Running mate & Colistier / Vice-président & \eng{She chose a youth leader as running mate.} \\
Pollster & Sondeur & \eng{Pollsters predict a tight race.} \\
\hline
\end{tabularx}
\end{center}

\courssub{Le processus électoral : actions, lieux et résultats (Process)}

Une campagne suit une chronologie précise. Maîtriser ces verbes et noms permet de narrer les faits (« reporting ») ou d'argumenter (« arguing »).

\begin{definition}[Processus : Verbes et noms indispensables]
\begin{itemize}
    \item \cle{To campaign} (v.) : Faire campagne. \eng{He campaigns tirelessly.}
    \item \cle{To canvass} (v.) : Faire du porte-à-porte, solliciter les voix activement. \eng{Volunteers canvass the neighbourhood.}
    \item \cle{To cast a ballot} (v. phrase) : Déposer son bulletin dans l'urne (locution figée).
    \item \cle{Polling station} (n.) : Bureau de vote (UK). US : \cle{polling place}.
    \item \cle{Ballot box} (n.) : Urne (litt. « boîte à bulletins »).
    \item \cle{Turnout} (n.) : Taux de participation (nombre/votants inscrits). \eng{High/Low turnout.}
    \item \cle{Landslide victory} (n.) : Victoire écrasante / à une large majorité.
    \item \cle{Runoff} (n.) / \cle{Run-off election} : Second tour (ballotage).
    \item \cle{Manifesto} (n. UK) / \cle{Platform} (n. US) : Programme électoral écrit.
\end{itemize}
\end{definition}

\begin{propriete}[Différenciation cruciale : Vote, Ballot, Poll]
Ces trois mots se traduisent souvent par « vote » ou « scrutin » en français, mais leur emploi diffère :
\begin{center}
\begin{tabular}{|l|l|l|}
\hline
\rowcolor{poporangeL} \textbf{Mot} & \textbf{Sens principal} & \textbf{Exemple traduit} \\
\hline
\cle{Vote} (n./v.) & Acte de voter / Suffrage / Résultat & \eng{The vote takes place Sunday.} (Le scrutin a lieu dimanche.) \\
\cle{Ballot} (n.) & Bulletins de vote \textbf{OU} Scrutin secret & \eng{A secret ballot.} (Un scrutin secret / vote à bulletin secret.) \\
\cle{Poll} (n.) & Sondage \textbf{OU} Bureau de vote (au pluriel \eng{the polls}) & \eng{Opinion polls show...} (Les sondages montrent...) / \eng{The polls open at 8am.} (Les bureaux ouvrent à 8h.) \\
\hline
\end{tabular}
\end{center}
\end{propriete}

\begin{exemplebox}[Exemples contextualisés]
\eng{The \textbf{turnout} in the Northwest region was historically low.} \\
→ La \textbf{participation} dans la région du Nord-Ouest était historiquement faible.

\eng{Voters queued for hours at the \textbf{polling station} to \textbf{cast their ballots}.} \\
→ Les électeurs ont fait la queue des heures au \textbf{bureau de vote} pour \textbf{glisser leur bulletin dans l'urne}.

\eng{The opposition claims the \textbf{incumbent} used state resources to \textbf{campaign}.} \\
→ L'opposition affirme que le \textbf{sortant} a utilisé des ressources de l'État pour \textbf{faire campagne}.
\end{exemplebox}

\courssub{Les enjeux et le vocabulaire de l'argumentation (Issues)}

Pour écrire une composition, il faut nommer les concepts abstraits : transparence, devoir, promesses.

\begin{definition}[Enjeux : Noms abstraits et verbes d'action civique]
\begin{itemize}
    \item \cle{Manifesto / Platform} : Programme électoral (promesses écrites).
    \item \cle{Pledge} (n./v.) : Engagement solennel, promesse formelle. \eng{He pledged to build roads.}
    \item \cle{Slogan} (n.) : Mot d'ordre, slogan de campagne.
    \item \cle{Rally} (n.) : Meeting politique, rassemblement populaire.
    \item \cle{Debate} (n.) : Débat télévisé ou public entre candidats.
    \item \cle{Corruption} (n.) : Corruption (faux ami : prononciation « ko-rup-chn », pas « kor-rup-sion »).
    \item \cle{Transparency} (n.) : Transparence (gestion claire des fonds).
    \item \cle{Accountability} (n.) : Redevabilité / Obligation de rendre des comptes (concept clé en gouvernance).
    \item \cle{Civic duty} (n.) : Devoir civique / citoyen.
    \item \cle{Eligible} (adj.) : Ayant le droit (âge, nationalité, inscription). \eng{Eligible voters.}
\end{itemize}
\end{definition}

\begin{propriete}[Collocations fortes (Blocs lexicaux figés)]
Ces associations de mots sont \textbf{obligatoires} pour un anglais naturel et noté au Bac. Ne les séparez pas.
\begin{center}
\begin{tabularx}{\linewidth}{|X|X|X|}
\hline
\rowcolor{popgreenL} \textbf{Collocation} & \textbf{Traduction} & \textbf{Exemple d'utilisation} \\
\hline
\cle{Free and fair elections} & Élections libres et transparentes & \eng{Observers monitor \textbf{free and fair elections}.} \\
\cle{Civic duty} & Devoir civique & \eng{Voting is a \textbf{civic duty}, not a favour.} \\
\cle{Grassroots campaign} & Campagne de terrain / de base & \eng{A strong \textbf{grassroots campaign} wins votes.} \\
\cle{To run for office} & Se présenter / briguer un mandat & \eng{He decided to \textbf{run for office} at 35.} \\
\cle{To win by a landslide} & Gagner haut la main / écrasante majorité & \eng{The party \textbf{won by a landslide}.} \\
\cle{Voter apathy} & Abstentionnisme / Désintérêt électoral & \eng{\textbf{Voter apathy} threatens democracy.} \\
\cle{To hold accountable} & Demander des comptes à & \eng{We must \textbf{hold leaders accountable}.} \\
\hline
\end{tabularx}
\end{center}
\end{propriete}

\courssub{Familles de mots (Word Formation) : \eng{Elect}, \eng{Campaign}, \eng{Vote}}

Au Bac, l'épreuve de « Use of English » ou la rédaction exige de transformer les mots. Voici les dérivations régulières.

\begin{center}
\begin{tabular}{|l|l|l|l|}
\hline
\rowcolor{poppurpleL} \textbf{Racine (Verbe)} & \textbf{Nom (Action/État)} & \textbf{Nom (Agent)} & \textbf{Adjectif} \\
\hline
\cle{Elect} & \cle{Election} & \cle{Elector} & \cle{Electoral} (adj.) \\
& \cle{Re-election} & & \cle{Elective} (éligible/électif) \\
\hline
\cle{Campaign} & \cle{Campaign} (n.) & \cle{Campaigner} & \cle{Campaign} (adj.: \eng{campaign trail}) \\
\hline
\cle{Vote} & \cle{Vote} (n.) & \cle{Voter} & \cle{Voting} (adj.: \eng{voting age}, \eng{voting card}) \\
\hline
\end{tabular}
\end{center}

\begin{attention}[Pièges de formation de mots (Erreurs Bac fréquentes)]
\begin{itemize}
    \item \textbf{*Electorial} $\rightarrow$ Faux. Correct : \cle{Electoral} (un seul « l » après « t »).
    \item \textbf{*Electer} $\rightarrow$ Faux. L'agent est \cle{Elector} (ou \cle{Voter}). \eng{To elect} est le verbe.
    \item \textbf{*Campaigner} (verbe) $\rightarrow$ Faux. \eng{Campaigner} est le \textbf{nom} (le militant). Le verbe est \cle{To campaign}.
    \item \textbf{*Votant} $\rightarrow$ N'existe pas en anglais. Participe présent : \eng{Voting} (utilisé comme adj. ou gérondif). L'agent : \cle{Voter}.
\end{itemize}
\end{attention}

\courssub{Mind-map lexicale : Visualiser le champ sémantique}

\begin{popfigure}
\begin{tikzpicture}[
    mindmap,
    grow cyclic,
    every node/.style={concept, align=center},
    concept color=popblue,
    level 1/.append style={level distance=4.5cm, sibling angle=90},
    level 2/.append style={level distance=3cm, sibling angle=45},
    text=white,
    font=\small\bfseries
]
\node {ELECTIONS}
    child[concept color=poporange] { node {PEOPLE}
        child { node {Candidate} }
        child { node {Voter / Elector} }
        child { node {Electorate} }
        child { node {Incumbent} }
        child { node {Running Mate} }
        child { node {Pollster} }
    }
    child[concept color=popgreen] { node {ACTIONS}
        child { node {Campaign / Canvass} }
        child { node {Cast a Ballot} }
        child { node {Run for Office} }
        child { node {Hold Accountable} }
        child { node {Pledge} }
    }
    child[concept color=poppurple] { node {CONCEPTS}
        child { node {Democracy} }
        child { node {Manifesto / Platform} }
        child { node {Transparency} }
        child { node {Accountability} }
        child { node {Civic Duty} }
        child { node {Turnout} }
    }
    child[concept color=poppink] { node {PLACE / TIME}
        child { node {Polling Station} }
        child { node {Ballot Box} }
        child { node {Election Day} }
        child { node {Runoff (2nd Round)} }
        child { node {Rally} }
    };
\end{tikzpicture}
\legende{Carte mentale du champ lexical « Elections \& Campaigning » — À mémoriser par branches.}
\end{popfigure}

\courssub{Exercice résolu : Du vocabulaire à la phrase}

\begin{exoresolu}[Complétion et traduction]
\textbf{Énoncé :} Complétez les phrases avec le mot approprié (forme correcte) puis traduisez.
\begin{enumerate}
    \item The \_\_\_\_\_\_\_\_\_\_ (elect) commission announced the results yesterday.
    \item She is a seasoned \_\_\_\_\_\_\_\_\_\_ (campaign) who has managed three presidential races.
    \item Low \_\_\_\_\_\_\_\_\_\_ (turnout) is often a sign of \_\_\_\_\_\_\_\_\_\_ (vote) apathy.
    \item The \_\_\_\_\_\_\_\_\_\_ (incumbent) deputy decided not to \_\_\_\_\_\_\_\_\_\_ (run) for re-election.
    \item Citizens must \_\_\_\_\_\_\_\_\_\_ (hold) their leaders \_\_\_\_\_\_\_\_\_\_ (accountable) for public spending.
\end{enumerate}
\tcblower
\textbf{Solution détaillée :}
\begin{enumerate}
    \item \textbf{Electoral} (adj. qualificatif de \eng{commission}). \textit{Règle :} \eng{Elect} $\rightarrow$ \eng{Electoral}. \\
    \eng{The \textbf{electoral} commission announced the results yesterday.} \\
    → La commission \textbf{électorale} a annoncé les résultats hier.
    \item \textbf{Campaigner} (nom d'agent). \textit{Piège :} pas « campaignist ». \\
    \eng{She is a seasoned \textbf{campaigner} who has managed three presidential races.} \\
    → C'est une \textbf{militante / stratège de campagne} chevronnée qui a géré trois courses présidentielles.
    \item \textbf{Turnout} (nom invariable) / \textbf{Voter} (adjectif attributif : \eng{voter apathy}). \textit{Note :} \eng{Voting apathy} est possible mais \eng{Voter apathy} est la collocation standard. \\
    \eng{Low \textbf{turnout} is often a sign of \textbf{voter} apathy.} \\
    → Une faible \textbf{participation} est souvent le signe d'un \textbf{désintérêt électoral}.
    \item \textbf{Incumbent} (adj. ou nom) / \textbf{Run} (verbe infinitif après \eng{to}). \textit{Note :} \eng{Run for re-election} = briguer un nouveau mandat. \\
    \eng{The \textbf{incumbent} deputy decided not to \textbf{run} for re-election.} \\
    → Le député \textbf{sortant} a décidé de ne pas \textbf{briguer} un nouveau mandat.
    \item \textbf{Hold} (infinitif après \eng{must}) / \textbf{Accountable} (adjectif, collocation \eng{hold sb accountable}). \\
    \eng{Citizens must \textbf{hold} their leaders \textbf{accountable} for public spending.} \\
    → Les citoyens doivent \textbf{demander des comptes à} leurs dirigeants pour les dépenses publiques.
\end{enumerate}
\end{exoresolu}

\courssub{Le saviez-vous ? : Culture électorale comparée}

\begin{savaistu}[Jour de vote : Une question de tradition]
\begin{itemize}
    \item \textbf{Royaume-Uni} : On vote traditionnellement un \textbf{jeudi} (\eng{Thursday}) depuis 1935 (pas de loi stricte, mais coutume). Les bureaux ouvrent de 7h à 22h.
    \item \textbf{États-Unis} : \eng{Election Day} est fixé par la loi au \textbf{mardi suivant le premier lundi de novembre} (années paires). Héritage du XIXe siècle (marchés le mercredi, dimanche église, lundi trajet).
    \item \textbf{Cameroun} : L'élection présidentielle a lieu un \textbf{dimanche} (souvent en octobre, tous les 7 ans). Les législatives et municipales sont souvent couplées. Le jour est décrété férié par décret présidentiel pour faciliter la participation.
    \item \textbf{Australie} : Le vote est \textbf{obligatoire} (\eng{compulsory voting}) pour les inscrits ; amende en cas d'abstention non justifiée. Le taux de participation y dépasse 90\%.
\end{itemize}
\end{savaistu}

\begin{exercice}[Entraînement personnel : Vocabulaire actif]
\textbf{1. Traduisez en anglais (utilisez les collocations de la leçon) :}
\begin{enumerate}
    \item Le taux de participation a été décevant lors du second tour.
    \item Les observateurs internationaux exigent des élections libres et transparentes.
    \item Le candidat sortant a promis la transparence dans son programme.
    \item Faire du porte-à-porte est essentiel pour une campagne de terrain.
    \item C'est un devoir civique de glisser son bulletin dans l'urne.
\end{enumerate}
\textbf{2. Complétez le tableau de dérivation (forme nom, adjectif, agent) :}
\begin{center}
\begin{tabular}{|l|c|c|c|}
\hline
\rowcolor{popblueL} Verbe & Nom (Action) & Nom (Agent) & Adjectif \\
\hline
To elect & & & \\
\hline
To campaign & & & \\
\hline
To vote & & & \\
\hline
\end{tabular}
\end{center}
\textbf{3. Corrigez les erreurs dans ces phrases d'élèves :}
\begin{enumerate}
    \item The *electorial* commission is independent.
    \item He *campaigners* well for his party.
    \item The *voters* (pour désigner le corps électoral) are angry.
    \item She *casted* her ballot at 10am.
    \item We need *fair and free* elections.
\end{enumerate}
\end{exercice}

\begin{corrige}[Exercice 1.7]
\textbf{1. Traduction :}
\begin{enumerate}
    \item \eng{Voter turnout was disappointing in the runoff.} (ou \eng{turnout})
    \item \eng{International observers demand free and fair elections.}
    \item \eng{The incumbent candidate pledged transparency in his manifesto.} (ou \eng{platform})
    \item \eng{Canvassing is essential for a grassroots campaign.}
    \item \eng{Casting one's ballot is a civic duty.} (ou \eng{Voting is a civic duty.})
\end{enumerate}
\textbf{2. Tableau :}
\begin{center}
\begin{tabular}{|l|c|c|c|}
\hline
\rowcolor{popblueL} Verbe & Nom (Action) & Nom (Agent) & Adjectif \\
\hline
To elect & Election / Re-election & Elector & Electoral / Elective \\
\hline
To campaign & Campaign & Campaigner & Campaign (adj.) \\
\hline
To vote & Vote & Voter & Voting \\
\hline
\end{tabular}
\end{center}
\textbf{3. Corrections :}
\begin{enumerate}
    \item \eng{Electoral} (un seul « l »).
    \item \eng{He campaigns well for his party.} (Verbe \eng{to campaign} à la 3e pers. sg ; \eng{campaigner} est un nom).
    \item \eng{The electorate} (collectif) \eng{is angry.}
    \item \eng{Cast} (verbe irrégulier : cast-cast-cast). \eng{She cast her ballot.}
    \item \eng{Free and fair elections} (collocation figée, ordre immuable).
\end{enumerate}
\end{corrige}

\coursec{Structure type de la composition argumentative}

Dans la section précédente, nous avons construit un solide vocabulaire thématique autour des élections et des campagnes (\eng{candidate, polling station, ballot paper, to cast a vote, manifesto}...). Mais posséder les mots ne suffit pas : au baccalauréat, un élève qui connaît cent mots d'anglais mais qui écrit un texte sans plan obtient une note médiocre. Ce que le correcteur cherche d'abord, c'est la \cle{structure}. Cette section vous donne le « squelette » de toute composition argumentative réussie : l'introduction, le développement et la conclusion, avec leurs sous-parties obligatoires.

\courssub{Observation : deux compositions, deux destins}

Lisons d'abord deux courts extraits de copies sur le sujet \eng{“Voting is a duty, not a choice.” Discuss}. Observez et comparez.

\begin{textebox}[Candidate A — a well-structured opening]
“Elections are the heartbeat of democracy.” This popular saying reminds us that without regular voting, no government can claim to represent its people. In Cameroon, as in many African countries, citizens go to the polls to choose their mayors, parliamentarians and president. However, some young people believe that voting is useless. This essay argues that voting is a duty, not a choice, because it legitimizes power and because it holds leaders accountable.
\end{textebox}

\begin{textebox}[Candidate B — a confused opening]
Voting is very important. Many people vote in many countries of the world. My uncle voted last year and he was very happy. There are many political parties and they do campaigns on television and radio. Democracy is good. I think everybody should vote because it is good for the country and also for the family and the future of children.
\end{textebox}

\textbf{Questions d'observation}
\begin{enumerate}
\item Which candidate announces a clear position (a thesis)? Underline it.
\item Which candidate announces the arguments that will follow?
\item Which introduction wanders from idea to idea without direction?
\end{enumerate}

Le texte A suit un chemin précis : une \cle{accroche} (\eng{hook}), un \cle{contexte} (\eng{background}), une \cle{thèse} (\eng{thesis statement}) et un \cle{plan} (\eng{outline}). Le texte B accumule des phrases vraies mais désorganisées : c'est l'erreur la plus fréquente des candidats francophones, qui « commencent à écrire » au lieu de « construire ».

\courssub{Vue d'ensemble : les trois parties obligatoires}

\begin{definition}[La composition argumentative (argumentative essay)]
Une \cle{composition argumentative} est un texte organisé en trois parties inséparables : l'\cle{introduction} (qui présente le sujet et la position de l'auteur), le \cle{développement} ou \eng{body} (qui défend cette position avec 2 ou 3 arguments illustrés) et la \cle{conclusion} (qui reformule la position et ouvre une perspective). Aucune de ces trois parties ne peut être supprimée sans pénalité lourde à l'examen.
\end{definition}

\begin{aretenir}[L'équilibre des parties — la règle des pourcentages]
Pour une composition du baccalauréat (cible : \cle{250 à 300 mots}), respectez cette répartition :
\begin{itemize}
\item \textbf{Introduction} : environ 10 à 15 \% du total (30 à 45 mots, un seul paragraphe).
\item \textbf{Développement} : environ 70 à 80 \% du total (180 à 230 mots, 2 ou 3 paragraphes).
\item \textbf{Conclusion} : environ 10 à 15 \% du total (30 à 45 mots, un seul paragraphe).
\end{itemize}
Une introduction d'une page et un développement de trois lignes est une copie déséquilibrée, donc mal notée.
\end{aretenir}

\begin{popfigure}
\begin{tikzpicture}[
  node distance=0.55cm,
  mainbox/.style={draw=popblue, thick, fill=popblueL, rounded corners=3pt, text width=11.5cm, align=left, inner sep=6pt},
  bodybox/.style={draw=poporange, thick, fill=poporangeL, rounded corners=3pt, text width=3.5cm, align=left, inner sep=5pt},
  conclbox/.style={draw=popgreen, thick, fill=popgreenL, rounded corners=3pt, text width=11.5cm, align=left, inner sep=6pt},
  arr/.style={-{Stealth[length=3mm]}, thick, popdark}
]
\node[mainbox, align=center] (intro){\textbf{INTRODUCTION (10--15 \%)}\\ \textbf{1.} Hook (accroche) \; \textbf{2.} Background (contexte) \; \textbf{3.} Thesis statement (position) \; \textbf{4.} Outline (plan)};
\node[bodybox, below=1.2cm of intro, align=center] (b2){\textbf{BODY 2}\\ Topic sentence\\ Example / Evidence\\ Explanation\\ Link};
\node[bodybox, left=0.4cm of b2, align=center] (b1){\textbf{BODY 1}\\ Topic sentence\\ Example / Evidence\\ Explanation\\ Link};
\node[bodybox, right=0.4cm of b2, align=center] (b3){\textbf{BODY 3 (optionnel)}\\ Topic sentence\\ Example / Evidence\\ Explanation\\ Link};
\node[conclbox, below=1.2cm of b2, align=center] (concl){\textbf{CONCLUSION (10--15 \%)}\\ \textbf{1.} Restate the thesis (reformulation) \; \textbf{2.} Summarize the main points \; \textbf{3.} Final thought / call to action};
\draw[arr] (intro.south) -- (b1.north);
\draw[arr] (intro.south) -- (b2.north);
\draw[arr] (intro.south) -- (b3.north);
\draw[arr] (b1.south) -- (concl.north);
\draw[arr] (b2.south) -- (concl.north);
\draw[arr] (b3.south) -- (concl.north);
\end{tikzpicture}
\legende{Organigramme de la structure d'une composition argumentative : de l'introduction en entonnoir vers les paragraphes de développement, puis la conclusion.}
\end{popfigure}

\courssub{L'introduction : la technique de l'entonnoir}

L'introduction fonctionne comme un \cle{entonnoir} (\eng{funnel}) : on part du général pour arriver au précis, c'est-à-dire à votre position personnelle. Elle contient quatre éléments, dans cet ordre.

\textbf{1. Le hook (l'accroche).} C'est la toute première phrase : elle doit capter l'attention du lecteur. Trois techniques efficaces : une \cle{citation}, un \cle{fait d'actualité} ou une \cle{statistique plausible}, une \cle{question rhétorique}.

\textbf{2. Le background (le contexte).} Deux ou trois phrases qui situent le sujet et définissent les termes clés. Ici, on peut évoquer le Cameroun, l'Afrique ou le monde anglophone.

\textbf{3. La thesis statement (la thèse).} C'est la phrase la plus importante de toute la copie : votre \cle{prise de position} claire et nuancée.

\textbf{4. L'outline (l'annonce de plan).} Une phrase qui annonce les deux ou trois arguments à venir. En anglais, elle est souvent intégrée à la thèse grâce à \eng{because} ou \eng{as}.

\begin{exemplebox}[Hooks et thèses : exemples traduits]
\begin{itemize}
\item \eng{“Democracy is the worst form of government, except for all the others,” Winston Churchill once said.} — « La démocratie est le pire des régimes, à l'exception de tous les autres », disait Winston Churchill. (hook par citation)
\item \eng{Every election year, thousands of young Cameroonians queue at polling stations to choose their future.} — Chaque année électorale, des milliers de jeunes Camerounais font la queue devant les bureaux de vote pour choisir leur avenir. (hook par fait d'actualité)
\item \eng{Although campaigning can be divisive, voting remains the cornerstone of democracy.} — Bien que la campagne puisse diviser, le vote reste la pierre angulaire de la démocratie. (thesis statement nuancée avec \eng{although})
\item \eng{Voting is a civic duty because it legitimizes power and holds leaders accountable.} — Voter est un devoir civique parce que cela légitime le pouvoir et rend les dirigeants responsables. (thèse + plan intégré : deux arguments annoncés)
\end{itemize}
\end{exemplebox}

\begin{methode}[Rédiger une thesis statement efficace en trois étapes]
\begin{enumerate}
\item \textbf{Identifiez le sujet} (\eng{topic}) : de quoi parle le sujet ? Ex. : \eng{voting}.
\item \textbf{Ajoutez votre opinion} (\eng{opinion}) : quelle est votre position ? Ex. : \eng{is a civic duty}.
\item \textbf{Ajoutez les mots-clés de vos arguments} avec \eng{because}, \eng{as} ou \eng{since} : cela annonce le plan. Ex. : \eng{because it legitimizes power and holds leaders accountable}.
\end{enumerate}
Formule complète : \textbf{Sujet + Opinion + Mots-clés des arguments}.\\
Résultat : \eng{Voting is a civic duty because it legitimizes power and holds leaders accountable.} Le correcteur sait déjà que le corps aura deux paragraphes : la légitimation du pouvoir, puis la responsabilisation des dirigeants.
\end{methode}

\courssub{Le développement : un paragraphe = une idée force}

Le développement (\eng{body}) comporte \cle{deux ou trois paragraphes distincts}. Règle d'or : \textbf{chaque paragraphe défend une seule idée force}. Si vous mélangez deux arguments dans le même paragraphe, vous perdez le lecteur — et des points.

\begin{definition}[Topic sentence et transition]
La \cle{topic sentence} (phrase sujet) est la première phrase d'un paragraphe : elle annonce clairement l'idée principale du paragraphe. Ex. : \eng{First of all, voting gives legitimacy to those who govern.}\\
Une \cle{transition} est un mot ou une phrase qui relie deux paragraphes et guide le lecteur : \eng{Furthermore,} (de plus), \eng{On the other hand,} (d'autre part), \eng{Consequently,} (par conséquent).
\end{definition}

\begin{propriete}[La structure PEEL d'un paragraphe de développement]
Chaque paragraphe du corps suit le schéma \cle{PEEL} :
\begin{itemize}
\item \textbf{P — Point} : la topic sentence, qui énonce l'argument. \eng{First of all, voting legitimizes power.}
\item \textbf{E — Evidence / Example} : une preuve, un fait, une statistique plausible ou un exemple de pays. \eng{In Ghana in 2016, a peaceful transfer of power followed the presidential election, which strengthened citizens' trust.} ou \eng{Statistics show that countries with high voter turnout enjoy greater political stability.}
\item \textbf{E — Explanation} : l'analyse — pourquoi cet exemple prouve-t-il l'argument ? \eng{This shows that when people vote massively, leaders cannot be accused of stealing power.}
\item \textbf{L — Link} : la transition vers le paragraphe suivant. \eng{Beyond legitimacy, voting also serves another essential purpose.}
\end{itemize}
Un paragraphe PEEL complet fait généralement 70 à 90 mots.
\end{propriete}

\begin{exemplebox}[Un paragraphe PEEL complet, décortiqué]
\eng{First of all, voting gives legitimacy to elected leaders.} (\textbf{P}) \eng{For instance, when a large majority of citizens cast their ballots, no one can seriously contest the result.} (\textbf{E}) \eng{In other words, the ballot box transforms individual opinions into collective authority: a president chosen by millions of voters speaks in the name of the whole nation.} (\textbf{E}) \eng{However, legitimacy alone is not enough; leaders must also be controlled.} (\textbf{L})\\
Traduction : « Tout d'abord, voter donne de la légitimité aux dirigeants élus. Par exemple, lorsqu'une large majorité de citoyens dépose son bulletin, personne ne peut sérieusement contester le résultat. En d'autres termes, les urnes transforment des opinions individuelles en autorité collective : un président choisi par des millions d'électeurs parle au nom de la nation entière. Cependant, la légitimité seule ne suffit pas ; les dirigeants doivent aussi être contrôlés. »
\end{exemplebox}

\courssub{La conclusion : l'entonnoir inversé}

La conclusion suit le mouvement inverse de l'introduction : on part du précis (votre thèse) pour revenir au général (une ouverture). Elle comporte trois éléments.

\textbf{1. Restatement of the thesis (reformulation).} Répétez votre position \emph{avec d'autres mots} — jamais un copier-coller de l'introduction. Si la thèse était \eng{Voting is a civic duty}, la reformulation peut être \eng{In conclusion, casting a ballot is far more than a right: it is a responsibility every citizen must assume.}

\textbf{2. Summary of the main points (synthèse).} Rappelez en une phrase vos deux ou trois arguments : \eng{As shown above, elections both legitimize power and force leaders to answer for their actions.}

\textbf{3. Final thought / opening (ouverture).} Terminez par une réflexion, un appel à l'action ou une perspective d'avenir : \eng{If every young Cameroonian voted, the future of the nation would truly belong to its youth.} — « Si chaque jeune Camerounais votait, l'avenir de la nation appartiendrait vraiment à sa jeunesse. »

\begin{attention}[Les trois erreurs fatales de structure]
\begin{itemize}
\item \textbf{Introduction trop longue} (plus d'un paragraphe, des exemples dès l'introduction) ou \textbf{sans thesis statement} : le correcteur ne sait pas ce que vous défendez.
\item \textbf{Conclusion qui introduit un nouvel argument} : la conclusion ne contient \emph{aucune} idée neuve ; elle reformule et ouvre, c'est tout.
\item \textbf{Paragraphes d'une seule phrase} : un paragraphe de développement contient au minimum 4 à 5 phrases (le schéma PEEL complet).
\end{itemize}
\end{attention}

\courssub{Le paragraphing : matérialiser la structure sur la copie}

La structure doit être \cle{visible} sur la copie. En anglais comme en français, on marque le changement de paragraphe de deux façons :
\begin{itemize}
\item soit une \cle{ligne blanche} entre les paragraphes ;
\item soit un \cle{retrait} (indentation) en début de première ligne, avec un léger saut de ligne.
\end{itemize}
Au baccalauréat, la ligne blanche est la solution la plus sûre : le correcteur voit immédiatement vos quatre ou cinq blocs (introduction, 2 ou 3 paragraphes de corps, conclusion). Ne numérotez jamais vos paragraphes et n'utilisez pas de puces : une composition est un texte continu, pas une liste.

\begin{center}
\begin{tabularx}{\linewidth}{|X|X|X|}
\hline
\rowcolor{popblueL} \textbf{Partie} & \textbf{Contenu obligatoire} & \textbf{Expressions utiles} \\
\hline
Introduction & Hook + background + thesis + plan & \eng{It is often said that... / This essay argues that...} \\
\hline
Body 1 & Argument 1 (PEEL) & \eng{First of all, / For instance, / This shows that...} \\
\hline
Body 2 & Argument 2 (PEEL) & \eng{Furthermore, / A striking example is... / In other words,...} \\
\hline
Body 3 (optionnel) & Contre-argument réfuté ou argument 3 & \eng{Admittedly,... however... / On the other hand,...} \\
\hline
Conclusion & Reformulation + synthèse + ouverture & \eng{In conclusion, / To sum up, / If... then...} \\
\hline
\end{tabularx}
\end{center}

\begin{savaistu}[L'entonnoir, une tradition anglo-saxonne]
La structure « hook — thesis — PEEL — conclusion » vient de la tradition rhétorique anglo-saxonne, où l'on annonce sa position \emph{dès le début} (\eng{get to the point}). La dissertation française, elle, préfère souvent garder la réponse pour la fin. Au baccalauréat d'anglais, c'est le modèle anglo-saxon qui est exigé : soyez direct, dites votre thèse dans l'introduction !
\end{savaistu}

\begin{exoresolu}[Remettre une introduction dans l'ordre]
\textbf{Énoncé.} Les quatre phrases suivantes forment une introduction, mais elles sont mélangées. Remettez-les dans l'ordre logique (hook → background → thesis → outline).\\
(a) \eng{This essay argues that voting is essential because it legitimizes leaders and protects citizens' rights.}\\
(b) \eng{In many African countries, including Cameroon, citizens regularly elect their mayors, members of parliament and president.}\\
(c) \eng{“The ballot is stronger than the bullet,” Abraham Lincoln famously declared.}\\
(d) \eng{Yet, in spite of this long tradition, many young people still refuse to register on electoral lists.}
\tcblower
\textbf{Solution détaillée.} L'ordre correct est \textbf{(c) → (b) → (d) → (a)}.\\
\textbf{(c)} est le \emph{hook} : une citation célèbre attire l'attention. \textbf{(b)} apporte le \emph{background} : le contexte camerounais et africain. \textbf{(d)} crée la tension qui justifie le sujet (\eng{yet} = pourtant) : c'est le problème à discuter. \textbf{(a)} est la \emph{thesis statement} avec plan intégré : \eng{because} annonce les deux arguments (légitimer les dirigeants, protéger les droits des citoyens).\\
Traduction de l'introduction reconstituée : « “Le bulletin de vote est plus fort que la balle”, déclara Abraham Lincoln. Dans de nombreux pays africains, dont le Cameroun, les citoyens élisent régulièrement leurs maires, députés et président. Pourtant, malgré cette longue tradition, beaucoup de jeunes refusent encore de s'inscrire sur les listes électorales. Cette dissertation soutient que voter est essentiel parce que cela légitime les dirigeants et protège les droits des citoyens. »
\end{exoresolu}

\begin{exercice}[À vous de jouer]
Sujet : \eng{“Young people should be more involved in choosing their leaders.”}\\
1. Rédigez une thesis statement complète (sujet + opinion + deux arguments annoncés avec \eng{because}).\\
2. Écrivez la topic sentence (P) de votre premier paragraphe de développement.\\
3. Proposez une phrase d'ouverture (final thought) pour votre conclusion.
\end{exercice}

\begin{corrige}[Exercice — éléments de réponse]
1. Exemple : \eng{Young people should take an active part in choosing their leaders because they represent the majority of the population and because they will live with the consequences of political decisions.} — « Les jeunes devraient participer activement au choix de leurs dirigeants parce qu'ils représentent la majorité de la population et qu'ils vivront avec les conséquences des décisions politiques. »\\
2. \eng{First of all, young voters form the largest age group in most African countries.} — « Tout d'abord, les jeunes électeurs forment le groupe d'âge le plus important dans la plupart des pays africains. »\\
3. \eng{If the youth of today vote tomorrow, the leaders of the future will finally speak their language.} — « Si la jeunesse d'aujourd'hui vote demain, les dirigeants du futur parleront enfin sa langue. »
\end{corrige}

\begin{attention}[Piège de traduction : « composition » et « développement »]
Attention aux faux amis ! En anglais scolaire, \eng{a composition} désigne bien une rédaction, mais on dit plus souvent \eng{an essay} pour une composition argumentative. Surtout, ne traduisez jamais « le développement » par \eng{the development} dans ce contexte : le terme correct est \eng{the body} (le corps du texte). \eng{Development} en anglais signifie « développement économique, croissance » (\eng{economic development}). De même, « un plan » au sens de « annonce de plan » se dit \eng{an outline}, pas \eng{a plan}.
\end{attention}

Vous disposez maintenant du squelette complet : introduction en entonnoir, paragraphes PEEL, conclusion en entonnoir inversé. La section suivante vous fournira les articulations de ce squelette : les connecteurs logiques et les temps verbaux qui donnent à votre composition sa fluidité et sa cohérence.

\coursec{Connecteurs logiques et Temps verbaux : Outils de cohérence}

Dans la section précédente, nous avons appris à organiser une composition argumentative en trois parties : introduction, développement et conclusion. Mais un bon plan ne suffit pas. Imaginez un mur de briques sans ciment : les briques sont bien empilées, mais le mur s'effondre au premier choc. Dans une composition, les \cle{linking words} (mots de liaison) sont ce ciment, et les \cle{temps verbaux} bien choisis sont le niveau du maçon. C'est précisément ce que le correcteur du Bac évalue sous la rubrique \eng{Coherence and Organisation}. Observons d'abord un court extrait de composition sur le vote, écrit par un élève de Terminale.

\begin{textebox}[Extract from a student's composition — observe the linking words]
Voting is a fundamental right in every democracy. \textbf{Moreover}, it is a duty that no responsible citizen should neglect. \textbf{Although} many young people believe that one vote cannot change anything, history proves the opposite. \textbf{For example}, in several African countries, elections \textbf{were won} by a very small number of votes. \textbf{Because} every ballot counts, citizens \textbf{must} register and vote. \textbf{However}, voting alone is not enough: voters \textbf{should} also choose their leaders carefully. \textbf{If} people voted wisely, corruption \textbf{would decrease}. \textbf{Therefore}, every citizen \textbf{should} take part in elections. \textbf{In conclusion}, voting matters because it shapes our future.
\end{textebox}

Glossaire : \emph{ballot} (n., bulletin de vote) ; \emph{to neglect} (v., négliger) ; \emph{wisely} (adv., sagement) ; \emph{to shape} (v., façonner, déterminer).

\textbf{Questions d'observation :} 1) Soulignez tous les mots de liaison. Quelle fonction remplit chacun d'eux ? 2) Quels temps verbaux et quels modaux remarquez-vous ? 3) Pourquoi l'auteur n'utilise-t-il jamais deux fois le même connecteur ?

\courssub{Les connecteurs logiques par fonction}

\begin{definition}[Linking words / Connecteurs logiques]
Un \cle{connecteur logique} (ou \emph{linking word}, \emph{linker}, \emph{transition word}) est un mot ou une expression qui relie deux idées et indique au lecteur la relation logique entre elles : ajout, opposition, cause, conséquence, illustration ou conclusion. Sans connecteurs, une composition est une simple liste de phrases ; avec eux, elle devient un raisonnement.
\end{definition}

Voici le tableau de référence à connaître par cœur pour le Bac. Apprenez au moins \textbf{trois connecteurs par fonction} afin de pouvoir varier.

\begin{center}
\begin{tabularx}{\linewidth}{|l|X|X|}
\hline
\rowcolor{popblueL} \textbf{Fonction} & \textbf{Connecteurs} & \textbf{Exemple} \\
\hline
Addition & \eng{Furthermore, Moreover, In addition, Besides, What is more, Also} & \eng{Moreover, voting is free and secret.} \\
\hline
Opposition / Concession & \eng{However, Nevertheless, On the other hand, Although, Even though, Despite, In spite of} & \eng{Although the campaign was long, voters stayed motivated.} \\
\hline
Cause / Raison & \eng{Because, Since, As, Because of, Due to} & \eng{Since voting is a right, it must be protected.} \\
\hline
Conséquence / Résultat & \eng{Therefore, Consequently, As a result, Thus, So} & \eng{Therefore, every vote counts.} \\
\hline
Illustration / Exemple & \eng{For example, For instance, Such as, Namely} & \eng{Some leaders, such as traditional rulers, influence voters.} \\
\hline
Conclusion & \eng{In conclusion, To sum up, Ultimately, All in all} & \eng{To sum up, voting is both a right and a duty.} \\
\hline
\end{tabularx}
\end{center}

\begin{exemplebox}[Exemples traduits]
\eng{Furthermore, the government should educate young voters.} « De plus, le gouvernement devrait éduquer les jeunes électeurs. »

\eng{On the other hand, some candidates make promises they cannot keep.} « D'un autre côté, certains candidats font des promesses qu'ils ne peuvent pas tenir. »

\eng{As a result, many citizens lost confidence in the process.} « Par conséquent, de nombreux citoyens ont perdu confiance dans le processus. »

\eng{For instance, in Buea, students organised a debate on elections.} « Par exemple, à Buea, des étudiants ont organisé un débat sur les élections. »

\eng{Ultimately, the choice belongs to the people.} « En fin de compte, le choix appartient au peuple. »
\end{exemplebox}

\begin{popfigure}
\begin{center}
\begin{tikzpicture}[font=\small]
\node[fill=popblue, text=white, rounded corners, minimum width=2.6cm, minimum height=0.9cm, align=center] (f1) at (0,4.5) {\textbf{FONCTION}};
\node[fill=popblue, text=white, rounded corners, minimum width=5.2cm, minimum height=0.9cm, align=center] (f2) at (4.4,4.5) {\textbf{CONNECTEURS}};
\node[fill=popblue, text=white, rounded corners, minimum width=5.2cm, minimum height=0.9cm, align=center] (f3) at (10.2,4.5) {\textbf{PONCTUATION / STRUCTURE}};
\node[fill=popgreenL, rounded corners, minimum width=2.6cm, minimum height=0.85cm, align=center] at (0,3.5) {Addition};
\node[fill=popgreenL!50, rounded corners, minimum width=5.2cm, minimum height=0.85cm, align=center] at (4.4,3.5) {Furthermore, Moreover, In addition, Besides};
\node[fill=popgreenL!50, rounded corners, minimum width=5.2cm, minimum height=0.85cm, align=center] at (10.2,3.5) {En tête de phrase + virgule};
\node[fill=poporangeL, rounded corners, minimum width=2.6cm, minimum height=0.85cm, align=center] at (0,2.55) {Opposition};
\node[fill=poporangeL!50, rounded corners, minimum width=5.2cm, minimum height=0.85cm, align=center] at (4.4,2.55) {However, Although, Despite, Nevertheless};
\node[fill=poporangeL!50, rounded corners, minimum width=5.2cm, minimum height=0.85cm, align=center] at (10.2,2.55) {However : . However, / ; however,};
\node[fill=poppurpleL, rounded corners, minimum width=2.6cm, minimum height=0.85cm, align=center] at (0,1.6) {Cause};
\node[fill=poppurpleL!50, rounded corners, minimum width=5.2cm, minimum height=0.85cm, align=center] at (4.4,1.6) {Because, Since, As, Due to};
\node[fill=poppurpleL!50, rounded corners, minimum width=5.2cm, minimum height=0.85cm, align=center] at (10.2,1.6) {Because + phrase ; Due to + nom};
\node[fill=poppinkL, rounded corners, minimum width=2.6cm, minimum height=0.85cm, align=center] at (0,0.65) {Conséquence};
\node[fill=poppinkL!50, rounded corners, minimum width=5.2cm, minimum height=0.85cm, align=center] at (4.4,0.65) {Therefore, Consequently, As a result, Thus};
\node[fill=poppinkL!50, rounded corners, minimum width=5.2cm, minimum height=0.85cm, align=center] at (10.2,0.65) {En tête de phrase + virgule};
\node[fill=popgoldL, rounded corners, minimum width=2.6cm, minimum height=0.85cm, align=center] at (0,-0.3) {Illustration};
\node[fill=popgoldL!50, rounded corners, minimum width=5.2cm, minimum height=0.85cm, align=center] at (4.4,-0.3) {For example, For instance, Such as};
\node[fill=popgoldL!50, rounded corners, minimum width=5.2cm, minimum height=0.85cm, align=center] at (10.2,-0.3) {For example, + phrase ; such as + nom};
\node[fill=popcream, rounded corners, minimum width=2.6cm, minimum height=0.85cm, align=center] at (0,-1.25) {Conclusion};
\node[fill=popcream!60, rounded corners, minimum width=5.2cm, minimum height=0.85cm, align=center] at (4.4,-1.25) {In conclusion, To sum up, Ultimately};
\node[fill=popcream!60, rounded corners, minimum width=5.2cm, minimum height=0.85cm, align=center] at (10.2,-1.25) {En tête de phrase + virgule};
\end{tikzpicture}
\end{center}
\legende{Tableau des connecteurs logiques par fonction, avec leur structure de ponctuation.}
\end{popfigure}

\courssub{Règles de ponctuation : trois familles de connecteurs}

Tous les connecteurs ne se ponctuent pas de la même manière. Il faut distinguer \textbf{trois familles grammaticales}.

\begin{propriete}[Famille 1 : les conjonctions de coordination — FANBOYS]
Les sept conjonctions de coordination anglaises se retiennent grâce au mot mnémotechnique \cle{FANBOYS} : \eng{For, And, Nor, But, Or, Yet, So}. Règle : quand elles relient \textbf{deux propositions indépendantes} (chacune avec sujet + verbe), on met une \textbf{virgule avant} la conjonction.

\eng{Citizens want change, so they vote in large numbers.} « Les citoyens veulent du changement, donc ils votent en grand nombre. »

\eng{The campaign was expensive, but it was peaceful.} « La campagne était coûteuse, mais elle était pacifique. »
\end{propriete}

\begin{propriete}[Famille 2 : les adverbes de liaison]
\eng{However, Therefore, Moreover, Furthermore, Consequently, Nevertheless} sont des \textbf{adverbes}, pas des conjonctions : ils ne peuvent pas relier deux propositions avec une simple virgule. Deux structures correctes :

1. \textbf{Point + majuscule + virgule} : \eng{Voting is a right. However, it is also a duty.}

2. \textbf{Point-virgule + adverbe + virgule} : \eng{Voting is a right; however, it is also a duty.}

La virgule après l'adverbe en début de phrase est \textbf{obligatoire}.
\end{propriete}

\begin{propriete}[Famille 3 : les prépositions et conjonctions de subordination]
\begin{itemize}
\item \eng{Because, Since, As, Although, Even though} = \textbf{conjonctions} : suivies d'une \textbf{proposition complète} (sujet + verbe). \eng{Although he was tired, he voted.}
\item \eng{Because of, Due to, Despite, In spite of} = \textbf{prépositions} : suivies d'un \textbf{nom} ou d'un \textbf{gérondif (V-ing)}, jamais d'une proposition complète. \eng{Despite the rain, voters queued.} / \eng{Because of the strike, the meeting was postponed.}
\end{itemize}
\end{propriete}

\begin{propriete}[However vs Although : le duel à maîtriser]
\cle{However} est un \textbf{adverbe} : il ouvre une \textbf{nouvelle phrase} (ou suit un point-virgule). \cle{Although} est une \textbf{conjonction de subordination} : il introduit une subordonnée \textbf{dans la même phrase}.

\eng{The election was difficult. However, people voted massively.} « L'élection fut difficile. Cependant, les gens ont voté massivement. »

\eng{Although the election was difficult, people voted massively.} « Bien que l'élection ait été difficile, les gens ont voté massivement. »

\eng{Despite} et \eng{In spite of} expriment la même concession, mais suivis d'un nom ou d'un V-ing : \eng{In spite of being busy, she went to the polling station.} « Bien qu'occupée, elle est allée au bureau de vote. »
\end{propriete}

\begin{attention}[Pièges fréquents des francophones]
\begin{itemize}
\item \textbf{Phrase fragmentée avec \eng{Because}} : \eng{Because voting is important.} seul est une \textbf{erreur} : la subordonnée attend sa proposition principale. Écrivez : \eng{Because voting is important, everyone should register.}
\item \textbf{Confusion \eng{Due to} / \eng{Because}} : \eng{Due to} est une préposition (+ nom) : \eng{The delay was due to the rain.} On ne dit pas \emph{due to he was ill} ; dites \eng{because he was ill}.
\item \textbf{Virgule oubliée} après \eng{Furthermore, Moreover, However, Therefore} en début de phrase : \eng{However, the results were contested.}
\item \textbf{Traduction mot à mot} : \eng{also} ne se place pas comme « aussi » en tête de phrase à tout bout de champ ; préférez \eng{Moreover} ou \eng{In addition} à l'écrit formel.
\item \textbf{Le cas \eng{So}} : \eng{So} est un caméléon. Comme coordinateur (FANBOYS), il prend une virgule \textbf{avant} : \eng{It rained, so we stayed home.} Comme adverbe de conséquence (début de phrase), il prend une virgule \textbf{après} : \eng{So, we decided to stay home.} Ne les confondez pas.
\end{itemize}
\end{attention}

\begin{methode}[Varier ses connecteurs pour une note maximale]
\begin{enumerate}
\item \textbf{Repérez} dans votre brouillon les répétitions : si trois phrases commencent par \eng{Also}, deux doivent changer.
\item \textbf{Substituez} selon le registre : \eng{Furthermore} (très formel, argument fort), \eng{Besides} (argument additionnel plus léger), \eng{What is more} (insistance, « qui plus est »).
\item \textbf{Alternez les structures} : une opposition avec \eng{However} au début d'une phrase, la suivante avec \eng{Although} à l'intérieur d'une phrase.
\item \textbf{Vérifiez la ponctuation} de chaque connecteur : virgule après l'adverbe en tête, point ou point-virgule avant \eng{however}.
\item \textbf{Comptez} : visez 8 à 12 connecteurs variés pour une composition de 250-300 mots.
\end{enumerate}
\end{methode}

\courssub{Les temps verbaux de la composition argumentative}

Le choix du temps verbal donne à votre texte sa valeur argumentative : fait général, conseil, hypothèse ou prédiction.

\begin{propriete}[Les quatre outils verbaux de l'argumentation]
\begin{itemize}
\item \textbf{Present simple} : faits généraux, vérités, définitions. \eng{Voting ensures political stability.} « Le vote garantit la stabilité politique. »
\item \textbf{Modaux \eng{must / should}} : obligation et conseil. \eng{Citizens must register before the deadline.} / \eng{Voters should compare the candidates' programmes.} « Les électeurs devraient comparer les programmes des candidats. »
\item \textbf{Second conditional} (\eng{If + past simple, ... would + V}) : hypothèse irréelle ou peu probable. \eng{If citizens voted massively, things would change.} « Si les citoyens votaient massivement, les choses changeraient. » Attention : \eng{If I were a candidate, I would listen to young people.} (\eng{were} pour toutes les personnes à l'écrit formel.)
\item \textbf{Future avec \eng{will}} : prédictions. \eng{The next election will determine the country's future.} « La prochaine élection déterminera l'avenir du pays. »
\end{itemize}
\end{propriete}

\begin{propriete}[La voix passive pour un ton objectif]
La voix passive (\eng{be + participe passé}) permet de mettre en avant le fait plutôt que la personne — très utile dans une composition formelle.

\eng{The president is elected by the people.} « Le président est élu par le peuple. »

\eng{Laws are passed by parliament.} « Les lois sont votées par le parlement. »

\eng{The results were announced on television.} « Les résultats ont été annoncés à la télévision. »
\end{propriete}

\begin{propriete}[Le discours rapporté pour citer des sources]
Pour appuyer un argument, citez des observateurs ou des rapports au \textbf{discours indirect}, avec recul du temps (concordance des temps) :

\eng{Observers stated that the election was fair.} « Les observateurs ont déclaré que l'élection était équitable. » (direct : \eng{The election is fair.} → indirect : \eng{was})

\eng{The report said that many young people had registered.} « Le rapport disait que de nombreux jeunes s'étaient inscrits. »
\end{propriete}

\begin{center}
\begin{tabularx}{\linewidth}{|X|X|X|}
\hline
\rowcolor{popblueL} \textbf{Français} & \textbf{English} & \textbf{Remarque} \\
\hline
De plus / En outre & Moreover / Furthermore & + virgule en tête de phrase \\
\hline
Cependant & However & adverbe : nouvelle phrase \\
\hline
Bien que + subjonctif & Although + indicatif & pas de subjonctif en anglais ! \\
\hline
Malgré + nom & Despite / In spite of + nom & jamais \emph{despite of} \\
\hline
À cause de & Because of / Due to & préposition + nom \\
\hline
Par conséquent & Therefore / Consequently & + virgule en tête \\
\hline
Par exemple & For example / For instance & + virgule \\
\hline
En conclusion & In conclusion / To sum up & éviter \emph{in a conclusion} \\
\hline
\end{tabularx}
\end{center}

\begin{exoresolu}[Relier et corriger]
\textbf{Énoncé.} Reliez ou corrigez les phrases suivantes à l'aide du connecteur indiqué, en respectant la ponctuation.

1. \eng{Young people should vote.} / \eng{They should also check the candidates' programmes.} (Moreover)

2. \eng{The campaign was noisy.} / \eng{It remained peaceful.} (Although)

3. \eng{The rain. Voters came early.} (Despite)

4. Corrigez : \eng{Due to the election was important, the media covered it widely.}
\tcblower
\textbf{Solution détaillée.}

1. \eng{Young people should vote. Moreover, they should check the candidates' programmes.} — \eng{Moreover} est un adverbe : point avant, virgule après. (« Les jeunes devraient voter. De plus, ils devraient examiner les programmes des candidats. »)

2. \eng{Although the campaign was noisy, it remained peaceful.} — \eng{Although} introduit la subordonnée dans la même phrase ; virgule entre les deux propositions. (« Bien que la campagne ait été bruyante, elle est restée pacifique. »)

3. \eng{Despite the rain, voters came early.} — \eng{Despite} est une préposition suivie du nom \eng{the rain}. (« Malgré la pluie, les électeurs sont venus tôt. »)

4. \eng{Due to} est une préposition : elle ne peut pas introduire \eng{the election was important} (proposition complète). Correction : \eng{Because the election was important, the media covered it widely.} ou \eng{Due to the importance of the election, the media covered it widely.}
\end{exoresolu}

\begin{exercice}[À vous de jouer]
Complétez avec un connecteur adapté et la ponctuation correcte : 1) \eng{Voting is free. \dots, every citizen can take part.} 2) \eng{\dots the long queues, voters waited patiently.} 3) \eng{Citizens must vote \dots\ their future depends on it.} 4) Transformez à la voix passive : \eng{The people elect the president.} 5) Rapportez au discours indirect : \eng{The observer said: “The vote is transparent.”}
\end{exercice}

\begin{corrige}[Exercice — À vous de jouer]
1) \eng{Therefore / Consequently} (virgule après). 2) \eng{Despite} (+ nom \eng{the long queues}). 3) \eng{because / since / as} (conjonction + proposition). 4) \eng{The president is elected by the people.} 5) \eng{The observer said (that) the vote was transparent.}
\end{corrige}

\begin{savaistu}[Les linking words au Bac : des points faciles à gagner]
Au baccalauréat camerounais, la grille d'évaluation de l'expression écrite réserve une part importante — souvent 4 à 5 points sur 20 — à la \eng{Coherence and Organisation}. Les correcteurs comptent littéralement vos connecteurs : un texte sans \eng{linking words} est plafonné, même avec de bonnes idées. À l'inverse, huit à dix connecteurs variés et bien ponctués signalent immédiatement une copie de qualité. C'est l'investissement le plus rentable de votre préparation !
\end{savaistu}

\begin{aretenir}[L'essentiel de la section]
\begin{itemize}
\item Trois familles de connecteurs : conjonctions de coordination (FANBOYS, virgule avant), adverbes de liaison (\eng{However, Therefore} : point ou point-virgule avant, virgule après), prépositions (\eng{Despite, Because of} + nom/V-ing).
\item \eng{Although} = conjonction dans la même phrase ; \eng{However} = adverbe ouvrant une nouvelle phrase.
\item Temps verbaux de l'argumentation : present simple (faits), \eng{must/should} (conseil), second conditional (hypothèses), \eng{will} (prédictions), voix passive (objectivité), discours indirect (sources).
\item Variez vos connecteurs : jamais trois fois le même ; c'est un critère majeur de la note au Bac.
\end{itemize}
\end{aretenir}

\coursec{Modèle commenté : "Why Voting Matters"}

Dans la section précédente, nous avons rassemblé nos deux « boîtes à outils » de la cohérence : les \cle{connecteurs logiques} (\eng{firstly, however, therefore, in conclusion}) et les \cle{temps verbaux} adaptés à chaque moment de l'argumentation. Il est temps maintenant de voir ces outils \emph{en action}. La meilleure façon d'apprendre à rédiger une composition de niveau Terminale est d'\textbf{analyser un modèle réussi} : non pas pour le recopier, mais pour comprendre comment il est construit, brique par brique. C'est exactement ce que nous allons faire avec la composition ci-dessous.

\courssub{Lecture du modèle : "Campaigning divides but voting unites"}

Lisez d'abord le texte en entier, sans dictionnaire, pour saisir l'idée générale. Puis relisez-le en suivant les codes couleur : en gras, les éléments de \textbf{structure} ; en italique, les \emph{connecteurs} ; les structures grammaticales ciblées sont signalées entre crochets.

\begin{textebox}[MODEL ESSAY — "Campaigning divides but voting unites" (environ 250 mots)]
\textbf{[Hook]} “The ballot is stronger than the bullet,” said Abraham Lincoln. \textbf{[Context]} In many democracies, campaign periods are noisy and tense: candidates make promises, opponents exchange insults, and communities sometimes turn against one another. \textbf{[Thesis]} \emph{However}, voting remains the ultimate peaceful tool for change because it legitimizes power and enforces accountability. \textbf{[Plan]} This essay will examine these two vital roles.

\textbf{[Body 1 — TS]} \emph{Firstly}, voting grants legitimacy to leaders. A president who \textbf{is elected} [passif] by the majority can govern with the confidence of the people. \emph{For instance}, in the 2016 elections in Ghana, power \textbf{was transferred} [passif, past simple] peacefully from one party to another because citizens had expressed their will clearly. \emph{Therefore}, elections transform a divided crowd into a united nation.

\textbf{[Body 2 — TS]} \emph{Secondly}, voting holds leaders accountable. If a government fails, citizens \textbf{can remove} [modal + conditionnel implicite] it at the next election without violence. \emph{Moreover}, the simple fear of losing votes pushes leaders to respect their promises.

\textbf{[Body 3 — TS]} \emph{Finally}, voting honours past struggles. Many Cameroonians and Africans fought for this right; ignoring it would betray their sacrifice.

\textbf{[Conclusion]} \emph{In conclusion}, far from being a mere right, voting is a sacred duty. It legitimizes power, enforces accountability and honours our history. If every young citizen voted, our democracy \textbf{would become} [conditionnel] truly unshakable.
\end{textebox}

\textbf{Glossaire (mots difficiles) :}

\begin{center}
\begin{tabularx}{\linewidth}{|l|l|X|}
\hline
\rowcolor{popblueL} \textbf{English} & \textbf{Nature} & \textbf{Français} \\
\hline
ballot & noun & bulletin de vote \\
\hline
bullet & noun & balle (de fusil) \\
\hline
to legitimize & verb & légitimer, conférer la légitimité \\
\hline
accountability & noun & reddition de comptes, responsabilité \\
\hline
to hold s.o. accountable & expression & obliger qqn à rendre des comptes \\
\hline
to transfer power & expression & transférer le pouvoir \\
\hline
sacred duty & noun phrase & devoir sacré \\
\hline
unshakable & adjective & inébranlable \\
\hline
\end{tabularx}
\end{center}

\textbf{Questions de compréhension :}
\begin{enumerate}
\item What is the writer's main thesis? Quote the sentence.
\item How many arguments support the thesis? List them in one word each.
\item Which historical example is used, and why is it effective?
\item What verb tense is used for the example? Why?
\item How does the conclusion go beyond simply repeating the thesis?
\end{enumerate}

\courssub{Carte annotée du modèle : voir la structure d'un coup d'œil}

La figure ci-dessous résume visuellement l'architecture de la composition. À gauche, le squelette du texte ; à droite, les commentaires codés par couleur : \textcolor{popgreen}{\textbf{vert = structure}}, \textcolor{popblue}{\textbf{bleu = vocabulaire}}, \textcolor{poporange}{\textbf{orange = grammaire}}, \textcolor{popink}{\textbf{rouge = connecteurs}}.

\begin{popfigure}
\begin{tikzpicture}[font=\small]
% Colonne gauche : squelette du texte
\node[draw=popline, fill=popcream, rounded corners, align=left, text width=6.2cm, inner sep=6pt] (hook) at (0,0) {\textbf{Hook} : “The ballot is stronger than the bullet.”};
\node[draw=popline, fill=popcream, rounded corners, align=left, text width=6.2cm, inner sep=6pt, below=0.25cm of hook] (thesis) {\textbf{Thesis} : voting = peaceful tool for change (legitimacy + accountability)};
\node[draw=popline, fill=popcream, rounded corners, align=left, text width=6.2cm, inner sep=6pt, below=0.25cm of thesis] (body1) {\textbf{TS1} : voting grants legitimacy \par \emph{Ex.} : Ghana 2016};
\node[draw=popline, fill=popcream, rounded corners, align=left, text width=6.2cm, inner sep=6pt, below=0.25cm of body1] (body2) {\textbf{TS2} : voting holds leaders accountable};
\node[draw=popline, fill=popcream, rounded corners, align=left, text width=6.2cm, inner sep=6pt, below=0.25cm of body2] (body3) {\textbf{TS3} : voting honours past struggles};
\node[draw=popline, fill=popcream, rounded corners, align=left, text width=6.2cm, inner sep=6pt, below=0.25cm of body3] (conc) {\textbf{Conclusion} : duty, not just a right + opening};
% Colonne droite : commentaires
\node[draw=popgreen, fill=popgreenL, rounded corners, align=left, text width=5.6cm, inner sep=6pt, right=1.1cm of thesis] (c1) {\textcolor{popgreen}{\textbf{Structure}} : thesis = réponse directe au sujet, annonce du plan};
\node[draw=popink, fill=poppinkL, rounded corners, align=left, text width=5.6cm, inner sep=6pt, right=1.1cm of body1] (c2) {\textcolor{popink}{\textbf{Connecteurs}} : \emph{firstly, for instance, therefore} guident le lecteur};
\node[draw=poporange, fill=poporangeL, rounded corners, align=left, text width=5.6cm, inner sep=6pt, right=1.1cm of body2] (c3) {\textcolor{poporange}{\textbf{Grammaire}} : passif (\eng{is elected}), modal (\eng{can remove})};
\node[draw=popblue, fill=popblueL, rounded corners, align=left, text width=5.6cm, inner sep=6pt, right=1.1cm of conc] (c4) {\textcolor{popblue}{\textbf{Vocabulaire}} : \eng{legitimacy, accountability, sacred duty} = lexique précis du thème};
% Flèches
\draw[-{Stealth}, popgreen, thick] (c1.west) -- (thesis.east);
\draw[-{Stealth}, popink, thick] (c2.west) -- (body1.east);
\draw[-{Stealth}, poporange, thick] (c3.west) -- (body2.east);
\draw[-{Stealth}, popblue, thick] (c4.west) -- (conc.east);
\end{tikzpicture}
\legende{Anatomie de la composition modèle : squelette à gauche, commentaires codés couleur à droite.}
\end{popfigure}

\courssub{Analyse de la progression argumentative}

\courssub{Introduction : les quatre mouvements obligatoires}

Observez comment l'introduction suit une progression en entonnoir, du général vers le précis :

\begin{enumerate}
\item \textbf{Le Hook (accroche)} : une citation d'Abraham Lincoln. \eng{“The ballot is stronger than the bullet.”} « Le bulletin de vote est plus fort que la balle. » L'accroche capte l'attention et annonce le thème de la paix.
\item \textbf{Le Context (mise en situation)} : deux phrases qui décrivent la réalité des campagnes électorales — bruit, tensions, promesses.
\item \textbf{La Thesis Statement (thèse)} : introduite par \eng{however}, elle répond directement au sujet et donne les deux arguments principaux : \eng{it legitimizes power and enforces accountability}.
\item \textbf{Le Plan (annonce)} : \eng{This essay will examine these two vital roles.} Une phrase simple qui guide le correcteur.
\end{enumerate}

\begin{propriete}[Les temps verbaux du modèle : une grammaire au service du sens]
Le modèle illustre parfaitement la règle vue en section 3 :
\begin{itemize}
\item \textbf{Present Simple} pour les vérités générales et la thèse : \eng{voting remains the ultimate peaceful tool} ; \eng{it legitimizes power}. « Le vote reste l'outil pacifique par excellence. »
\item \textbf{Past Simple} pour l'exemple historique daté : \eng{power was transferred peacefully in 2016}. « Le pouvoir fut transféré pacifiquement en 2016. » (Le français utilise le passé composé ou le passé simple ; l'anglais exige le Past Simple dès qu'une date est mentionnée.)
\item \textbf{Passif} pour mettre l'accent sur l'action et non sur l'agent : \eng{a president who is elected} ; \eng{power was transferred}. Forme : \cle{be} + participe passé.
\item \textbf{Modaux} pour la nuance : \eng{citizens can remove it} (capacité) ; \eng{ignoring it would betray their sacrifice} (conditionnel, hypothèse morale).
\item \textbf{Conditionnel de type 2} pour l'ouverture : \eng{If every young citizen voted, our democracy would become truly unshakable.} Forme : \cle{if + past simple, ... would + base verbale}.
\end{itemize}
\end{propriete}

\courssub{Développement : trois arguments, trois paragraphes PEEL}

Chaque paragraphe du corps suit la structure \cle{PEEL} : \textbf{P}oint (topic sentence), \textbf{E}xplanation, \textbf{E}xample, \textbf{L}ink (lien avec la thèse).

\begin{center}
\begin{tabularx}{\linewidth}{|l|X|X|}
\hline
\rowcolor{popblueL} \textbf{Paragraphe} & \textbf{Topic Sentence (anglais)} & \textbf{Fonction argumentative} \\
\hline
Body 1 & \eng{voting grants legitimacy to leaders} & Argument 1 : la \textbf{légitimité} — le vote rend le pouvoir acceptable. \\
\hline
Body 2 & \eng{voting holds leaders accountable} & Argument 2 : la \textbf{reddition de comptes} — le vote sanctionne sans violence. \\
\hline
Body 3 & \eng{voting honours past struggles} & Argument 3 : l'\textbf{hommage} — voter respecte les luttes passées. \\
\hline
\end{tabularx}
\end{center}

Remarquez la progression : de l'argument le plus institutionnel (légitimité) au plus émotionnel (sacrifice des anciens). Cette gradation donne de la force persuasive à l'ensemble.

\begin{exemplebox}[Expressions clés du modèle, traduites]
\begin{itemize}
\item \eng{The ballot is stronger than the bullet.} → « Le bulletin de vote est plus fort que la balle. »
\item \eng{Voting grants legitimacy to leaders.} → « Le vote confère la légitimité aux dirigeants. »
\item \eng{Voting holds leaders accountable.} → « Le vote oblige les dirigeants à rendre des comptes (les tient pour responsables). »
\item \eng{Far from being a mere right, voting is a sacred duty.} → « Loin d'être un simple droit, le vote est un devoir sacré. »
\item \eng{Power was transferred peacefully.} → « Le pouvoir a été transféré pacifiquement. »
\end{itemize}
\end{exemplebox}

\courssub{Conclusion : les trois gestes du bon concluant}

La conclusion du modèle accomplit trois gestes précis, à imiter :

\begin{enumerate}
\item \textbf{Reformulation de la thèse} (sans la copier) : \eng{far from being a mere right, voting is a sacred duty} — la thèse « outil pacifique » devient « devoir sacré » : même idée, mots nouveaux, intensité supérieure.
\item \textbf{Synthèse des trois arguments} en une seule phrase : \eng{It legitimizes power, enforces accountability and honours our history.} Notez la \emph{triade} (trois verbes parallèles) : procédé rhétorique très apprécié des correcteurs.
\item \textbf{Ouverture} vers l'engagement civique au-delà du vote : \eng{If every young citizen voted...} — une hypothèse au conditionnel qui projette le lecteur vers l'avenir.
\end{enumerate}

\begin{attention}[Ne copiez jamais le modèle !]
Le danger numéro un après l'étude d'un modèle est le \textbf{plagiat mécanique} : recopier des phrases entières le jour de l'examen. Le correcteur reconnaît immédiatement les formules apprises par cœur, surtout si le reste de la copie est faible. La bonne méthode : \textbf{analysez} le modèle pour en extraire la \emph{structure} (hook + context + thesis + plan ; PEEL ; conclusion en trois gestes) et les \emph{procédés} (citation, triade, conditionnel d'ouverture), puis réutilisez-les avec \textbf{vos propres idées} et \textbf{votre propre vocabulaire}. La structure est un moule ; le contenu doit être le vôtre.
\end{attention}

\begin{methode}[L'annotation active : apprendre à décortiquer un texte en 5 gestes]
Pour toute composition modèle (ou tout texte d'examen), procédez ainsi au crayon :
\begin{enumerate}
\item \textbf{Entourez} la Thesis Statement dans l'introduction (une seule phrase, en général après \eng{however} ou \eng{therefore}).
\item \textbf{Numérotez} les Topic Sentences (TS1, TS2, TS3) au début de chaque paragraphe du corps.
\item \textbf{Surlignez} les connecteurs (\eng{firstly, for instance, moreover, in conclusion}) et vérifiez qu'ils varient.
\item \textbf{Encadrez} le vocabulaire du thème « élections » (\eng{ballot, legitimacy, accountability}) et comptez les répétitions.
\item \textbf{Vérifiez} la structure PEEL de chaque paragraphe : Point, Explanation, Example, Link. S'il manque une brique, le paragraphe est déséquilibré.
\end{enumerate}
\end{methode}

\courssub{Auto-évaluation guidée du modèle}

Notons maintenant le modèle avec la grille simplifiée de l'examen, comme le ferait un correcteur.

\begin{center}
\begin{tabularx}{\linewidth}{|l|X|l|}
\hline
\rowcolor{popblueL} \textbf{Critère} & \textbf{Observation sur le modèle} & \textbf{Note} \\
\hline
Contenu /10 & Thèse claire, 3 arguments pertinents et gradués, exemple précis (Ghana 2016), ouverture originale. & 9/10 \\
\hline
Organisation /5 & 5 paragraphes nets, PEEL respecté, connecteurs variés, plan annoncé puis tenu. & 5/5 \\
\hline
Langue /5 & Passif, modaux et conditionnel corrects ; vocabulaire riche ; quasi aucune faute. & 5/5 \\
\hline
\end{tabularx}
\end{center}

\begin{exoresolu}[Repérer les procédés dans un extrait]
\textbf{Énoncé.} Voici la conclusion d'une copie d'élève : \eng{In conclusion, voting is important. It is good for democracy. Also, young people should vote.} Identifiez deux faiblesses par rapport au modèle étudié, puis réécrivez cette conclusion en trois phrases (reformulation + synthèse + ouverture).
\tcblower
\textbf{Solution détaillée.}
\textbf{Faiblesses :} (1) La thèse est répétée mot pour mot (\eng{voting is important}) au lieu d'être reformulée avec intensité ; le vocabulaire est pauvre (\eng{good, important}). (2) Il n'y a ni synthèse des arguments ni ouverture ; le connecteur \eng{also} est trop faible pour une conclusion.
\textbf{Réécriture possible :} \eng{In conclusion, voting is not merely a right but the heartbeat of democracy. It legitimizes leaders, forces them to render accounts and keeps the memory of past struggles alive. If every citizen of voting age cast a ballot in Cameroon, no leader could ever ignore the people's voice.} « En conclusion, voter n'est pas seulement un droit mais le battement de cœur de la démocratie. Cela légitime les dirigeants, les force à rendre des comptes et maintient vivante la mémoire des luttes passées. Si chaque citoyen en âge de voter glissait un bulletin dans l'urne au Cameroun, aucun dirigeant ne pourrait jamais ignorer la voix du peuple. »
Notez la triade de verbes, le présent simple pour les vérités générales et le conditionnel (\eng{if... cast... could ignore}) pour l'ouverture : exactement les procédés du modèle, avec des idées et des mots nouveaux.
\end{exoresolu}

\begin{exercice}[À votre tour : annotation active]
Reprenez le modèle « Campaigning divides but voting unites » et appliquez les 5 gestes de la méthode : (1) entourez la thèse ; (2) numérotez TS1, TS2, TS3 ; (3) surlignez six connecteurs différents ; (4) encadrez cinq mots du lexique électoral ; (5) vérifiez que le paragraphe 1 du corps contient bien les quatre briques PEEL. Puis notez la copie sur la grille (Contenu/10, Organisation/5, Langue/5) et justifiez chaque note en une phrase.
\end{exercice}

\begin{corrige}[Exercice : Annotation active et évaluation]
\textbf{1. Thesis Statement (entourée) :} \eng{However, voting remains the ultimate peaceful tool for change because it legitimizes power and enforces accountability.}

\textbf{2. Topic Sentences :}
\begin{itemize}
\item TS1 (Body 1) : \eng{Firstly, voting grants legitimacy to leaders.}
\item TS2 (Body 2) : \eng{Secondly, voting holds leaders accountable.}
\item TS3 (Body 3) : \eng{Finally, voting honours past struggles.}
\end{itemize}

\textbf{3. Six connecteurs différents surlignés :} \emph{However} (introduction), \emph{Firstly}, \emph{For instance}, \emph{Therefore} (Body 1), \emph{Secondly}, \emph{Moreover} (Body 2), \emph{Finally} (Body 3), \emph{In conclusion} (Conclusion). (Au moins six présents).

\textbf{4. Cinq mots du lexique électoral encadrés :} \eng{ballot, campaign, candidates, voting, elections, legitimacy, accountability, power, citizens, vote}. (Choix varié possible).

\textbf{5. Vérification PEEL Body 1 :}
\begin{itemize}
\item \textbf{P}oint : \eng{voting grants legitimacy to leaders.}
\item \textbf{E}xplanation : \eng{A president who is elected by the majority can govern with the confidence of the people.}
\item \textbf{E}xample : \eng{For instance, in the 2016 elections in Ghana, power was transferred peacefully...}
\item \textbf{L}ink : \eng{Therefore, elections transform a divided crowd into a united nation.} (Retour à la thèse : paix/unité).
\end{itemize}
\textbf{Structure PEEL respectée.}

\textbf{6. Notation sur la grille (Justifiée) :}
\begin{itemize}
\item \textbf{Contenu /10 : 9/10.} La thèse est nuancée (campagne divise, vote unit), les trois arguments sont distincts et progressifs (institutionnel $\to$ contrôle $\to$ mémoire), l'exemple du Ghana 2016 est daté et pertinent, l'ouverture sur la jeunesse camerounaise ancre le propos localement.
\item \textbf{Organisation /5 : 5/5.} Structure en 5 paragraphes canonique. Introduction en entonnoir (Hook, Context, Thesis, Plan). Corps en PEEL $\times$ 3. Conclusion en 3 gestes (Reformulation, Triade, Ouverture). Connecteurs variés et précis. Plan annoncé et respecté.
\item \textbf{Langue /5 : 5/5.} Maîtrise du passif (\eng{is elected, was transferred}), des modaux (\eng{can remove}), du conditionnel type 2 (\eng{would become}). Vocabulaire précis (\eng{legitimacy, accountability, sacred duty, unshakable}). Orthographe et ponctuation britanniques impeccables. Aucune erreur de concordances des temps.
\end{itemize}
\end{corrige}

Vous savez désormais \emph{lire} une composition comme un correcteur : vous repérez la thèse, les topic sentences, les connecteurs, les temps verbaux et les procédés rhétoriques. Dans la section suivante, vous passerez de l'analyse à la production : l'atelier de rédaction guidée vous mènera du brouillon à la copie finale, étape par étape.

\coursec{Atelier de rédaction guidée : Du brouillon à la copie finale}

Dans la section précédente, nous avons lu et commenté le modèle \emph{Why Voting Matters} : nous avons vu comment une thèse claire, des arguments étayés d'exemples et des connecteurs bien choisis produisent une composition remarquée. Il est temps maintenant de passer de la lecture à l'\textbf{action}. Cet atelier vous guide, étape par étape, du sujet vierge à la copie finale, exactement comme le jour du baccalauréat. Suivez la méthode, chronométrez-vous, et vous transformerez le stress de la page blanche en un automatisme rassurant.

\courssub{Étape 1 : Analyser le sujet avant d'écrire (5 minutes)}

La première cause d'échec au Bac n'est pas la grammaire : c'est le \cle{hors-sujet}. Avant d'écrire la moindre phrase, soulignez les \textbf{mots-clés} du sujet et identifiez le \textbf{verbe d'instruction}.

\begin{definition}[Les verbes d'instruction au Bac]
Le verbe d'instruction vous dit \emph{quoi faire} du sujet :
\begin{itemize}
\item \eng{Discuss} : présenter les deux faces d'une question (pour et contre) avant de conclure.
\item \eng{Argue / Give your opinion} : défendre une position claire avec des arguments.
\item \eng{Explain} : montrer les causes, les mécanismes, les raisons d'un phénomène.
\item \eng{Describe / Narrate} : raconter ou décrire (rare en Tle A, plutôt réservé aux récits).
\end{itemize}
\end{definition}

\begin{exemplebox}[Décortiquons un sujet type Bac]
Sujet : \eng{Write a composition on the importance of peaceful campaigns and credible elections.}
\begin{itemize}
\item Mots-clés : \eng{peaceful campaigns} (campagnes pacifiques), \eng{credible elections} (élections crédibles), \eng{importance} (il faut montrer \emph{pourquoi} c'est important).
\item Instruction implicite : \eng{explain / argue} — il ne s'agit PAS de raconter une élection, mais d'expliquer pourquoi ces deux conditions comptent.
\end{itemize}
\end{exemplebox}

\begin{attention}[Erreurs fatales au Bac]
Le correcteur élimine ou pénalise lourdement :
\begin{itemize}
\item Le \textbf{hors-sujet} : raconter le déroulement d'une élection à Douala au lieu d'\emph{analyser} les conditions d'un scrutin libre.
\item L'absence de \textbf{paragraphes} : un bloc unique sans alinéas est illisible.
\item Le \textbf{langage SMS} : \eng{u}, \eng{thx}, \eng{b4}, \eng{gonna} sont interdits dans une copie d'examen.
\item L'oubli de l'\textbf{introduction} ou de la \textbf{conclusion}.
\item Une copie de \textbf{moins de 200 mots} quand 250 sont demandés.
\end{itemize}
\end{attention}

\courssub{Étape 2 : Brainstorming et mind-mapping (5 minutes)}

Pendant cinq minutes, notez en vrac, au brouillon, toutes les idées, exemples et mots de vocabulaire liés au sujet. Ne jugez rien : triez ensuite.

\begin{exemplebox}[Brainstorming réussi sur le sujet « free, fair and peaceful elections »]
Idées : transparence, médias libres, éducation civique, sécurité, participation des jeunes. Exemples : le vote biométrique au Nigeria, les observateurs internationaux au Cameroun, le débat télévisé aux USA, le vote par correspondance au Royaume-Uni. Vocabulaire à caser : \eng{ballot box}, \eng{voter turnout}, \eng{transparency}, \eng{rigging}, \eng{civic duty}, \eng{peaceful}, \eng{credible}.
\end{exemplebox}

\begin{center}
\begin{tabularx}{\linewidth}{|X|X|X|}
\hline
\rowcolor{popblueL} English & Français & Exemple d'emploi \\
\hline
transparency & transparence & \eng{Transparency builds trust in the results.} \\
\hline
independent media & médias indépendants & \eng{Independent media inform voters objectively.} \\
\hline
civic education & éducation civique & \eng{Civic education teaches young people tolerance.} \\
\hline
voter turnout & taux de participation & \eng{Voter turnout was high among the youth.} \\
\hline
electoral commission & commission électorale & \eng{The electoral commission must remain neutral.} \\
\hline
to rig an election & truquer une élection & \eng{Rigging an election destroys democracy.} \\
\hline
polling station & bureau de vote & \eng{Each polling station should be well secured.} \\
\hline
\end{tabularx}
\end{center}

\courssub{Étape 3 : Le plan détaillé — la technique du « Skeleton Outline » (10 minutes)}

\begin{methode}[La technique du « Skeleton Outline » (plan-squelette)]
Ne rédigez pas encore ! Écrivez uniquement l'ossature de votre copie :
\begin{enumerate}
\item La \textbf{Thesis Statement} exacte (une phrase complète, votre position).
\item Les 2 ou 3 \textbf{Topic Sentences} (première phrase de chaque paragraphe du corps).
\item Un \textbf{exemple} par argument, noté en mots-clés.
\item La \textbf{phrase d'ouverture} de l'introduction (le « hook »).
\end{enumerate}
Ensuite seulement, développez chaque os en phrases complètes. Cette technique garantit la structure : vous ne pouvez plus « partir en vrille », car votre colonne vertébrale est déjà écrite.
\end{methode}

\begin{exemplebox}[Skeleton Outline pour le sujet type Bac]
\textbf{Hook :} \eng{Elections are the heartbeat of democracy.}\\
\textbf{Thesis :} \eng{Transparency, independent media, and civic education are the essential conditions for a free, fair and peaceful election.}\\
\textbf{Arg 1 :} \eng{First of all, transparency guarantees credible results.} — exemple : \emph{biometric voting in Nigeria, international observers}.\\
\textbf{Arg 2 :} \eng{Moreover, independent media allow voters to make informed choices.} — exemple : \emph{fact-checking during US campaigns}.\\
\textbf{Arg 3 :} \eng{Finally, civic education promotes tolerance and peace.} — exemple : \emph{youth programmes in Cameroonian schools}.\\
\textbf{Conclusion :} \eng{To conclude, without these three pillars, no election can truly reflect the people's will.}
\end{exemplebox}

\courssub{Étape 4 : Rédiger la copie — gérer son temps (50 minutes)}

L'épreuve de composition dure environ deux heures. Voici le chronogramme idéal à respecter :

\begin{popfigure}
\begin{center}
\begin{tikzpicture}[font=\small]
% Frise : 90 minutes, échelle 1 min = 0.11 cm
\draw[-{Stealth}, thick, popdark] (0,0) -- (10.6,0) node[right]{temps};
% Préparation (bleu) : 0-20 min
\fill[popblue!70] (0,0) rectangle (2.2,0.8);
\node[white, align=center] at (1.1,0.4) {\textbf{Préparation}};
\node[align=center, popblue] at (0.28,-0.5) {\tiny Analyse\\ \tiny 5 min};
\node[align=center, popblue] at (0.83,-0.5) {\tiny Brainstorm\\ \tiny 5 min};
\node[align=center, popblue] at (1.65,-0.5) {\tiny Plan\\ \tiny 10 min};
\draw[popblue, thick] (0,0) -- (0,0.8); \draw[popblue, thick] (0.55,0) -- (0.55,0.8); \draw[popblue, thick] (1.1,0) -- (1.1,0.8); \draw[popblue, thick] (2.2,0) -- (2.2,0.8);
% Écriture (vert) : 20-70 min
\fill[popgreen!70] (2.2,0) rectangle (7.7,0.8);
\node[white, align=center] at (4.95,0.4) {\textbf{Rédaction (50 min)}};
\node[align=center, popgreen!60!black] at (3.3,-0.5) {\tiny Intro 10 min};
\node[align=center, popgreen!60!black] at (5.5,-0.5) {\tiny Corps 25 min};
\node[align=center, popgreen!60!black] at (7.15,-0.5) {\tiny Conclu\\ \tiny 5 min};
\draw[popgreen!60!black, thick] (3.3,0) -- (3.3,0.8); \draw[popgreen!60!black, thick] (6.05,0) -- (6.05,0.8);
% Contrôle (orange) : 70-80 min
\fill[poporange!80] (7.7,0) rectangle (8.8,0.8);
\node[white, align=center] at (8.25,0.4) {\tiny \textbf{Relecture}};
\node[align=center, poporange!80!black] at (8.25,-0.5) {\tiny 10 min};
% Marge : 80-90 min
\fill[poppurple!50] (8.8,0) rectangle (9.9,0.8);
\node[white, align=center] at (9.35,0.4) {\tiny \textbf{Marge}};
\node[align=center, poppurple] at (9.35,-0.5) {\tiny 10 min};
% Graduations
\foreach \x/\t in {0/0, 2.2/20, 7.7/70, 8.8/80, 9.9/90} \draw[popdark] (\x,0) -- (\x,-0.12) node[below=14pt, popdark]{\tiny \t{} min};
\end{tikzpicture}
\end{center}
\legende{Chronogramme d'une épreuve de composition : bleu = préparation, vert = écriture, orange = contrôle, violet = marge de sécurité.}
\end{popfigure}

\begin{aretenir}[Règles d'or pendant la rédaction]
\begin{itemize}
\item Écrivez \textbf{lisiblement} : une copie propre gagne des points avant même d'être lue.
\item Un \textbf{alinéa} par paragraphe ; sautez une ligne entre introduction, corps et conclusion.
\item Ne cherchez pas la phrase parfaite : écrivez d'abord, corrigez à la relecture.
\item Surveillez la montre : si l'introduction dépasse 10 minutes, passez au corps.
\end{itemize}
\end{aretenir}

\courssub{Étape 5 : La relecture active — la checklist du correcteur (10 minutes)}

La relecture n'est pas un simple coup d'œil : c'est une \textbf{vérification point par point}, stylo à la main. Voici la checklist à recopier dans votre tête (ou en marge de votre brouillon) :

\begin{aretenir}[Checklist de relecture — Mémo]
\begin{itemize}
\item[$\square$] \textbf{Sujet respecté ?} Chaque paragraphe répond-il à la question posée ?
\item[$\square$] \textbf{3 parties ?} Introduction, développement (2-3 paragraphes), conclusion.
\item[$\square$] \textbf{Thesis claire ?} Ma position est-elle annoncée dès l'introduction ?
\item[$\square$] \textbf{2-3 arguments + exemples ?} Chaque argument est-il illustré ?
\item[$\square$] \textbf{Connecteurs variés ?} \eng{First of all, moreover, however, finally, to conclude}...
\item[$\square$] \textbf{S à la 3e personne du singulier ?} \eng{The government ensure\textbf{s}}, \eng{it protect\textbf{s}}.
\item[$\square$] \textbf{Temps et modaux corrects ?} Présent de vérité générale, \eng{should}, \eng{must}, passif si besoin.
\item[$\square$] \textbf{Vocabulaire précis ?} \eng{election} et non \eng{selection} ; \eng{to vote for} et non \eng{to vote someone}.
\item[$\square$] \textbf{Ponctuation anglaise ?} Pas d'espace avant le point ou la virgule ; majuscule après le point.
\item[$\square$] \textbf{250 mots minimum ?} Comptez rapidement (environ 10 mots par ligne).
\item[$\square$] \textbf{Propreté ?} Ratures propres, pas de langage SMS.
\end{itemize}
\end{aretenir}

\begin{attention}[Le piège n°1 des francophones : le « s » de la 3e personne]
C'est la faute la plus fréquente et la plus visible pour un correcteur. Au présent simple, \textbf{he / she / it} (et tout nom singulier) exigent un \textbf{-s} au verbe :\\
\eng{The electoral commission \textbf{organises} the vote.} « La commission électorale organise le scrutin. »\\
\eng{Every citizen \textbf{has} the right to vote.} « Chaque citoyen a le droit de vote. »\\
Attention aussi aux auxiliaires : \eng{The government \textbf{does} not listen.} (et non \eng{do not}, car \eng{the government} est singulier).
\end{attention}

\courssub{Sujet d'entraînement et exercice résolu}

\begin{textebox}[Sujet type Bac — à traiter en conditions réelles]
\eng{In your opinion, what are the essential conditions for a free, fair and peaceful election? Illustrate with examples from your country or abroad. (250--300 words)}\\[4pt]
« À votre avis, quelles sont les conditions essentielles d'une élection libre, équitable et pacifique ? Illustrez avec des exemples de votre pays ou de l'étranger. (250 à 300 mots) »\\[4pt]
\textbf{Glossaire :} \emph{free} (adj.) libre ; \emph{fair} (adj.) équitable ; \emph{peaceful} (adj.) pacifique ; \emph{to illustrate} (v.) illustrer ; \emph{abroad} (adv.) à l'étranger.\\[4pt]
\textbf{Questions de préparation :} 1. Quel est le verbe d'instruction ? 2. Combien d'arguments allez-vous développer ? 3. Quels exemples camerounais ou étrangers pouvez-vous mobiliser ?
\end{textebox}

\begin{exoresolu}[Corriger un extrait de copie d'élève]
Énoncé : voici le début du corps d'une copie sur le sujet ci-dessus. Relevez et corrigez les cinq fautes.\\
\eng{First of all, transparency are very important for credible elections. The government must to organise the vote correctly. Every citizen have the right to check the results. For example, in Nigeria the biometric system reduce fraud. Medias should also inform the voters.}
\tcblower
Solution détaillée :
\begin{enumerate}
\item \eng{transparency \textbf{are}} $\rightarrow$ \eng{transparency \textbf{is}} : sujet singulier.
\item \eng{must \textbf{to organise}} $\rightarrow$ \eng{must \textbf{organise}} : après un modal, base verbale sans \eng{to}.
\item \eng{Every citizen \textbf{have}} $\rightarrow$ \eng{Every citizen \textbf{has}} : 3e personne du singulier.
\item \eng{the biometric system \textbf{reduce}} $\rightarrow$ \eng{... \textbf{reduces}} : encore le « s » de la 3e personne.
\item \eng{\textbf{Medias}} $\rightarrow$ \eng{\textbf{The media}} : \eng{media} est déjà le pluriel de \eng{medium} ; pas de « s ».
\end{enumerate}
Version corrigée : \eng{First of all, transparency is very important for credible elections. The government must organise the vote correctly. Every citizen has the right to check the results. For example, in Nigeria the biometric system reduces fraud. The media should also inform the voters.}
\end{exoresolu}

\begin{exercice}[À vous de jouer — Skeleton Outline chronométré]
En 10 minutes, rédigez le plan-squelette complet (hook, thesis, 3 topic sentences, 3 exemples en mots-clés, phrase de conclusion) pour le sujet suivant : \eng{Discuss the role of young people in ensuring peaceful elections in Cameroon.} Puis rédigez uniquement l'introduction complète (4 à 5 phrases).
\end{exercice}

\begin{corrige}[Exercice — éléments de réponse]
Hook possible : \eng{Young people represent the majority of voters in Cameroon.} Thesis : \eng{By registering to vote, rejecting violence and spreading reliable information, young people are the key to peaceful elections.} Topic sentences : \eng{First of all, young voters must register and participate massively.} / \eng{Moreover, the youth should refuse to be used as tools of violence.} / \eng{Finally, young people can fight fake news on social media.} Exemples : campus campaigns in Yaoundé, peace concerts in Buea, fact-checking groups online. Conclusion : \eng{To conclude, the future of democracy lies in the hands of the youth.}
\end{corrige}

\begin{experience}[Simulation d'examen en classe]
Par binômes, échangez vos brouillons après l'étape 3 et vérifiez mutuellement le Skeleton Outline : la thesis est-elle claire ? Chaque argument a-t-il son exemple ? Puis, après la rédaction complète (50 minutes, en silence, sans dictionnaire), échangez les copies et appliquez la checklist de relecture à la copie de votre camarade avant de rendre à l'enseignant. Ce double regard est le meilleur entraînement à l'auto-correction.
\end{experience}

\begin{savaistu}[Pourquoi 250--300 mots ?]
Au baccalauréat camerounais, la composition évalue votre capacité à développer une pensée organisée en anglais, pas à remplir des pages. Les correcteurs savent qu'en dessous de 200 mots, il est impossible de présenter une introduction, deux arguments développés et une conclusion. Visez donc environ 25 à 30 lignes : c'est la longueur naturelle d'une argumentation complète et maîtrisée.
\end{savaistu}

Vous disposez désormais d'une méthode complète : analyser, brasser les idées, construire le squelette, rédiger, relire. La dernière section vous présentera la grille d'évaluation officielle et les erreurs types des francophones pour viser la note maximale.

\coursec{Grille d'évaluation \& Erreurs types des francophones (Pièges du Bac)}

Vous avez rédigé votre composition dans l'atelier précédent : félicitations ! Mais une composition ne vaut que par sa \emph{version finale}. Au Baccalauréat, ce n'est pas votre brouillon qui est noté, c'est votre copie — et elle est notée selon des critères précis que les correcteurs appliquent méthodiquement. Cette dernière section vous donne les clés du jeu : comprendre la grille officielle, traquer les erreurs typiques des candidats francophones, et adopter les stratégies de relecture des meilleurs candidats.

\courssub{La grille d'évaluation officielle (type MINESEC / Bac)}

Au Baccalauréat, l'expression écrite est notée sur 20 points répartis en cinq critères. Connaître cette grille, c'est savoir exactement où se gagnent — et où se perdent — les points.

\begin{propriete}[Grille Bac (simplifiée) : les 5 critères sur 20 points]
\begin{itemize}
  \item \cle{Contenu / Pertinence} (6 pts) : idées pertinentes par rapport au sujet, exemples concrets (contexte camerounais ou international), thèse clairement défendue, arguments développés et non simplement listés.
  \item \cle{Organisation / Cohérence} (4 pts) : structure Introduction / Développement / Conclusion respectée, paragraphes distincts, connecteurs logiques variés et bien placés, progression des idées.
  \item \cle{Correction grammaticale} (5 pts) : accords sujet-verbe, temps verbaux cohérents, syntaxe correcte (ordre des mots, prépositions, modaux).
  \item \cle{Richesse lexicale} (3 pts) : vocabulaire précis et varié, champ lexical des élections maîtrisé (\eng{turnout, polling station, ballot, manifesto, landslide victory}...), pas de répétitions.
  \item \cle{Orthographe / Ponctuation / Présentation} (2 pts) : orthographe correcte, ponctuation efficace, écriture lisible, longueur respectée (environ 250--300 mots).
\end{itemize}
\textbf{Total : 20 points.}
\end{propriete}

\begin{attention}[Ce que la grille signifie concrètement]
Beaucoup d'élèves croient que seule la grammaire compte. Faux ! La grammaire ne pèse que 5 points sur 20. Un élève qui écrit des phrases simples mais correctes, avec une thèse claire, des exemples concrets et des connecteurs bien placés, obtiendra une meilleure note qu'un élève qui tente des phrases complexes pleines de fautes. \textbf{La clarté bat la complexité.}
\end{attention}

La figure ci-dessous présente la grille sous forme de tableau à quatre niveaux de performance. Utilisez-la pour situer votre propre production.

\begin{popfigure}
\begin{center}
\begin{tikzpicture}[
  cell/.style={minimum width=3.2cm, minimum height=0.62cm, align=center, font=\scriptsize},
  header/.style={cell, fill=popblue, text=white, font=\small\bfseries},
  crit/.style={cell, fill=popblueL, font=\bfseries},
  lev1/.style={cell, fill=popgreenL},
  lev2/.style={cell, fill=popblueL!60},
  lev3/.style={cell, fill=poporangeL},
  lev4/.style={cell, fill=poppinkL},
  x=3.45cm, y=0.62cm
]
% En-têtes
\node[header] at (0,5) {Critère};
\node[header, fill=popgreen] at (1,5) {Excellent};
\node[header, fill=popblue!70] at (2,5) {Good};
\node[header, fill=poporange] at (3,5) {Fair};
\node[header, fill=poppink] at (4,5) {Weak};
% Ligne 1 : Contenu
\node[crit] at (0,4) {Contenu (6)};
\node[lev1, align=center] at (1,4){Strong thesis,\\convincing examples};
\node[lev2, align=center] at (2,4){Relevant ideas,\\few examples};
\node[lev3, align=center] at (3,4){Ideas listed,\\not developed};
\node[lev4, align=center] at (4,4){Off-topic,\\no argument};
% Ligne 2 : Organisation
\node[crit] at (0,3) {Organisation (4)};
\node[lev1, align=center] at (1,3){Clear I/D/C, logical\\flow, good connectors};
\node[lev2, align=center] at (2,3){Visible paragraphs,\\some connectors};
\node[lev3, align=center] at (3,3){Weak structure,\\few connectors};
\node[lev4, align=center] at (4,3){No paragraphs,\\jumpy ideas};
% Ligne 3 : Grammaire
\node[crit] at (0,2) {Grammaire (5)};
\node[lev1, align=center] at (1,2){Almost no errors,\\tenses controlled};
\node[lev2, align=center] at (2,2){Minor errors,\\meaning clear};
\node[lev3, align=center] at (3,2){Frequent errors,\\meaning strained};
\node[lev4, align=center] at (4,2){Errors block\\understanding};
% Ligne 4 : Vocabulaire
\node[crit] at (0,1) {Vocabulaire (3)};
\node[lev1, align=center] at (1,1){Rich, precise\\election lexis};
\node[lev2, align=center] at (2,1){Correct but\\repetitive};
\node[lev3, align=center] at (3,1){Limited,\\some misuse};
\node[lev4, align=center] at (4,1){Very poor,\\French calques};
% Ligne 5 : Forme
\node[crit] at (0,0) {Forme (2)};
\node[lev1, align=center] at (1,0){Neat, well\\punctuated};
\node[lev2, align=center] at (2,0){Readable, few\\slips};
\node[lev3, align=center] at (3,0){Untidy, weak\\punctuation};
\node[lev4, align=center] at (4,0){Illegible, no\\punctuation};
\end{tikzpicture}
\end{center}
\legende{Grille d'évaluation simplifiée : 5 critères, 4 niveaux de performance.}
\end{popfigure}

\courssub{Le Top 10 des erreurs des candidats francophones}

Observation d'abord. Voici deux phrases extraites de copies réelles types. L'une contient une erreur, l'autre est correcte. Saurez-vous dire laquelle ?

\eng{People is tired of empty promises.} — \eng{People are tired of empty promises.}

La seconde est correcte : \eng{people} est un pluriel en anglais. Voici maintenant la liste complète des dix pièges qui coûtent le plus de points aux candidats francophones, avec la règle et la correction pour chacun.

\begin{center}
\begin{tabularx}{\linewidth}{|l|X|X|}\hline
\rowcolor{popblueL} \textbf{N} & \textbf{Erreur fréquente} & \textbf{Correction et règle} \\ \hline
1 & \eng{He vote for change.} & \eng{He votes for change.} — Le \textbf{-s} de la 3\textsuperscript{e} personne du singulier au présent simple est obligatoire. \\ \hline
2 & \eng{make one's duty} & \eng{do one's duty} ; mais \eng{make a decision}. \textbf{Do} = accomplir une tâche ; \textbf{make} = créer, produire un résultat. \\ \hline
3 & \eng{since five years} & \eng{for five years} / \eng{since 2018}. \textbf{Since} + point de départ ; \textbf{for} + durée. \\ \hline
4 & \eng{Although he is poor, but he votes.} & \eng{Although he is poor, he votes.} — Jamais \eng{although} et \eng{but} dans la même phrase : un seul connecteur suffit. \\ \hline
5 & \eng{Despite of the rain...} & \eng{Despite the rain...} ou \eng{In spite of the rain...} — \eng{despite} ne prend jamais \eng{of}. \\ \hline
6 & \eng{People is...} & \eng{People are...} — \eng{people} est toujours pluriel. \\ \hline
7 & \eng{informations, advices} & \eng{information, advice} — noms \textbf{incomptables}, jamais de -s. \\ \hline
8 & \eng{I look forward to vote.} & \eng{I look forward to voting.} — Après \eng{look forward to} et \eng{be used to}, le verbe prend \textbf{-ing} (ici \eng{to} est une préposition, pas l'infinitif). \\ \hline
9 & \eng{i vote on monday in cameroon.} & \eng{I vote on Monday in Cameroon.} — Majuscules obligatoires : \eng{I}, jours, mois, pays, nationalités, langues (\eng{English}). \\ \hline
10 & Phrases kilométriques sans virgule & Une phrase = une idée. Ponctuez : virgule après un connecteur initial (\eng{Furthermore, ...}), point après chaque argument complet. \\ \hline
\end{tabularx}
\end{center}

\begin{attention}[Pièges francophones : focus sur les constructions verbales]
Ces six erreurs reviennent dans presque toutes les copies faibles. Apprenez les corrections par cœur :
\begin{enumerate}
  \item \eng{He must to vote.} $\rightarrow$ \eng{He must vote.} — Après un modal (\eng{must, can, should, will}), l'infinitif est \textbf{sans to}.
  \item \eng{The government must to ensure free elections.} $\rightarrow$ \eng{The government must ensure free elections.}
  \item \eng{I am agree with you.} $\rightarrow$ \eng{I agree with you.} — \eng{agree} est un verbe d'état : pas de \eng{be} devant.
  \item \eng{It depends of the candidates.} $\rightarrow$ \eng{It depends on the candidates.} — La préposition est \eng{on}, pas \eng{of} (calque du français « dépendre de »).
  \item \eng{The elections was fair.} $\rightarrow$ \eng{The elections were fair.} — \eng{elections} est pluriel.
  \item \eng{A big number of voters was absent.} $\rightarrow$ \eng{A large number of voters were absent.} — L'accord se fait avec \eng{voters} (pluriel), et on dit \eng{a large number}, pas \eng{a big number}.
\end{enumerate}
\end{attention}

\begin{exemplebox}[Erreurs corrigées : constructions après certains verbes]
\begin{itemize}
  \item \eng{The candidate promised to reduce taxes.} « Le candidat a promis de réduire les impôts. » — Correct : \eng{promise} + infinitif avec \eng{to}.
  \item \eng{The candidate promised reducing taxes.} — Incorrect : jamais de \textbf{-ing} après \eng{promise}.
  \item \eng{Despite he was tired, he voted.} — Incorrect : \eng{despite} n'est jamais suivi d'un sujet + verbe conjugué.
  \item \eng{Despite being tired, he voted.} « Malgré sa fatigue, il a voté. » — Correct : \eng{despite} + nom ou \textbf{-ing}.
  \item \eng{Although he was tired, he voted.} « Bien qu'il soit fatigué, il a voté. » — Correct : \eng{although} + sujet + verbe conjugué.
\end{itemize}
\end{exemplebox}

\begin{exoresolu}[Corriger un paragraphe de copie]
\textbf{Énoncé.} Le paragraphe suivant, extrait d'une copie d'élève, contient huit erreurs. Repérez-les et corrigez-les.

\eng{Since five years, the people of my town is asking for change. Although the government made many promises, but nothing changed. The citizens must to demand more informations before voting. Despite of the difficulties, I am agree that every young person should votes on election day.}

\tcblower
\textbf{Solution détaillée.}
\begin{enumerate}
  \item \eng{Since five years} $\rightarrow$ \eng{For five years} (durée = \eng{for}).
  \item \eng{the people ... is asking} $\rightarrow$ \eng{the people ... are asking} (\eng{people} = pluriel).
  \item \eng{Although ... but} $\rightarrow$ supprimer \eng{but} : \eng{Although the government made many promises, nothing changed.}
  \item \eng{must to demand} $\rightarrow$ \eng{must demand} (modal + infinitif sans \eng{to}).
  \item \eng{informations} $\rightarrow$ \eng{information} (incomptable).
  \item \eng{Despite of} $\rightarrow$ \eng{Despite} (jamais \eng{of}).
  \item \eng{I am agree} $\rightarrow$ \eng{I agree} (verbe d'état sans \eng{be}).
  \item \eng{should votes} $\rightarrow$ \eng{should vote} (pas de -s après un modal).
\end{enumerate}
\textbf{Version corrigée :} \eng{For five years, the people of my town have been asking for change. Although the government made many promises, nothing changed. The citizens must demand more information before voting. Despite the difficulties, I agree that every young person should vote on election day.}
\end{exoresolu}

\courssub{Stratégies de relecture et d'auto-correction}

Relire sa copie n'est pas « regarder » sa copie. Une relecture efficace est \emph{ciblée} : on cherche des catégories d'erreurs précises, une par une.

\begin{methode}[Auto-correction ciblée en quatre passages]
\begin{enumerate}
  \item \textbf{Passage 1 — Les verbes.} Sur votre brouillon, entourez tous les verbes conjugués. Vérifiez chacun : accord avec le sujet (le -s de la 3\textsuperscript{e} personne ?), temps cohérent, infinitif sans \eng{to} après les modaux.
  \item \textbf{Passage 2 — Les connecteurs.} Entourez tous les connecteurs (\eng{firstly, however, furthermore, although, despite...}). Vérifiez : ponctuation correcte (virgule après le connecteur initial), variété (pas trois \eng{and} de suite), pas de doublet \eng{although...but}.
  \item \textbf{Passage 3 — Vos erreurs habituelles.} Chaque élève a ses 2 ou 3 fautes récurrentes. Cherchez \emph{spécifiquement} les vôtres dans le Top 10 ci-dessus.
  \item \textbf{Passage 4 — La fluidité.} Relisez mentalement à voix haute. Si vous devez reprendre votre souffle au milieu d'une phrase, elle est trop longue : coupez-la.
\end{enumerate}
\textbf{Astuce professionnelle :} relisez votre texte \emph{à l'envers}, phrase par phrase, en commençant par la dernière. Cela casse le sens et force votre cerveau à voir la grammaire au lieu de lire « ce que vous vouliez écrire ».
\end{methode}

\begin{methode}[S'auto-évaluer avec la grille]
Reprenez la composition rédigée dans l'atelier de la section 5 et attribuez-vous une note sur 20, critère par critère :
\begin{enumerate}
  \item \textbf{Contenu /6} : ma thèse est-elle visible dans l'introduction ? Ai-je au moins deux arguments développés avec des exemples concrets ?
  \item \textbf{Organisation /4} : ai-je trois parties nettes et des paragraphes ? Mes connecteurs sont-ils variés ?
  \item \textbf{Grammaire /5} : combien d'erreurs du Top 10 ai-je trouvées en relisant ? (Chaque erreur coûte environ 0,5 point.)
  \item \textbf{Vocabulaire /3} : ai-je utilisé au moins cinq mots du champ lexical des élections sans me répéter ?
  \item \textbf{Forme /2} : ma copie est-elle lisible, ponctuée, de longueur suffisante ?
\end{enumerate}
Justifiez chaque sous-note par une phrase, puis fixez-vous un objectif d'amélioration pour la prochaine composition (exemple : « gagner 1 point en grammaire en éliminant l'oubli du -s »).
\end{methode}

\begin{exercice}[Auto-évaluation de votre composition]
Reprenez votre production de la section 5. À l'aide de la grille simplifiée (figure ci-dessus), notez-vous sur 20 et rédigez en français un court commentaire justifiant chaque sous-note (cinq phrases). Identifiez ensuite les trois erreurs du Top 10 que vous commettez le plus souvent et copiez la règle correspondante dans votre cahier de révisions.
\end{exercice}

\begin{corrige}[Exercice — éléments de réponse]
Il s'agit d'une auto-évaluation personnelle : il n'existe pas de corrigé unique. Voici un \textbf{modèle de justification} attendu : « Contenu : 5/6 — ma thèse est claire et j'ai deux exemples (les élections locales à Bamenda, le vote des jeunes), mais mon second argument manque de développement. Organisation : 3/4 — trois parties visibles, mais j'ai répété \eng{and} quatre fois. Grammaire : 3/5 — j'ai trouvé deux oublis de -s et un \eng{despite of}. Vocabulaire : 2/3 — six mots du champ lexical, mais \eng{vote} est répété. Forme : 2/2 — copie lisible et ponctuée. Total : 15/20. Objectif : éliminer l'oubli du -s et varier les connecteurs. »
\end{corrige}

\courssub{Conseils pratiques pour le jour J}

\begin{aretenir}[Le jour de l'examen : les bons réflexes]
\begin{itemize}
  \item \textbf{Gestion du temps} : pour une composition de 250--300 mots, comptez environ 5 minutes pour analyser le sujet et faire un plan, 25--30 minutes pour rédiger au brouillon puis au propre, et 5 minutes minimum pour la relecture ciblée.
  \item \textbf{Brouillon obligatoire} : ne rédigez jamais directement au propre. Le plan et les corrections se font au brouillon.
  \item \textbf{Présentation} : encre bleue ou noire uniquement, écriture soignée et régulière, alinéas visibles pour chaque paragraphe.
  \item \textbf{Pas de correcteur blanc} : raturez proprement d'un trait simple. Une copie propre avec deux ratures vaut mieux qu'une copie maculée.
  \item \textbf{Longueur} : respectez le nombre de mots demandé ; une copie trop courte perd des points de contenu, une copie interminable multiplie les fautes.
\end{itemize}
\end{aretenir}

\begin{savaistu}[L'effet Halo : pourquoi la première impression compte]
Les correcteurs du Baccalauréat lisent des centaines de copies, souvent rapidement. Les psychologues appellent \emph{effet de halo} le fait qu'une première impression positive influence toute l'évaluation. Une copie propre, clairement paragrapheée, avec une thèse visible dès l'introduction et des connecteurs en début de phrase (\eng{Firstly, Furthermore, However, In conclusion}), met immédiatement le correcteur en confiance : il lira la suite avec bienveillance. À l'inverse, une première ligne illisible ou une introduction sans thèse le dispose à chercher les fautes. Soigner les cinq premières lignes de votre copie est donc un investissement rentable — c'est aussi vrai en anglais qu'en français !
\end{savaistu}

\begin{aretenir}[L'essentiel de la section]
\begin{itemize}
  \item La composition est notée sur 20 selon cinq critères : contenu (6), organisation (4), grammaire (5), vocabulaire (3), forme (2).
  \item Les dix pièges francophones (oubli du -s, \eng{make/do}, \eng{since/for}, \eng{although...but}, \eng{despite of}, \eng{people is}, incomptables, \eng{to} + -ing, majuscules, ponctuation) doivent être vérifiés systématiquement à la relecture.
  \item La relecture se fait en passages ciblés : verbes, connecteurs, erreurs personnelles, fluidité.
  \item Le jour J : plan au brouillon, encre bleue ou noire, écriture soignée, relecture obligatoire.
\end{itemize}
\end{aretenir}

\newpage

\coursec{Activité d'intégration}

\coursec{Lesson 48 — Writing: Writes compositions on campaigning and voting for leaders}

\courssub{Objectifs de l'activité}

Cette activité d'intégration simule une \cle{condition réelle d'examen} (type épreuve d'anglais du Baccalauréat, section Essay Writing). À l'issue de la séance, l'élève sera capable de :

\begin{itemize}
\item \textbf{Comprendre} un sujet de composition et en dégager les attendus (analyse de la consigne, mots-clés, type de texte demandé) ;
\item \textbf{Planifier} sa production en 10 minutes : brainstorming, plan en trois parties (introduction -- développement -- conclusion) ;
\item \textbf{Rédiger} une composition argumentative de 250 à 300 mots sur le vote des jeunes, en mobilisant le vocabulaire du thème \emph{Elections \& Campaigning} ;
\item \textbf{Organiser} ses idées grâce aux connecteurs logiques (addition, cause, conséquence, opposition, illustration, conclusion) ;
\item \textbf{Employer} correctement les temps verbaux adaptés (présent simple pour les vérités générales, present perfect pour les bilans, futur et conditionnel pour les perspectives) ;
\item \textbf{Relire et corriger} sa copie à l'aide d'une grille d'évaluation, puis évaluer la copie d'un pair (\emph{peer assessment}) ;
\item \textbf{Gérer son temps} : 10 min de plan, 50 min de rédaction, 15 min de relecture, 15 min d'évaluation croisée.
\end{itemize}

\courssub{Situation}

Vous êtes élève en Terminale A au Lycée de Buea (South West Region). Le Conseil des Jeunes de votre arrondissement organise des élections pour renouveler ses représentants. Vous êtes candidat(e). Le journal local des jeunes, \emph{Youth Voice}, publie chaque mois une tribune libre (\emph{opinion column}) rédigée par un candidat. Le rédacteur en chef vous a envoyé la note suivante :

\begin{textebox}[Note from the Editor of \emph{Youth Voice}]
Dear Candidate,

Thank you for accepting to write for our “Youth Voice” opinion column. As you know, voter turnout among young people aged 18–25 remains worryingly low in many African countries, including Cameroon. Many young citizens say that “politics is not for them” or that “one vote changes nothing”.

Your task: write a composition of 250–300 words entitled “Young people must vote to shape their future”. Your article should:

\begin{itemize}
\item define what is at stake for the youth;
\item give two clear arguments (for example: the demographic weight of young voters; education and employment policies);
\item illustrate your point with one concrete example (Cameroon, Ghana, Kenya, or another country you know);
\item end with a strong call to action.
\end{itemize}

Deadline: Friday, 5 p.m. Remember: our readers are young people like you. Be clear, convincing and honest.

The Editor
\end{textebox}

\begin{attention}[Comprendre la consigne avant d'écrire]
À l'examen, la première erreur est le \cle{hors-sujet}. Soulignez les mots-clés du sujet : \eng{must vote} (il faut \emph{convaincre}, pas seulement décrire), \eng{young people} (le public cible), \eng{250–300 words} (la longueur obligatoire), \eng{two arguments} (ni un, ni trois), \eng{one example} (une illustration concrète), \eng{call to action} (une conclusion impérative). Chaque élément non traité fait perdre des points.
\end{attention}

\courssub{Consignes}

\textbf{Tâche 1 — Analyse du sujet (5 min, à l'oral ou par écrit).} Répondez en anglais aux questions suivantes : Who is the writer? Who are the readers? What is the purpose of the text (to inform, to describe, to persuade)? How many paragraphs do you need? Quel temps verbal domine dans un texte qui expose des vérités générales ?

\textbf{Tâche 2 — Brainstorming et plan (10 min).} Sur votre brouillon, complétez le plan suivant avec vos idées en anglais :

\begin{enumerate}
\item \textbf{Introduction} : une phrase d'accroche (\emph{hook}) + définition de l'enjeu + annonce du plan implicite.
\item \textbf{Paragraph 2} : argument 1 (\emph{demographic argument}) + connecteur + explication.
\item \textbf{Paragraph 3} : argument 2 (\emph{policies on education and employment}) + exemple concret (Cameroun ou Ghana).
\item \textbf{Conclusion} : rappel de la thèse + appel à l'action (\emph{call to action}).
\end{enumerate}

\textbf{Tâche 3 — Rédaction (50 min).} Rédigez votre composition de 250 à 300 mots au propre. Contraintes obligatoires : utilisez \textbf{au moins six connecteurs logiques différents}, \textbf{au moins huit mots du vocabulaire des élections} (voir Ressources), \textbf{au moins une phrase au conditionnel} et \textbf{une au present perfect}.

\textbf{Tâche 4 — Relecture avec grille (15 min).} Relisez votre copie en cochant chaque critère de la grille d'évaluation (voir Ressources). Corrigez au stylo d'une autre couleur : orthographe, accords sujet-verbe, temps, connecteurs.

\textbf{Tâche 5 — Évaluation croisée (15 min).} Échangez votre copie avec un voisin. Évaluez sa production avec la grille simplifiée : notez sur 20 et écrivez deux points forts (\emph{strengths}) et un axe d'amélioration (\emph{one point to improve}) en anglais. Rendez ensuite les deux copies au professeur pour la correction finale.

\courssub{Ressources à mobiliser}

\textbf{Vocabulaire utile (Elections \& Campaigning).}

\begin{center}
\begin{tabularx}{\linewidth}{|X|X|X|}\hline
\rowcolor{popblueL} English & Français & Exemple \\ \hline
to vote / a vote & voter / un vote & Every vote counts. \\ \hline
a voter / the electorate & un électeur / l'électorat & Young voters are the majority. \\ \hline
voter turnout & taux de participation & Turnout among the youth is low. \\ \hline
to run for office / a candidate & se porter candidat / un candidat & She is running for the Youth Council. \\ \hline
a campaign / to campaign & une campagne / faire campagne & They campaigned on education. \\ \hline
a polling station & un bureau de vote & Go to your polling station early. \\ \hline
a ballot paper / to cast a ballot & un bulletin / voter (déposer) & He cast his ballot peacefully. \\ \hline
to register (to vote) & s'inscrire (sur les listes) & Have you registered yet? \\ \hline
a policy & une politique (mesure) & Employment policies affect us. \\ \hline
democracy / a citizen & démocratie / un citoyen & Voting is a citizen's right. \\ \hline
\end{tabularx}
\end{center}

\textbf{Connecteurs logiques par fonction.}

\begin{center}
\begin{tabularx}{\linewidth}{|X|X|}\hline
\rowcolor{popblueL} Fonction & Connecteurs \\ \hline
Introduction du sujet & Nowadays, It is widely believed that, In many countries, \\ \hline
Addition & First of all, Moreover, Furthermore, In addition, \\ \hline
Cause / conséquence & Because, Since, Therefore, As a result, Consequently, \\ \hline
Opposition & However, Whereas, Although, Nevertheless, \\ \hline
Illustration & For example, For instance, Such as, \\ \hline
Conclusion & To conclude, In conclusion, All in all, \\ \hline
\end{tabularx}
\end{center}

\begin{propriete}[Temps verbaux de la composition argumentative]
\begin{itemize}
\item \cle{Présent simple} : vérités générales et opinions. \eng{Young people represent the majority of the population.} « Les jeunes représentent la majorité de la population. »
\item \cle{Present perfect} : bilan d'une période qui continue. \eng{Turnout has fallen in recent years.} « La participation a baissé ces dernières années. »
\item \cle{Futur (will)} : prévisions et promesses. \eng{Our vote will decide tomorrow's policies.} « Notre vote décidera des politiques de demain. »
\item \cle{Conditionnel (would)} : hypothèses et arguments. \eng{If young people voted massively, leaders would listen to them.} « Si les jeunes votaient massivement, les dirigeants les écouteraient. »
\end{itemize}
\end{propriete}

\textbf{Grille d'évaluation simplifiée (sur 20).}

\begin{center}
\begin{tabularx}{\linewidth}{|X|X|}\hline
\rowcolor{popblueL} Critère & Points \\ \hline
Respect de la consigne (thème, 250–300 mots, 2 arguments, 1 exemple, appel à l'action) & /5 \\ \hline
Organisation (introduction, développement, conclusion ; connecteurs) & /5 \\ \hline
Vocabulaire (richesse et exactitude du lexique des élections) & /4 \\ \hline
Grammaire (temps, accords, ordre des mots) et orthographe & /4 \\ \hline
Force de conviction et originalité & /2 \\ \hline
\end{tabularx}
\end{center}

\courssub{Proposition de production et corrigé commenté}

\begin{exoresolu}[Modèle de composition : “Young people must vote to shape their future”]
Rédigez la composition demandée par le rédacteur en chef de \emph{Youth Voice} (250–300 mots).

\tcblower

\textbf{Production modèle (environ 280 mots) :}

\eng{Nowadays, many young citizens say that politics does not concern them. However, they are wrong: elections decide the schools they will attend, the jobs they will find and the country they will live in. Young people must vote, because their future is at stake.}

\eng{First of all, the youth has an enormous demographic weight. In Cameroon, as in most African countries, people under twenty-five represent the majority of the population. Therefore, if all young voters registered and cast their ballots, they would become the strongest political force in the country. No candidate could ignore them any longer.}

\eng{Moreover, voting is the only peaceful way to influence the policies that concern us directly, such as education and employment. For example, in Ghana, massive youth participation in recent elections has pushed the government to invest more in technical training and jobs for young graduates. As a result, leaders there now listen carefully to students' demands. Unfortunately, in many countries, voter turnout among the youth has remained low for years, and consequently their priorities are often forgotten.}

\eng{To conclude, staying at home on election day means letting other people choose our future for us. So, dear readers, register today, go to your polling station, and cast your ballot. Your vote is your voice: use it!}

\medskip

\textbf{Commentaire des choix de langue :}

\begin{itemize}
\item \textbf{Introduction} : la phrase d'accroche (\eng{Nowadays, many young citizens say that...}) part d'une opinion répandue pour mieux la réfuter avec \eng{However}. L'enjeu est défini (\eng{their future is at stake}) et la thèse est annoncée clairement.
\item \textbf{Connecteurs} : huit connecteurs différents structurent le texte (\eng{However, First of all, Therefore, Moreover, For example, As a result, Unfortunately, To conclude}) : chaque paragraphe s'ouvre sur un signal logique.
\item \textbf{Temps verbaux} : présent simple pour les vérités générales (\eng{elections decide}), conditionnel pour l'hypothèse argumentative (\eng{if all young voters registered..., they would become}) — attention à la concordance : \emph{if + prétérit, would + base verbale} ; present perfect pour le bilan (\eng{has pushed}, \eng{has remained}).
\item \textbf{Vocabulaire thématique} : \eng{voters, registered, cast their ballots, candidate, policies, voter turnout, polling station, election day} — dix mots du champ lexical des élections, employés avec les bonnes collocations (\eng{cast a ballot}, \eng{voter turnout}).
\item \textbf{Conclusion} : appel à l'action à l'impératif (\eng{register, go, cast, use}) et formule finale mémorable (\eng{Your vote is your voice}), adaptée au public jeune du journal.
\end{itemize}
\end{exoresolu}

\begin{attention}[Erreurs types des francophones à traquer pendant la relecture]
\begin{itemize}
\item \eng{to assist to a meeting} : faux ami. Dites \eng{to attend a meeting} (« assister à »). \emph{Assist} signifie « aider ».
\item \eng{the youths are...} : \emph{youth} au sens de « la jeunesse » est singulier collectif ; \emph{young people} est plus sûr.
\item \eng{If I would vote...} : jamais \emph{would} après \emph{if}. Dites \eng{If I voted, things would change.}
\item \eng{He don't vote} : n'oubliez pas le \textbf{-s} de la troisième personne au présent simple : \eng{He doesn't vote.}
\item \eng{an unique opportunity} : devant le son « you », on dit \eng{a unique opportunity}.
\item Ordre des mots : pas d'adverbe entre le verbe et son complément. Écrivez \eng{He always votes} et non \eng{He votes always}.
\end{itemize}
\end{attention}

\courssub{Bilan de l'activité}

Cette simulation a permis de mobiliser, dans des conditions proches de l'examen, l'ensemble des compétences de la leçon : analyse d'une consigne, planification rapide, rédaction structurée en trois parties, emploi raisonné des connecteurs et des temps verbaux, usage du lexique électoral, puis relecture critique et évaluation par les pairs. L'évaluation croisée a en outre développé l'\cle{autonomie} de l'élève : savoir repérer les erreurs sur la copie d'autrui, c'est apprendre à les repérer sur la sienne.

Pour vous auto-évaluer, posez-vous les questions suivantes :

\begin{itemize}
\item Ai-je respecté le nombre de mots et traité \textbf{tous} les points de la consigne ?
\item Mon introduction définit-elle l'enjeu et annonce-t-elle ma position ?
\item Chaque paragraphe s'ouvre-t-il sur un connecteur logique approprié ?
\item Ai-je utilisé au moins un conditionnel et un present perfect correctement formés ?
\item Ma conclusion contient-elle un véritable appel à l'action ?
\end{itemize}

Si vous avez répondu « oui » à ces cinq questions, vous êtes prêt(e) pour l'épreuve de composition du Baccalauréat. Sinon, reprenez la grille, identifiez votre point faible, et réécrivez le paragraphe concerné : en anglais comme en français, \emph{good writing is rewriting} — « bien écrire, c'est réécrire ».

\newpage

\coursec{Méthodes et exercices résolus}

\courssub{Fiches méthode et exercices résolus}

\begin{methode}[Planifier une composition argumentative sur les élections]
Pour réussir une composition sur \eng{campaigning and voting for leaders}, il faut d'abord \cle{planifier} avant d'écrire. Voici la démarche en cinq étapes :
\begin{enumerate}
\item \textbf{Analyser le sujet} : souligner les mots-clés (\eng{discuss}, \eng{give your opinion}, \eng{explain why}) et identifier le type de texte demandé (\eng{opinion essay}, \eng{for/against essay}, \eng{speech}).
\item \textbf{Brainstormer} : noter en cinq minutes toutes les idées liées au thème (\eng{free and fair elections}, \eng{vote buying}, \eng{youth participation}, \eng{peaceful campaigns}).
\item \textbf{Choisir une thèse} : décider clairement de sa position. Une composition argumentative sans opinion claire perd des points.
\item \textbf{Construire le plan en trois parties} :
\begin{itemize}
\item \cle{Introduction} : accroche (\eng{hook}) + contexte + annonce du plan (\eng{thesis statement}).
\item \cle{Development} : deux ou trois paragraphes, chacun = une idée principale (\eng{topic sentence}) + exemple + explication.
\item \cle{Conclusion} : rappel de la thèse + ouverture ou appel à l'action (\eng{call to action}).
\end{itemize}
\item \textbf{Prévoir le vocabulaire et les connecteurs} : préparer une mini-liste (\eng{firstly, moreover, however, therefore, in conclusion}) et le lexique électoral (\eng{ballot box, polling station, candidate, manifesto}).
\end{enumerate}
Réflexe d'examen : consacrer environ dix minutes au plan pour quarante minutes de rédaction.
\end{methode}

\begin{exoresolu}[Analyse et planification d'un sujet type Bac]
\textbf{Énoncé.} Your school is organising a debate on the topic : \emph{“Young people in Cameroon should vote in every election.”} You want to write a composition for the school magazine. (a) Identify the type of essay required. (b) Write a thesis statement. (c) Produce a brief plan (introduction, two body paragraphs, conclusion) with the main idea of each part.
\tcblower
\textbf{Solution détaillée.}
\begin{enumerate}
\item \textbf{(a) Type d'essai.} La formulation \eng{should vote} invite à prendre position : il s'agit d'un \cle{opinion essay} (essai d'opinion). On ne se contente pas de décrire ; on défend une thèse.
\item \textbf{(b) Thesis statement.} Elle doit être claire, tenir en une phrase, et annoncer les arguments :\\
\eng{Young people in Cameroon should vote in every election because voting is a civic duty and because it gives them a voice in decisions that shape their future.}\\
« Les jeunes Camerounais devraient voter à chaque élection parce que voter est un devoir civique et parce que cela leur donne une voix dans les décisions qui façonnent leur avenir. »
\item \textbf{(c) Plan.}
\begin{itemize}
\item \textbf{Introduction} : accroche (\eng{In many polling stations in Yaoundé, very few young faces can be seen.}) + contexte (faible participation des jeunes) + thesis statement.
\item \textbf{Body paragraph 1} : \eng{Voting is a civic duty.} Idée : les droits impliquent des responsabilités ; exemple : les élections locales qui décident des écoles et des routes.
\item \textbf{Body paragraph 2} : \eng{Voting gives young people a voice.} Idée : chômage des jeunes, éducation ; exemple : un \eng{manifesto} qui promet des emplois.
\item \textbf{Conclusion} : rappel de la thèse + appel à l'action : \eng{Register, collect your voter's card, and make your voice heard.}
\end{itemize}
\end{enumerate}
\textbf{Conclusion méthodologique.} Un plan écrit au brouillon, même schématique, garantit une copie cohérente et évite les hors-sujets.
\end{exoresolu}

\begin{methode}[Rédiger introduction, paragraphes et conclusion]
\textbf{1. L'introduction (trois phrases).}
\begin{enumerate}
\item \cle{Hook} (accroche) : une question, un fait ou une scène : \eng{Have you ever wondered what would happen if nobody voted?}
\item \cle{Context} : situer le sujet : \eng{In Cameroon, elections decide who will lead our councils, regions and country.}
\item \cle{Thesis statement} : annoncer sa position et le plan.
\end{enumerate}
\textbf{2. Le paragraphe de développement (modèle P.E.E.).}
\begin{enumerate}
\item \textbf{Point} : \eng{topic sentence} — une seule idée par paragraphe : \eng{Firstly, campaigns allow voters to compare candidates.}
\item \textbf{Evidence} : exemple concret : \eng{During the campaign season in Douala, candidates present their manifestos at public meetings.}
\item \textbf{Explanation} : lien avec la thèse : \eng{This means that citizens can choose leaders who truly represent their needs.}
\end{enumerate}
\textbf{3. La conclusion.}
\begin{enumerate}
\item Reformuler la thèse avec d'autres mots (\eng{To sum up, ...}).
\item Terminer par une ouverture ou un \eng{call to action} : \eng{Let every citizen go to the polls and defend democracy.}
\end{enumerate}
Interdit : introduire une idée nouvelle dans la conclusion.
\end{methode}

\begin{exoresolu}[Rédaction d'un paragraphe de développement]
\textbf{Énoncé.} Using the P.E.E. model, write one body paragraph (5–7 lines) for the topic : \emph{“Free and fair elections are important for Cameroon.”} Begin with a suitable linking expression.
\tcblower
\textbf{Solution détaillée.}
\begin{enumerate}
\item \textbf{Choix du connecteur} : pour un premier argument, on utilise \eng{Firstly} ou \eng{First of all}.
\item \textbf{Point} : une phrase-sujet claire au présent simple (vérité générale).
\item \textbf{Evidence} : un exemple camerounais concret et crédible.
\item \textbf{Explanation} : une phrase de conclusion qui relie l'exemple à la thèse.
\end{enumerate}
\textbf{Paragraphe rédigé :}\\
\eng{Firstly, free and fair elections guarantee that leaders are chosen by the people and not by force or corruption. For example, when voters in Bamenda can cast their ballots secretly and when the votes are counted transparently, the result reflects the real will of the community. As a result, elected leaders feel responsible towards the citizens who trusted them, and they work harder to provide schools, roads and health centres. This shows that transparent voting is the foundation of good governance in Cameroon.}\\
« Premièrement, des élections libres et équitables garantissent que les dirigeants sont choisis par le peuple et non par la force ou la corruption. Par exemple, lorsque les électeurs de Bamenda peuvent voter secrètement et que les voix sont comptées de manière transparente, le résultat reflète la volonté réelle de la communauté. Par conséquent, les élus se sentent responsables envers les citoyens qui leur ont fait confiance et travaillent davantage à fournir des écoles, des routes et des centres de santé. Cela montre qu'un vote transparent est le fondement de la bonne gouvernance au Cameroun. »\\
\textbf{Justifications grammaticales.} Présent simple pour les vérités générales (\eng{guarantee, reflects, feel}) ; \eng{when} + présent pour les faits habituels ; connecteurs \eng{for example}, \eng{as a result}, \eng{this shows that} pour la cohérence.
\end{exoresolu}

\begin{methode}[Utiliser les connecteurs logiques et les bons temps verbaux]
\textbf{1. Les connecteurs par fonction.}
\begin{center}
\begin{tabularx}{\linewidth}{|X|X|}\hline
\rowcolor{popblueL} \textbf{Fonction} & \textbf{Connecteurs utiles} \\ \hline
Ajouter une idée & \eng{firstly, moreover, in addition, furthermore} \\ \hline
Opposer & \eng{however, whereas, on the other hand, although} \\ \hline
Donner un exemple & \eng{for example, for instance, such as} \\ \hline
Cause / conséquence & \eng{because, since, therefore, as a result, consequently} \\ \hline
Conclure & \eng{to sum up, in conclusion, all in all} \\ \hline
Donner son opinion & \eng{in my opinion, I believe that, it seems to me that} \\ \hline
\end{tabularx}
\end{center}
\textbf{2. Les temps verbaux d'une composition électorale.}
\begin{itemize}
\item \cle{Present simple} : vérités générales et faits : \eng{Every citizen has the right to vote.}
\item \cle{Present perfect} : bilan d'actions récentes ou lien avec le présent : \eng{The government has organised several campaigns to register young voters.}
\item \cle{Past simple} : événements datés et terminés : \eng{The last elections took place peacefully.}
\item \cle{Will / be going to} : promesses et prédictions de campagne : \eng{The candidate promises that she will build new classrooms.}
\item \cle{Modaux} : \eng{should, must, can} pour le devoir civique : \eng{Every young Cameroonian should register to vote.}
\end{itemize}
Réflexe : vérifier à la relecture que chaque connecteur placé en tête de phrase est suivi d'une virgule (\eng{However, ...}).
\end{methode}

\begin{exoresolu}[Connecteurs et temps verbaux]
\textbf{Énoncé.} (a) Complete the sentences with : \emph{however, therefore, moreover, for example, in conclusion}. (1) Campaigns inform voters about candidates ; \dots, they sometimes spread false promises. (2) Many young people do not register ; \dots, their voices are not heard. (3) The candidate promised new hospitals ; \dots, she pledged to fight corruption. (4) Some candidates organise rallies ; \dots, they visit markets and schools. (5) \dots, voting is both a right and a duty. (b) Put the verbs in brackets in the correct tense : \emph{Last year, our school (organise) a mock election. Since then, many students (decide) to register when they turn eighteen.}
\tcblower
\textbf{Solution détaillée.}
\begin{enumerate}
\item \textbf{(a) Connecteurs.}
\begin{enumerate}
\item \eng{however} : opposition entre « informer » et « fausses promesses ».
\item \eng{therefore} : conséquence logique (ne pas s'inscrire $\Rightarrow$ ne pas être entendu).
\item \eng{moreover} : ajout d'une deuxième promesse dans le même sens.
\item \eng{for example} : les visites de marchés et d'écoles illustrent les activités de campagne.
\item \eng{In conclusion} : phrase finale de synthèse (majuscule en tête de phrase).
\end{enumerate}
\item \textbf{(b) Temps verbaux.}
\begin{itemize}
\item \eng{Last year, our school \textbf{organised} a mock election.} : \cle{past simple} obligatoire avec un repère daté (\eng{last year}). « L'année dernière, notre école a organisé une élection simulée. »
\item \eng{Since then, many students \textbf{have decided} to register when they turn eighteen.} : \cle{present perfect} imposé par \eng{since then} (bilan d'une période qui continue jusqu'à maintenant). « Depuis, beaucoup d'élèves ont décidé de s'inscrire à leurs dix-huit ans. »
\end{itemize}
\end{enumerate}
\textbf{Conclusion.} Le choix du temps dépend du repère temporel : date passée = past simple ; \eng{since, for, already, just} = present perfect.
\end{exoresolu}

\begin{methode}[Relire, corriger et viser la grille d'évaluation]
Une copie de Bac est notée sur la \cle{task fulfilment} (respect du sujet), la \cle{coherence} (plan et connecteurs), le \cle{vocabulary} (lexique riche et correct) et la \cle{grammar/accuracy} (temps, accords, orthographe). D'où la méthode de relecture en quatre passages :
\begin{enumerate}
\item \textbf{Passage 1 — Contenu} : ai-je répondu au sujet ? Ai-je une thèse claire et deux ou trois arguments ?
\item \textbf{Passage 2 — Cohérence} : chaque paragraphe commence-t-il par un connecteur ? Y a-t-il une introduction et une conclusion ?
\item \textbf{Passage 3 — Grammaire} : accord sujet-verbe (\eng{he votes}, jamais \eng{he vote}), temps homogènes, articles (\eng{a, an, the}), pluriels.
\item \textbf{Passage 4 — Orthographe et présentation} : mots pièges (\eng{government, campaign, believe, responsible}), majuscules (\eng{I}, noms propres), paragraphes bien séparés.
\end{enumerate}
Réflexe final : compter les mots (en général 250 à 300 mots demandés) et garder une écriture lisible.
\end{methode}

\begin{exoresolu}[Correction d'un extrait de copie d'élève]
\textbf{Énoncé.} The following extract contains \textbf{six mistakes}. Find them and rewrite the passage correctly :\\
\emph{“Firstly, the youngs people should votes in elections because is a civic duty. In Cameroon, many citizen don't registered on the electoral lists. However, the government has organised campaigns last year to encourage them.”}
\tcblower
\textbf{Solution détaillée — les six erreurs.}
\begin{enumerate}
\item \eng{the youngs people} $\rightarrow$ \eng{young people} : \eng{young} est un adjectif ; en anglais, un adjectif ne prend jamais de \eng{-s}.
\item \eng{should votes} $\rightarrow$ \eng{should vote} : après un modal, on met la base verbale, sans \eng{-s}.
\item \eng{because is a civic duty} $\rightarrow$ \eng{because it is a civic duty} : en anglais, le sujet est obligatoire (contrairement au français « parce que c'est »).
\item \eng{many citizen} $\rightarrow$ \eng{many citizens} : pluriel obligatoire après \eng{many}.
\item \eng{don't registered} $\rightarrow$ \eng{do not register} : après \eng{do/don't}, on met la base verbale.
\item \eng{has organised campaigns last year} $\rightarrow$ \eng{organised campaigns last year} : un repère daté (\eng{last year}) impose le past simple, jamais le present perfect.
\end{enumerate}
\textbf{Passage corrigé :}\\
\eng{Firstly, young people should vote in elections because it is a civic duty. In Cameroon, many citizens do not register on the electoral lists. However, the government organised campaigns last year to encourage them.}\\
« Premièrement, les jeunes devraient voter aux élections parce que c'est un devoir civique. Au Cameroun, beaucoup de citoyens ne s'inscrivent pas sur les listes électorales. Cependant, le gouvernement a organisé des campagnes l'année dernière pour les encourager. »\\
\textbf{Conclusion.} La relecture ciblée (accords, modaux, temps) permet de récupérer facilement plusieurs points.
\end{exoresolu}

\begin{exoresolu}[Composition complète guidée — sujet type Bac]
\textbf{Énoncé.} Write a composition of about 250 words on : \emph{“Why every citizen should take part in campaigning and voting for leaders.”} Use an introduction, two body paragraphs and a conclusion.
\tcblower
\textbf{Solution détaillée — démarche puis modèle.}
\begin{enumerate}
\item \textbf{Plan au brouillon.} Thèse : participer aux campagnes et voter est essentiel. Argument 1 : les campagnes informent les électeurs. Argument 2 : voter, c'est choisir son avenir. Conclusion : appel à l'action.
\item \textbf{Connecteurs prévus} : \eng{firstly, for example, moreover, however, therefore, in conclusion}.
\item \textbf{Temps dominants} : présent simple (vérités générales), modaux \eng{should/must}.
\end{enumerate}
\textbf{Modèle de composition :}\\
\eng{Have you ever asked yourself who decides the future of your school, your road or your hospital? In every democracy, it is the citizens who decide, through campaigning and voting. In my opinion, every citizen should take part in campaigns and elections because campaigns inform voters and because voting gives people the power to choose their leaders.}\\
\eng{Firstly, campaigns are a school of democracy. During the campaign period, candidates travel across the country, from Douala to Bamenda, to present their manifestos and answer questions. For example, a candidate who promises to improve health centres must explain how he or she will do it. Therefore, citizens can compare the candidates and make a responsible choice.}\\
\eng{Moreover, voting is the most powerful weapon of the ordinary citizen. When people stay at home on election day, they allow others to decide for them. However, when everyone goes to the polling station and casts a ballot, the result truly reflects the will of the people. A single vote may look small, but millions of votes can change the destiny of a nation.}\\
\eng{In conclusion, campaigning and voting are not a waste of time ; they are the foundation of good leadership. Every Cameroonian citizen should therefore register, listen to the candidates and vote freely. Let us all go to the polls and build the future we want.}\\
\textbf{Commentaire.} Le texte respecte le plan en trois parties, emploie six connecteurs variés, mobilise le lexique électoral (\eng{manifesto, polling station, cast a ballot, register}) et garde des temps homogènes. C'est ce niveau d'exigence qui vaut l'excellence au Bac.
\end{exoresolu}

\begin{attention}[Les 5 erreurs qui coûtent des points au Bac]
\begin{enumerate}
\item \textbf{Oublier le sujet dans la phrase.} Le français dit « parce que c'est important » ; l'anglais exige toujours un sujet exprimé : \eng{because \textbf{it} is important}. Jamais \eng{because is important}.
\item \textbf{Conjuguer après un modal ou après \eng{do}.} Après \eng{should, must, can, don't}, on met la base verbale : \eng{citizens should \textbf{vote}} (pas \eng{should to vote} ni \eng{should votes}), \eng{they don't \textbf{register}}.
\item \textbf{Confondre past simple et present perfect.} Avec un repère daté (\eng{last year, in 2020, yesterday}), le past simple est obligatoire : \eng{the elections \textbf{took} place last year}. Le present perfect s'emploie avec \eng{since, for, already, just, never}.
\item \textbf{Faux amis et calques du français.} \eng{Actually} signifie « en fait » (et non « actuellement », qui se dit \eng{currently}) ; « faire une campagne » se dit \eng{to run a campaign} (pas \eng{to make a campaign}) ; « s'inscrire sur les listes électorales » = \eng{to register on the electoral roll} ; « un manifeste électoral » = \eng{a manifesto} (et non \eng{a program} dans ce contexte).
\item \textbf{Orthographe et pluriels du lexique électoral.} Les mots pièges : \eng{\textbf{gover}nment} (avec un \eng{n}), \eng{\textbf{be}lieve}, \eng{res\textbf{pon}sible}, \eng{cam\textbf{pai}gn} (\eng{g} muet), \eng{\textbf{elec}tion}. Et jamais de \eng{-s} aux adjectifs : \eng{young people}, pas \eng{youngs people}.
\end{enumerate}
Réflexe final : réserver les cinq dernières minutes de l'épreuve à la relecture ciblée de ces cinq points.
\end{attention}

\newpage

\coursec{Exercices}

\courssub{Je vérifie mes connaissances}

\begin{exercice}[QCM : vocabulaire et grammaire]
Pour chaque question, entoure la bonne réponse.
\begin{enumerate}
\item A person who presents himself or herself for election is a \ldots
\begin{itemize}
\item[A.] voter \item[B.] candidate \item[C.] ballot \item[D.] polling
\end{itemize}
\item The piece of paper on which you mark your choice is called a \ldots
\begin{itemize}
\item[A.] manifesto \item[B.] campaign \item[C.] ballot paper \item[D.] polling station
\end{itemize}
\item Choose the correct sentence :
\begin{itemize}
\item[A.] The government must to ensure free elections.
\item[B.] The government must ensure free elections.
\item[C.] The government must ensuring free elections.
\item[D.] The government must ensures free elections.
\end{itemize}
\item \eng{Although} exprime :
\begin{itemize}
\item[A.] une conséquence \item[B.] une addition \item[C.] une concession \item[D.] une cause
\end{itemize}
\item The day when people go to vote is called \ldots
\begin{itemize}
\item[A.] election day \item[B.] the counting day \item[C.] the campaign day \item[D.] the registration day
\end{itemize}
\end{enumerate}
\end{exercice}

\begin{exercice}[Vrai ou faux ? Justifie ta réponse]
Dis si chaque affirmation est vraie (V) ou fausse (F) et justifie en une phrase en anglais.
\begin{enumerate}
\item An argumentative composition is made up of an introduction, a body and a conclusion.
\item \eng{Moreover} is used to oppose two ideas.
\item In English, the word \eng{election} takes an \eng{s} in the plural : \eng{elections}.
\item \eng{Actually} means \eng{actuellement}.
\item The conclusion of a composition should introduce a brand new argument.
\end{enumerate}
\end{exercice}

\begin{exercice}[Vocabulaire : la journée du vote]
Complète chaque phrase avec le mot qui convient, choisi dans la liste : \emph{polling station, turnout, manifesto, rigging, ballot box, observers}.
\begin{enumerate}
\item Voters went to the nearest \ldots\ldots\ldots\ldots to cast their votes.
\item The party published its \ldots\ldots\ldots\ldots, a document presenting its programme.
\item The \ldots\ldots\ldots\ldots was high : almost 80 \% of registered voters came to vote.
\item International \ldots\ldots\ldots\ldots watched the elections to make sure they were free and fair.
\item After voting, each voter drops the ballot paper into the \ldots\ldots\ldots\ldots.
\item \ldots\ldots\ldots\ldots an election means manipulating the results dishonestly.
\end{enumerate}
\end{exercice}

\begin{exercice}[Conjugaison et temps verbaux]
Mets le verbe entre parenthèses à la forme qui convient.
\begin{enumerate}
\item Every citizen \ldots\ldots\ldots\ldots (have) the right to vote.
\item Last year, thousands of young people \ldots\ldots\ldots\ldots (register) on the electoral lists for the first time.
\item If the elections are transparent, the population \ldots\ldots\ldots\ldots (accept) the results.
\item While the votes \ldots\ldots\ldots\ldots (count), the candidates were waiting nervously.
\item Hate speech must \ldots\ldots\ldots\ldots (punish) by law.
\item By next October, the electoral commission \ldots\ldots\ldots\ldots (publish) the official lists.
\end{enumerate}
\end{exercice}

\courssub{Je m'entraîne}

\begin{exercice}[Traduction]
Traduis les phrases suivantes en anglais.
\begin{enumerate}
\item Voter est à la fois un droit et un devoir pour chaque citoyen.
\item Bien que la campagne électorale soit souvent animée, elle doit rester pacifique.
\item De nombreux jeunes ne s'intéressent pas à la politique ; par conséquent, ils ne s'inscrivent pas sur les listes électorales.
\item Les candidats promettent toujours beaucoup de choses, mais ils tiennent rarement leurs promesses.
\item Il est important que les élections soient libres et transparentes.
\end{enumerate}
\end{exercice}

\begin{exercice}[Combinaison de phrases avec des connecteurs]
Voici cinq phrases simples :
\begin{itemize}
\item Voting is a right. Voting is a duty. Many people ignore it. Abstention weakens democracy. Citizens should register and vote.
\end{itemize}
Combine-les en \textbf{un seul paragraphe cohérent} de 4 à 5 phrases, en utilisant obligatoirement les connecteurs suivants : \emph{Although, Moreover, Consequently, Therefore}. Souligne les connecteurs utilisés.
\end{exercice}

\begin{exercice}[Texte à trous : un discours de campagne]
Complète le texte suivant avec les mots de la liste : \emph{accountability, biometric, campaign, citizens, peacefully, promises, register, transparency}.

\begin{textebox}[A candidate's speech]
Dear fellow \ldots\ldots\ldots\ldots, my name is Grace Ebai and I am asking for your vote. During this \ldots\ldots\ldots\ldots, I have travelled across our region to listen to your concerns. I make only one promise : I will not make \ldots\ldots\ldots\ldots I cannot keep. Our elections must be organised \ldots\ldots\ldots\ldots, without violence or hate speech. I support the use of \ldots\ldots\ldots\ldots registration to avoid fraud and guarantee \ldots\ldots\ldots\ldots in the counting of votes. Leaders must also accept \ldots\ldots\ldots\ldots : they must explain their decisions to the people. So I call on all young people : go and \ldots\ldots\ldots\ldots on the electoral lists, because your vote is your voice!
\end{textebox}
\end{exercice}

\begin{exercice}[Appariements : connecteurs et fonctions]
Associe chaque connecteur (1--8) à sa fonction (a--h).
\begin{center}
\begin{tabularx}{\linewidth}{|l|X|}
\hline
\rowcolor{popblueL} \textbf{Connecteur} & \textbf{Fonction} \\
\hline
1. First of all & a. exprimer une conséquence \\
\hline
2. However & b. introduire le premier point \\
\hline
3. Moreover & c. conclure \\
\hline
4. For instance & d. exprimer une opposition \\
\hline
5. Consequently & e. ajouter une idée \\
\hline
6. Although & f. donner un exemple \\
\hline
7. In conclusion & g. exprimer une concession \\
\hline
8. In my opinion & h. donner son avis personnel \\
\hline
\end{tabularx}
\end{center}
\end{exercice}

\courssub{Je me prépare au Bac}

\begin{exercice}[Compréhension de texte et questions de langue]
Lis le texte suivant, puis réponds aux questions.

\begin{textebox}[Youth and the ballot box]
In many African countries, young people under thirty make up the majority of the population, yet they are often the least likely to vote. In Cameroon, as elsewhere, some young citizens say that politics does not concern them, or that their vote will not change anything. Others complain that registration procedures are too complicated, or that polling stations are too far from their homes.

However, history shows that elections can transform societies. When young people stay away from the ballot box, they leave important decisions — about education, employment and health — entirely in the hands of older generations. Moreover, politicians rarely listen to groups that do not vote, because they have no reason to fear losing their support.

Civil society organisations are now trying to reverse this trend. They organise debates in schools, explain registration procedures on social media and encourage first-time voters to participate. “Voting is not a favour you do for politicians,” says one youth leader in Yaoundé. “It is the only weapon that every citizen, rich or poor, possesses equally.”
\end{textebox}

\textbf{A. Comprehension questions} — Answer in your own words.
\begin{enumerate}
\item According to the text, give two reasons why young people do not vote. (2 marks)
\item What happens when young people stay away from elections? (2 marks)
\item What are civil society organisations doing to encourage young voters? Mention two actions. (2 marks)
\item Explain in your own words what the youth leader means by “the only weapon that every citizen possesses equally”. (2 marks)
\end{enumerate}

\textbf{B. Language work}
\begin{enumerate}
\item Find in the text : (a) a connector expressing opposition ; (b) a connector expressing addition. (2 marks)
\item Rewrite the following sentence in the passive voice : \eng{Politicians rarely listen to groups that do not vote.} $\rightarrow$ \eng{Groups that do not vote \ldots} (2 marks)
\item Complete with the correct preposition : \eng{Many young people are not interested \ldots politics.} / \eng{They complain \ldots the registration procedures.} (2 marks)
\item Give the noun corresponding to the verb \eng{to register} and the adjective corresponding to the noun \eng{education}. (2 marks)
\end{enumerate}
\end{exercice}

\begin{exercice}[Traduction (version)]
Traduis le passage suivant en anglais. (10 marks)

\begin{textebox}[Texte à traduire]
Les élections sont un moment important dans la vie d'un pays. Malheureusement, la campagne électorale est souvent marquée par la violence et les discours de haine. Pourtant, voter devrait être une fête de la démocratie. Si les citoyens s'inscrivent massivement sur les listes électorales et votent dans le calme, les résultats seront acceptés par tout le monde. C'est pourquoi les chefs traditionnels, les enseignants et les médias doivent sensibiliser la population, surtout les jeunes, à l'importance du vote pacifique.
\end{textebox}
\end{exercice}

\begin{exercice}[Sujet de composition type Bac]
\textbf{Topic :} \emph{Campaigning is often marred by violence and hate speech. Write an article for your school magazine suggesting ways to ensure peaceful elections in your country.}

\textbf{Instructions :}
\begin{itemize}
\item Write an \textbf{article} of \textbf{250--300 words}.
\item Give your article a catchy \textbf{title}.
\item Organise your work in three parts : an \textbf{introduction} presenting the problem, a \textbf{body} of two paragraphs (each paragraph = one suggestion + explanation + example), and a \textbf{conclusion} with a call to action.
\item Use at least \textbf{five} of these words and expressions : \emph{transparency, accountability, youth, biometric registration, hate speech, polling station, peaceful, civic education}.
\item Use at least \textbf{four} different connectors (e.g. \emph{First of all, Moreover, However, Consequently, In conclusion}).
\item Pay attention to verb tenses, subject--verb agreement and prepositions.
\end{itemize}

\textbf{Bonus (analyse de sujets) :} Pour chacun des trois sujets ci-dessous, indique : le format attendu (article, speech, essay), le public cible, le registre (formal / semi-formal), les mots-clés, une thèse possible et deux arguments.
\begin{enumerate}
\item \emph{“Should voting be compulsory for all citizens?” Write an essay giving your opinion.}
\item \emph{As president of your school's Civic Club, write a speech encouraging your schoolmates to register on the electoral lists.}
\item \emph{Write an article for a national newspaper on the role of social media in election campaigns.}
\end{enumerate}
\end{exercice}

\newpage

\coursec{Corrigés des exercices}

\courssub{Corrigés des exercices — Leçon 48}

\begin{corrige}[Exercice 1 — QCM : vocabulaire et grammaire]
\begin{enumerate}
\item \textbf{B. candidate}. \eng{A candidate is a person who presents himself or herself for election.} « Un candidat est une personne qui se présente à une élection. » \eng{A voter} (« un électeur ») is a person who votes ; \eng{a ballot} (« un bulletin de vote / un scrutin ») is the vote itself ; \eng{polling} (« le vote / le scrutin ») refers to the act of voting.
\item \textbf{C. ballot paper}. \eng{The ballot paper} (« le bulletin de vote ») \eng{is the piece of paper on which the voter marks his or her choice.} \eng{A manifesto} (« un manifeste / un programme électoral ») is a party's programme ; \eng{a polling station} (« un bureau de vote ») is the place where people vote.
\item \textbf{B. The government must ensure free elections.} Après le modal \eng{must}, on utilise toujours la base verbale (\cle{bare infinitive}), sans \eng{to} et sans \eng{-s} : \eng{must ensure}. Les réponses A (\eng{must to}), C (\eng{must ensuring}) et D (\eng{must ensures}) sont donc toutes fautives.
\item \textbf{C. une concession}. \eng{Although} (« bien que / quoique ») introduit une idée qui semble contredire la principale : \eng{Although voting is a duty, many citizens abstain.} (« Bien que voter soit un devoir, de nombreux citoyens s'abstiennent. »)
\item \textbf{A. election day}. \eng{Election day} (« le jour du scrutin ») \eng{is the day when people go to vote.} \eng{Registration day} (« le jour de l'inscription ») concerne l'inscription sur les listes, pas le vote.
\end{enumerate}
\end{corrige}

\begin{corrige}[Exercice 2 — Vrai ou faux ?]
\begin{enumerate}
\item \textbf{V (True).} \eng{An argumentative composition always has three parts : an introduction, a body and a conclusion.} (« Une composition argumentative comporte toujours trois parties : une introduction, un développement, une conclusion. »)
\item \textbf{F (False).} \eng{Moreover is used to add an idea, not to oppose two ideas ; to oppose ideas we use \textbf{however} or \textbf{although}.} (« \eng{Moreover} sert à ajouter une idée ; l'opposition s'exprime avec \eng{however} ou \eng{although}. »)
\item \textbf{V (True).} \eng{\textbf{Election} is a countable noun, so it takes an \textbf{s} in the plural : \textbf{elections}.} (Nom dénombrable régulier : \eng{an election} $\rightarrow$ \eng{elections}.)
\item \textbf{F (False).} \eng{\textbf{Actually} is a false friend : it means \textbf{in fact} or \textbf{really}, whereas \textbf{actuellement} is translated by \textbf{currently} or \textbf{at present}.} (Faux ami classique : \eng{actually} = « en fait / réellement » ; « actuellement » = \eng{currently}.)
\item \textbf{F (False).} \eng{The conclusion should sum up the arguments and give a final opinion or a call to action ; it must not introduce a new argument.} (« La conclusion récapitule les arguments et donne un avis final ou un appel à l'action ; elle n'introduit jamais d'argument nouveau. »)
\end{enumerate}
\end{corrige}

\begin{corrige}[Exercice 3 — Vocabulaire : la journée du vote]
\begin{enumerate}
\item \cle{polling station} — \eng{Voters went to the nearest polling station to cast their votes.} (« Les électeurs se sont rendus au bureau de vote le plus proche pour voter. »)
\item \cle{manifesto} — \eng{The party published its manifesto, a document presenting its programme.} (« Le parti a publié son manifeste, un document présentant son programme. »)
\item \cle{turnout} — \eng{The turnout was high : almost 80 \% of registered voters came to vote.} (« Le taux de participation était élevé : près de 80 \% des électeurs inscrits sont venus voter. »)
\item \cle{observers} — \eng{International observers watched the elections to make sure they were free and fair.} (« Des observateurs internationaux ont surveillé les élections pour s'assurer qu'elles étaient libres et équitables. ») Pluriel obligatoire : le verbe \eng{watched} est sans \eng{-s}.
\item \cle{ballot box} — \eng{After voting, each voter drops the ballot paper into the ballot box.} (« Après avoir voté, chaque électeur glisse son bulletin dans l'urne. »)
\item \cle{Rigging} — \eng{Rigging an election means manipulating the results dishonestly.} (« Truquer une élection signifie manipuler les résultats malhonnêtement. ») Majuscule en début de phrase ; le gérondif \eng{rigging} est ici sujet de la phrase.
\end{enumerate}
\end{corrige}

\begin{corrige}[Exercice 4 — Conjugaison et temps verbaux]
\begin{enumerate}
\item \textbf{has} — \eng{Every citizen \textbf{has} the right to vote.} Présent simple, 3\textsuperscript{e} personne du singulier (\eng{every citizen} = singulier) : verbe + \eng{s} ; \eng{have} devient irrégulièrement \eng{has}.
\item \textbf{registered} — \eng{Last year, thousands of young people \textbf{registered} on the electoral lists for the first time.} Prétérit simple, marqué par \eng{last year} ; \eng{register} est régulier : \eng{register} $\rightarrow$ \eng{registered} (l'accent tonique porte sur la première syllabe, donc pas de redoublement du \eng{r}).
\item \textbf{will accept} — \eng{If the elections are transparent, the population \textbf{will accept} the results.} Conditionnel de type 1 : \eng{if} + présent simple, futur avec \eng{will} dans la principale. \textbf{Règle d'or :} jamais \eng{will} dans la subordonnée introduite par \eng{if}.
\item \textbf{were being counted} — \eng{While the votes \textbf{were being counted}, the candidates were waiting nervously.} Prétérit continu à la voix passive (\cle{past continuous passive}) : \eng{was/were being} + participe passé. Les votes « sont en train d'être comptés » (passif + durée).
\item \textbf{be punished} — \eng{Hate speech must \textbf{be punished} by law.} Passif après un modal : modal + \eng{be} + participe passé (\cle{passive infinitive}). \eng{Punish} est régulier : \eng{punished}.
\item \textbf{will have published} — \eng{By next October, the electoral commission \textbf{will have published} the official lists.} Futur antérieur (\cle{future perfect} : \eng{will have} + participe passé), imposé par \eng{by} + date future : l'action sera achevée avant octobre.
\end{enumerate}
\end{corrige}

\begin{corrige}[Exercice 5 — Traduction (thème)]
\begin{enumerate}
\item \eng{Voting is both a right and a duty for every citizen.} — « à la fois... et... » = \eng{both... and...} ; le gérondif \eng{voting} en sujet.
\item \eng{Although the election campaign is often lively, it must remain peaceful.} — \textbf{Attention :} pas de \eng{but} après \eng{although} ; « animée » = \eng{lively} ou \eng{heated}.
\item \eng{Many young people are not interested in politics ; consequently, they do not register on the electoral lists.} — Préposition : \eng{interested \textbf{in}} ; « par conséquent » = \eng{consequently} / \eng{therefore}.
\item \eng{Candidates always promise many things, but they rarely keep their promises.} — « tenir une promesse » = \eng{to \textbf{keep} a promise} (faux ami : \eng{to hold a promise} n'existe pas).
\item \eng{It is important that elections \textbf{should be} free and transparent.} — ou \eng{It is important for elections to be free and transparent.} Après \eng{it is important that}, le subjonctif anglais (\eng{should be} ou base verbale \eng{be}) est la forme la plus élégante.
\end{enumerate}
\textbf{Points de vigilance :} ne jamais traduire « s'intéresser à » par \eng{interested for} ; « sur les listes électorales » = \eng{on the electoral lists} (préposition \eng{on}).
\end{corrige}

\begin{corrige}[Exercice 6 — Combinaison de phrases avec des connecteurs]
\textbf{Modèle de paragraphe attendu} (les connecteurs sont mis en évidence) :

\eng{Voting is not only a right ; it is also a duty. \textbf{Although} many people ignore it, this duty concerns every citizen. \textbf{Moreover}, abstention weakens democracy, because decisions are then taken by a minority. \textbf{Consequently}, the results of elections may not reflect the real will of the people. \textbf{Therefore}, all citizens should register on the electoral lists and vote on election day.}

(« Voter n'est pas seulement un droit, c'est aussi un devoir. Bien que beaucoup de gens l'ignorent, ce devoir concerne chaque citoyen. De plus, l'abstention affaiblit la démocratie, car les décisions sont alors prises par une minorité. Par conséquent, les résultats des élections peuvent ne pas refléter la volonté réelle du peuple. C'est pourquoi tous les citoyens devraient s'inscrire sur les listes électorales et voter le jour du scrutin. »)

\textbf{Points de vigilance :}
\begin{itemize}
\item \eng{Although} et \eng{but} ne se combinent \textbf{jamais} dans la même phrase.
\item \eng{Consequently} et \eng{Therefore} sont suivis d'une virgule en début de phrase.
\item Vérifier l'accord sujet--verbe : \eng{abstention weakens} (singulier), \eng{citizens should register} (pluriel).
\end{itemize}
\end{corrige}

\begin{corrige}[Exercice 7 — Texte à trous : un discours de campagne]
Texte complété :

\eng{Dear fellow \textbf{citizens}, my name is Grace Ebai and I am asking for your vote. During this \textbf{campaign}, I have travelled across our region to listen to your concerns. I make only one promise : I will not make \textbf{promises} I cannot keep. Our elections must be organised \textbf{peacefully}, without violence or hate speech. I support the use of \textbf{biometric} registration to avoid fraud and guarantee \textbf{transparency} in the counting of votes. Leaders must also accept \textbf{accountability} : they must explain their decisions to the people. So I call on all young people : go and \textbf{register} on the electoral lists, because your vote is your voice!}

\textbf{Justifications :}
\begin{itemize}
\item \eng{fellow citizens} : expression consacrée pour s'adresser à la foule (« chers concitoyens »).
\item \eng{promises} au pluriel : \eng{make promises I cannot keep} (« faire des promesses que je ne peux pas tenir »).
\item \eng{peacefully} : adverbe en \eng{-ly} car il modifie le verbe \eng{organised} (participe passé utilisé comme adjectif verbal).
\item \eng{biometric registration} : « l'inscription biométrique » ; l'adjectif précède le nom en anglais.
\item \eng{register} : base verbale après \eng{go and} (structure \eng{go and} + BV).
\end{itemize}
\end{corrige}

\begin{corrige}[Exercice 8 — Appariements : connecteurs et fonctions]
\begin{center}
\begin{tabularx}{\linewidth}{|l|X|X|}
\hline
\rowcolor{popblueL} \textbf{Connecteur} & \textbf{Fonction} & \textbf{Exemple} \\
\hline
1. First of all & b. introduire le premier point & \eng{First of all, voting is a right.} \\
\hline
2. However & d. exprimer une opposition & \eng{However, many citizens abstain.} \\
\hline
3. Moreover & e. ajouter une idée & \eng{Moreover, abstention weakens democracy.} \\
\hline
4. For instance & f. donner un exemple & \eng{For instance, some schools organise debates.} \\
\hline
5. Consequently & a. exprimer une conséquence & \eng{Consequently, results may be contested.} \\
\hline
6. Although & g. exprimer une concession & \eng{Although voting is free, some people do not vote.} \\
\hline
7. In conclusion & c. conclure & \eng{In conclusion, every vote counts.} \\
\hline
8. In my opinion & h. donner son avis personnel & \eng{In my opinion, voting should be encouraged.} \\
\hline
\end{tabularx}
\end{center}
\textbf{Points de vigilance :} \eng{However} s'emploie en général en début de phrase suivi d'une virgule ; \eng{Although} introduit une subordonnée et ne peut pas être suivi de \eng{but}.
\end{corrige}

\begin{corrige}[Exercice 9 — Compréhension de texte et questions de langue]
\textbf{A. Comprehension questions} (réponses modèles, en propres mots) :
\begin{enumerate}
\item \eng{Some young people think that politics does not concern them and that their vote will not change anything. Others find the registration procedures too complicated or live too far from polling stations.} (2 marks : 1 mark par raison correctement formulée.)
\item \eng{When young people stay away from elections, they let older generations take all the important decisions about education, employment and health. In addition, politicians do not listen to them because they are not afraid of losing their votes.} (2 marks.)
\item \eng{They organise debates in schools and they explain registration procedures on social media. They also encourage first-time voters to participate.} (2 marks : deux actions distinctes suffisent.)
\item \eng{The youth leader means that voting is a power that belongs to everybody equally : whether you are rich or poor, your vote has exactly the same value, so it is the fairest tool citizens have to defend themselves and change society.} (2 marks : l'idée d'égalité et de pouvoir doit apparaître.)
\end{enumerate}

\textbf{B. Language work} :
\begin{enumerate}
\item (a) Opposition : \eng{However} (ou \eng{yet}). (b) Addition : \eng{Moreover}. (2 marks.)
\item \eng{Groups that do not vote are rarely listened to (by politicians).} Passif au présent simple : \eng{are} + participe passé \eng{listened}, avec la préposition \eng{to} conservée après \eng{listen}. \textbf{Nécessaire :} ne jamais écrire \eng{are rarely listened} sans \eng{to}. (2 marks.)
\item \eng{Many young people are not interested \textbf{in} politics.} / \eng{They complain \textbf{about} the registration procedures.} (2 marks.)
\item \eng{to register} $\rightarrow$ le nom : \eng{registration} ; \eng{education} $\rightarrow$ l'adjectif : \eng{educational}. (2 marks.)
\end{enumerate}
\textbf{Barème indicatif total : 16 marks} (compréhension 8, langue 8). Pénaliser les réponses recopiées mot à mot du texte en compréhension : l'énoncé exige \eng{in your own words}.
\end{corrige}

\begin{corrige}[Exercice 10 — Traduction (version)]
\textbf{Traduction modèle :}

\eng{Elections are an important moment in the life of a country. Unfortunately, the election campaign is often marked by violence and hate speech. Yet voting should be a celebration of democracy. If citizens register massively on the electoral lists and vote peacefully, the results will be accepted by everybody. That is why traditional rulers, teachers and the media must raise the awareness of the population, especially young people, about the importance of peaceful voting.}

\textbf{Barème indicatif (10 marks) :} 1 mark par segment correct (5 segments de sens $\times$ 2 marks) ; demi-point retiré par faute de grammaire ou de vocabulaire dans un segment.

\textbf{Points de vigilance :}
\begin{itemize}
\item « Pourtant » = \eng{yet} ou \eng{however}, jamais \eng{but} en tête de phrase élaborée.
\item « s'inscrivent massivement » : \eng{register massively} ou \eng{register in large numbers}.
\item « les chefs traditionnels » = \eng{traditional rulers} (et non \eng{traditional chiefs cooks} !) ; \eng{chief} seul est accepté.
\item « sensibiliser » = \eng{to raise awareness} ou \eng{to sensitise} (orthographe britannique) ; éviter le calque \eng{to sensibilize}.
\item Futur après \eng{if} + présent : \eng{the results will be accepted} (passif avec \eng{will be} + participe passé).
\end{itemize}
\end{corrige}

\begin{corrige}[Exercice 11 — Sujet de composition type Bac]
\textbf{Barème indicatif (20 marks) :}
\begin{itemize}
\item Respect du format article (titre accrocheur, 250--300 mots, trois parties distinctes) : 4 marks.
\item Contenu : deux suggestions pertinentes, expliquées et illustrées : 6 marks.
\item Cohérence : au moins quatre connecteurs variés et correctement employés : 4 marks.
\item Lexique : au moins cinq mots du champ lexical imposé, bien utilisés : 3 marks.
\item Correction grammaticale (temps, accords, prépositions, orthographe) : 3 marks.
\end{itemize}

\textbf{Proposition de rédaction modèle :}

\eng{\textbf{Let Peace Win on Election Day}}

\eng{Elections should be a celebration of democracy. However, in many countries, including ours, campaigns are often marred by violence and hate speech. How can we make sure that election day remains peaceful ? Here are two suggestions.}

\eng{First of all, civic education must be reinforced, especially among the youth. If young citizens understand that a vote is a voice and not a weapon, they will refuse to be manipulated by violent groups. For instance, schools and community radio stations could organise debates and workshops before every election, in order to explain the rules of the democratic game.}

\eng{Moreover, transparency is the best medicine against violence. Biometric registration and the presence of independent observers at every polling station can reduce fraud, and when people trust the results, they have no reason to fight. Consequently, leaders must accept accountability : they must publish results quickly and punish anyone who spreads hate speech.}

\eng{In conclusion, peaceful elections are everybody's business. Government, candidates, the media and ordinary citizens must work together. So, on election day, let us all go out, vote calmly and let peace win !}

(Environ 200 mots dans le corps ; à développer jusqu'à 250--300 mots en détaillant les exemples et arguments.)

\textbf{Points de vigilance :}
\begin{itemize}
\item Un \textbf{article} exige un titre, un ton semi-formel et des phrases directes au lecteur ; ne pas écrire \eng{Dear Sir} (ce serait une lettre formelle).
\item Ne pas confondre \eng{peaceful} (adjectif) et \eng{peacefully} (adverbe).
\item Accords : \eng{transparency is}, \eng{leaders must accept}, \eng{citizens must work}.
\item Éviter le faux ami : « sensibiliser » = \eng{raise awareness}, pas \eng{sensibilize}.
\end{itemize}

\textbf{Bonus — Analyse des trois sujets types :}
\begin{enumerate}
\item \emph{``Should voting be compulsory for all citizens?''} — Format : \textbf{essay} (opinion essay) ; public : le correcteur / lectorat général ; registre : \textbf{formel} ; mots-clés : \emph{compulsory, right, duty, abstention, democracy} ; thèse possible : \eng{Voting should be compulsory because democracy only works when everyone participates.} Arguments : (1) un vote obligatoire rend les résultats plus représentatifs ; (2) cela rappelle que voter est un devoir civique, pas seulement un droit.
\item \emph{Speech for the Civic Club} — Format : \textbf{speech} (avec salutation : \eng{Dear fellow students,}) ; public : camarades de lycée ; registre : \textbf{semi-formel}, chaleureux et direct (\eng{you}, questions rhétoriques) ; mots-clés : \emph{register, electoral lists, your vote is your voice, first-time voters} ; thèse : \eng{Every student who is old enough should register, because our future is decided at the ballot box.} Arguments : (1) les décisions sur l'éducation et l'emploi concernent d'abord les jeunes ; (2) s'inscrire est simple et gratuit.
\item \emph{Article on social media in election campaigns} — Format : \textbf{article} de presse ; public : lecteurs d'un journal national ; registre : \textbf{formel / semi-formel} ; mots-clés : \emph{social media, campaign, fake news, hate speech, awareness} ; thèse : \eng{Social media are a powerful campaign tool, but they must be used responsibly.} Arguments : (1) ils permettent de toucher les jeunes électeurs à moindre coût ; (2) ils propagent aussi les rumeurs et les discours de haine, d'où la nécessité d'une régulation.
\end{enumerate}
\end{corrige}

\newpage

\coursec{Fiche bilan}

\begin{popfigure}
\begin{tikzpicture}[mindmap, grow cyclic, every node/.style={concept, text=white, font=\small\bfseries}, concept color=popblue, level 1/.append style={level distance=3.4cm, sibling angle=60}, level 2/.append style={level distance=2.2cm}]
\node[concept, font=\normalsize\bfseries] {Writing: Campaigning \& Voting}
  child[concept color=poporange] { node {Vocabulary} 
    child { node[font=\scriptsize] {election, ballot, candidate} }
    child { node[font=\scriptsize] {campaign, rally, vote} }
  }
  child[concept color=popgreen] { node {Structure} 
    child { node[font=\scriptsize] {introduction} }
    child { node[font=\scriptsize] {body (2--3 arguments)} }
    child { node[font=\scriptsize] {conclusion} }
  }
  child[concept color=poppurple] { node {Connectors} 
    child { node[font=\scriptsize] {firstly, moreover} }
    child { node[font=\scriptsize] {however, therefore} }
  }
  child[concept color=poppink] { node {Tenses} 
    child { node[font=\scriptsize] {present simple} }
    child { node[font=\scriptsize] {future: will / going to} }
  }
  child[concept color=popgold] { node {Method} 
    child { node[font=\scriptsize] {plan $\rightarrow$ draft $\rightarrow$ revise} }
  }
  child[concept color=popdark] { node {Pitfalls} 
    child { node[font=\scriptsize] {faux amis, word order} }
  };
\end{tikzpicture}
\legende{Carte mentale de la leçon : rédiger une composition sur la campagne électorale et le vote.}
\end{popfigure}

\begin{bilan}
\textbf{1. Lexique clé : Elections \& Campaigning}
\begin{itemize}
  \item \eng{an election} (une élection) ; \eng{to vote / a vote} (voter / un vote) ; \eng{a voter} (un électeur)
  \item \eng{a candidate} (un candidat) ; \eng{a political party} (un parti politique) ; \eng{a manifesto} (un programme électoral)
  \item \eng{a campaign} (une campagne) ; \eng{to campaign for} (faire campagne pour) ; \eng{a rally} (un meeting)
  \item \eng{a ballot paper} (un bulletin de vote) ; \eng{a ballot box} (une urne) ; \eng{a polling station} (un bureau de vote)
  \item \eng{to elect / to be elected} (élire / être élu) ; \eng{a leader / leadership} (un dirigeant / le leadership)
  \item \eng{democracy} (la démocratie) ; \eng{a citizen / citizenship} (un citoyen / la citoyenneté) ; \eng{a right / a duty} (un droit / un devoir)
\end{itemize}

\textbf{2. Règles de grammaire à connaître par c\oe ur}
\begin{center}
\begin{tabularx}{\linewidth}{|l|X|X|}
\hline
\rowcolor{popblueL} Point & Règle & Exemple \\
\hline
Présent simple & Vérités générales, arguments (3\up{e} pers. : \cle{-s}) & \eng{Voting strengthens democracy.} \\
\hline
Futur (\cle{will}) & Promesses, prédictions de campagne & \eng{Our candidate will create jobs.} \\
\hline
\cle{be going to} & Projets, intentions décidées & \eng{I am going to vote on election day.} \\
\hline
Passif (\cle{be + participe}) & Quand l'agent importe peu & \eng{Leaders are elected by the people.} \\
\hline
Modaux & \cle{should, must, can} pour conseiller et convaincre & \eng{Every citizen should vote.} \\
\hline
Comparatifs & Comparer les candidats & \eng{She is more experienced than her rival.} \\
\hline
\end{tabularx}
\end{center}

\textbf{3. Connecteurs logiques indispensables}
\begin{itemize}
  \item Ajout : \eng{firstly, secondly, moreover, furthermore, in addition}
  \item Opposition : \eng{however, whereas, on the other hand, although}
  \item Cause / conséquence : \eng{because, since, therefore, consequently, as a result}
  \item Conclusion : \eng{to conclude, in conclusion, all in all}
\end{itemize}

\textbf{4. Méthode express (3 étapes)}
\begin{enumerate}
  \item \cle{Planifier} : lire le sujet, noter 2 ou 3 arguments + exemples (contexte camerounais possible).
  \item \cle{Rédiger} : introduction (sujet + thèse), paragraphes de développement (un argument + connecteur + exemple chacun), conclusion (bilan + opinion).
  \item \cle{Relire} : accords sujet--verbe, temps, connecteurs, orthographe, nombre de mots.
\end{enumerate}
\end{bilan}

\begin{aretenir}[Checklist avant le Bac]
Avant de rendre votre copie, vérifiez chaque point :
\begin{itemize}
  \item[$\square$] J'ai bien compris le sujet et je réponds exactement à la question posée.
  \item[$\square$] Ma composition a une \cle{introduction}, un \cle{développement} (2 ou 3 paragraphes) et une \cle{conclusion}.
  \item[$\square$] Chaque paragraphe contient un argument, un connecteur et un exemple.
  \item[$\square$] J'ai utilisé au moins 4 connecteurs logiques variés (\eng{firstly, however, therefore, in conclusion}).
  \item[$\square$] J'ai employé le vocabulaire du thème : \eng{election, candidate, campaign, vote, democracy}.
  \item[$\square$] Les verbes à la 3\up{e} personne du singulier prennent bien un \cle{-s} au présent simple.
  \item[$\square$] J'ai évité les faux amis : \eng{to elect} (pas \eng{to choose} pour « élire »), \eng{a leader} (pas \eng{a chief} pour « dirigeant »).
  \item[$\square$] Je n'ai pas traduit mot à mot : « voter pour » = \eng{to vote for}, « faire campagne » = \eng{to campaign}.
  \item[$\square$] Mes temps verbaux sont cohérents (présent pour les arguments, \eng{will} pour les promesses).
  \item[$\square$] J'ai relu l'orthographe : \eng{government, democracy, candidate, responsible}.
  \item[$\square$] Ma conclusion donne clairement mon opinion personnelle.
  \item[$\square$] Mon écriture est lisible et le nombre de mots est respecté.
\end{itemize}
\end{aretenir}

\findecours

\end{document}
